%% Volume II, Chapter 15: Difference Equations and Recurrences
\chapter{Difference Equations and Recurrences}
\label{ch:difference-equations}

\abstract{A computer never meets a continuum. Whether it solves a differential equation, sums a series or evaluates an orthogonal polynomial, what it runs is a recurrence, a rule that produces the next number from a few previous ones. Linear recurrences have a theory that runs line by line parallel to that of linear differential equations, with powers in place of exponentials and the unit disc in place of the left half-plane. Their stability, the question whether small errors grow, decides whether an algorithm built on them can be trusted, and the answer may depend on nothing more than the direction in which the recurrence is run. This chapter develops that theory and tests it on examples that return in later volumes.}

Recurrences have run through this book from its first part. \Vcref{1}{ch:counting} solved the recurrence of the Tower of Hanoi and that of the Fibonacci numbers, found Binet's formula and packed sequences into generating functions; \cref{ch:vectors-matrices} wrote the Fibonacci recurrence as the power of a matrix, \cref{ch:determinants} met a three-term recurrence for tridiagonal determinants, and \cref{ch:canonical-forms} proved, through the Jordan form of the companion matrix, that every linear recurrence with constant coefficients is solved by sums of terms $n^jr^n$. Part~II of this volume added recurrences of another origin. Euler's method replaces a differential equation by the recurrence $y_{k+1}=y_k+hf(t_k,y_k)$, and whether this recurrence reproduces the decay of the solutions decided whether the method could be used at all (\cref{sec:euler,sec:systems-numerics}). The power series method of \cref{ch:series-solutions} turned differential equations into recurrences for coefficients, and the functions it produced, the Legendre polynomials and the Bessel functions, satisfy three-term recurrences in their degree and order, whose use in floating-point arithmetic that chapter left open. This chapter gathers these threads into a theory of \term{difference equations}, the discrete counterparts of the differential equations of \cref{ch:first-order-equations,ch:linear-higher-order,ch:linear-ode-systems}.

The parallel is close. Differentiation becomes the difference operator $\Delta$, integration becomes summation, the exponential $e^{\lambda t}$ becomes the power $r^n$, the Wronskian becomes the Casorati determinant, the matrix exponential $e^{tA}$ becomes the matrix power $A^n$, and the condition $\operatorname{Re}\lambda<0$ for decay becomes $|r|<1$. Most proofs carry over with little change, and several become simpler, because a recurrence can be solved by running it: existence and uniqueness of solutions, which cost \cref{ch:first-order-equations} a Lipschitz condition and successive approximations, are immediate. But the parallel also breaks down, and the places where it does are revealing. A difference equation can forget its past, a first-order recurrence can oscillate, and a simple nonlinear recurrence can behave chaotically, none of which happens for the corresponding differential equations. And there is a question that has no counterpart in the continuous theory at all: a recurrence that determines a sequence uniquely may still be useless for computing it, because it amplifies the rounding errors, and running it in the opposite direction may cure the defect.

\Cref{sec:difference-operators} introduces the shift and difference operators and the calculus of finite differences. \Cref{sec:first-order-recurrences} treats recurrences of the first order, with their fixed points and the cobweb diagram. \Cref{sec:constant-coefficient-recurrences} solves linear recurrences with constant coefficients completely, with the Casorati determinant and the discrete method of undetermined coefficients, and \cref{sec:recurrence-gf} returns to generating functions. \Cref{sec:difference-systems} studies systems $\mathbf x_{n+1}=A\mathbf x_n$ and their stability, side by side with the systems of differential equations of \cref{ch:linear-ode-systems}. \Cref{sec:three-term} treats three-term recurrences with variable coefficients, their minimal and dominant solutions and the direction in which they may be run, with the Bessel functions and the orthogonal polynomials as examples. \Cref{sec:ode-to-difference} looks at the recurrences produced by discretising differential equations, \cref{sec:nonlinear-recurrences} gives a glimpse of nonlinear recurrences, from Newton's method to the logistic map, and \cref{sec:recurrence-computation} tests the theory by computation.

The sequences of this chapter are written $(y_n)_{n\ge0}$, or $(y_n)$ for short, and their terms are complex numbers unless said otherwise. The index $n$ plays the role of the time $t$ of \cref{ch:first-order-equations,ch:linear-higher-order,ch:linear-ode-systems} and of the variable $x$ of \cref{ch:series-solutions}; it is always a natural number. When a differential equation is discretised, with the step $h$ and the times $t_n=t_0+nh$, its approximations are written $y_n$, so that the index $k$ of the approximations $y_k$ in those chapters becomes $n$ here, and the iterates $\mathbf x_k$ of a matrix become $\mathbf x_n$. Vector-valued sequences are written in bold, $(\mathbf x_n)$. A recurrence of order $k$ is written with its highest index first, $y_{n+k}+c_{k-1}y_{n+k-1}+\cdots+c_0y_n=0$, in analogy with $y^{(k)}+c_{k-1}y^{(k-1)}+\cdots+c_0y=0$; the recurrence $a_n=c_1a_{n-1}+\cdots+c_ka_{n-k}$ of \vcref{1}{prop:linear-recurrence-k} and \cref{cor:linear-recurrence} is the case in which the coefficient $c_j$ there is $-c_{k-j}$ here.

\section{Sequences and Difference Operators}
\label{sec:difference-operators}

\subsection{Shift and Difference}

The set $\mathbb C^{\mathbb N}$ of all complex sequences is a vector space under termwise addition and multiplication by numbers, and \cref{sec:linear-maps} introduced two linear operators on it, the left shift $S_-$ and the forward difference $S_--\operatorname{id}$. They are the basic operators of this chapter, and we give them the names that are customary in the calculus of finite differences.

\begin{definition}[Shift and difference operators]
\label{def:shift-difference}
The \term{shift operator} $E$ and the \term{forward difference operator} $\Delta$ on $\mathbb C^{\mathbb N}$ are
\[
(Ey)_n\coloneqq y_{n+1},\qquad(\Delta y)_n\coloneqq y_{n+1}-y_n\qquad(n\in\mathbb N),
\]
so that $E=S_-$ and $\Delta=E-I$, where $I$ is the identity. Their powers are $E^0=\Delta^0=I$, $E^{k+1}=EE^k$ and $\Delta^{k+1}=\Delta\Delta^k$, and for a polynomial $q(\lambda)=\sum_jq_j\lambda^j$ we write $q(E)\coloneqq\sum_jq_jE^j$.
\end{definition}

We write $\Delta y_n$ for $(\Delta y)_n$ when no confusion is possible, and we apply the operators also to sequences given by a formula, as in $\Delta n^2=(n+1)^2-n^2=2n+1$. The same operator acts on functions of a real variable by $(\Delta f)(x)=f(x+1)-f(x)$; \cref{sec:matrix-of-map} found its matrix on $\mathcal P_3$. A sequence $y_n=f(n)$ given by a function has the difference $\Delta y_n=(\Delta f)(n)$, so the two uses agree. Since $(E^ky)_n=y_{n+k}$, a polynomial in $E$ forms a fixed combination of consecutive terms: $q(E)y$ is the sequence $\sum_jq_jy_{n+j}$. Powers of $\Delta$ are combinations of shifts, and conversely.

\begin{proposition}[Powers of $\Delta$ and $E$]
\label{prop:difference-powers}
For every sequence $(y_n)$ and all $k,n\in\mathbb N$,
\[
\Delta^ky_n=\sum_{j=0}^k(-1)^{k-j}\binom kj\,y_{n+j},\qquad y_{n+k}=\sum_{j=0}^k\binom kj\,\Delta^jy_n .
\]
\end{proposition}

\begin{proof}
The operators $E$ and $I$ commute, and the binomial theorem (\vcref{1}{thm:binomial}) holds for any two commuting linear operators, since its proof uses only the distributive law and the commutativity of the two summands. Hence $\Delta^k=(E-I)^k=\sum_j\binom kj(-1)^{k-j}E^j$ and $E^k=(I+\Delta)^k=\sum_j\binom kj\Delta^j$; applying both sides to $y$ and taking the term with index $n$ gives the two formulas.
\end{proof}

For instance $\Delta^2y_n=y_{n+2}-2y_{n+1}+y_n$, the second difference of \cref{sec:special-matrices}, shifted by one. The rules for differences resemble those for derivatives, with one shift in the product rule.

\begin{proposition}[Rules for differences]
\label{prop:difference-rules}
Let $(y_n)$ and $(z_n)$ be sequences, $\alpha,\beta,r\in\mathbb C$ and $k\ge1$.
\begin{enumerate}
\item $\Delta(\alpha y+\beta z)=\alpha\Delta y+\beta\Delta z$.
\item $\Delta(yz)_n=y_n\,\Delta z_n+z_{n+1}\,\Delta y_n$.
\item $\Delta n^{\underline k}=k\,n^{\underline{k-1}}$ and $\Delta\binom nk=\binom n{k-1}$.
\item $Er^n=r\cdot r^n$ and $\Delta r^n=(r-1)\,r^n$.
\end{enumerate}
\end{proposition}

\begin{proof}
(i) is clear. (ii) $y_{n+1}z_{n+1}-y_nz_n=y_n(z_{n+1}-z_n)+z_{n+1}(y_{n+1}-y_n)$. (iii) is \veqref{1}{eq:falling-difference}, and dividing it by $k!$ gives the second identity, which is also Pascal's rule (\vcref{1}{prop:pascal}). (iv) $r^{n+1}=r\cdot r^n$, and $r^{n+1}-r^n=(r-1)r^n$.
\end{proof}

Part (iii) is the discrete form of $(x^k)'=kx^{k-1}$, with the falling factorial $n^{\underline k}=n(n-1)\cdots(n-k+1)$ in the place of the power, as \vcref{1}{sec:recurrences} observed; this is why falling factorials, and not powers, are the natural polynomials of the calculus of differences. Part (iv) is the discrete form of $(e^{\lambda t})'=\lambda e^{\lambda t}$: the geometric sequences $r^n$ are the eigenvectors of $E$, as the exponentials are eigenfunctions of $D$ (\cref{sec:eigen-definitions}). Since $r^n=e^{n\ln r}$ for $r>0$, the rate $\lambda$ of an exponential and the ratio $r$ of a geometric sequence correspond through $r=e^\lambda$.

\subsection{Summation}

Summation is to differences what integration is to derivatives.

\begin{theorem}[Fundamental theorem of summation]
\label{thm:summation}
\leavevmode
\begin{enumerate}
\item $\Delta y=0$ if and only if $(y_n)$ is constant.
\item Let $(y_n)$ be a sequence and let $(Y_n)$ satisfy $\Delta Y=y$. Then for all $a\le b$ in $\mathbb N$,
\[
\sum_{n=a}^{b-1}y_n=Y_b-Y_a .
\]
The sequence $Y_n\coloneqq\sum_{m=0}^{n-1}y_m$ satisfies $\Delta Y=y$, and every other solution of $\Delta Z=y$ differs from it by a constant.
\end{enumerate}
\end{theorem}

\begin{proof}
(i) $\Delta y=0$ means $y_{n+1}=y_n$ for all $n$, so $y_n=y_0$ by induction. (ii) The sum is $\sum_{n=a}^{b-1}(Y_{n+1}-Y_n)=Y_b-Y_a$ by telescoping (\vcref{1}{prop:telescoping}). For $Y_n=\sum_{m<n}y_m$, $Y_{n+1}-Y_n=y_n$. If also $\Delta Z=y$, then $\Delta(Z-Y)=0$, and $Z-Y$ is constant by (i).
\end{proof}

A sequence $Y$ with $\Delta Y=y$ is a discrete antiderivative of $y$, and the theorem says that sums are computed by finding one, as integrals are computed by antiderivatives (\vcref{1}{thm:ftc2}). By \cref{prop:difference-rules}\,(iii), the antiderivative of $n^{\underline k}$ is $n^{\underline{k+1}}/(k+1)$, which gives \veqref{1}{eq:falling-sum} once more, and by (iv) the antiderivative of $r^n$ is $r^n/(r-1)$ for $r\ne1$, which gives the geometric sum $\sum_{n=0}^{N-1}r^n=(r^N-1)/(r-1)$. Integration by parts has a counterpart as well.

\begin{proposition}[Summation by parts]
\label{prop:summation-by-parts}
For all sequences $(u_n)$, $(v_n)$ and all $a\le b$ in $\mathbb N$,
\[
\sum_{n=a}^{b-1}u_n\,\Delta v_n=u_bv_b-u_av_a-\sum_{n=a}^{b-1}v_{n+1}\,\Delta u_n .
\]
\end{proposition}

\begin{proof}
By \cref{prop:difference-rules}\,(ii), $\Delta(uv)_n=u_n\Delta v_n+v_{n+1}\Delta u_n$. Summing from $a$ to $b-1$ and using \cref{thm:summation}\,(ii) with $Y=uv$ gives $u_bv_b-u_av_a=\sum_{n=a}^{b-1}u_n\Delta v_n+\sum_{n=a}^{b-1}v_{n+1}\Delta u_n$.
\end{proof}

The shift in the last sum, $v_{n+1}$ instead of $v_n$, is the only difference from integration by parts. The formula, often called Abel's transformation, was used by Abel in his study of the binomial series~\cite{abel1826}. As an example, take $u_n=n$ and $v_n=2^n$, so that $\Delta v_n=2^n$ and $\Delta u_n=1$:
\[
\sum_{n=0}^{N-1}n\,2^n=N2^N-\sum_{n=0}^{N-1}2^{n+1}=N2^N-(2^{N+1}-2)=(N-2)\,2^N+2,
\]
which for $N=3$ gives $0+2+8=10$, as it should.

\subsection{Newton's Forward Difference Formula}

Taylor's formula reconstructs a polynomial from its derivatives at one point. Its discrete counterpart reconstructs every sequence from its differences at $0$.

\begin{theorem}[Newton's forward difference formula]
\label{thm:newton-forward}
\leavevmode
\begin{enumerate}
\item For every sequence $(y_n)$ and every $n\in\mathbb N$,
\[
y_n=\sum_{k=0}^n\binom nk\,\Delta^ky_0 .
\]
\item Let $d\in\mathbb N$. A sequence satisfies $\Delta^{d+1}y=0$ if and only if there is a polynomial $P$ of degree at most $d$ with $y_n=P(n)$ for all $n$, and then
\[
P(x)=\sum_{k=0}^d\Delta^ky_0\,\binom xk=\sum_{k=0}^d\frac{\Delta^ky_0}{k!}\,x^{\underline k} .
\]
\end{enumerate}
\end{theorem}

\begin{proof}
(i) is the second formula of \cref{prop:difference-powers} with $n=0$ and $k$ replaced by $n$. (ii) If $y_n=P(n)$ with a polynomial $P$ of degree $m\le d$, then $\Delta y_n=(\Delta P)(n)$, and $\Delta P$ is a polynomial of degree $m-1$ if $m\ge1$, since the terms $x^m$ cancel in $P(x+1)-P(x)$ and the leading coefficient of $P$ is multiplied by $m$, and it is $0$ if $P$ is constant. So $\Delta^{d+1}P=0$. Conversely, if $\Delta^{d+1}y=0$, then $\Delta^ky=\Delta^{k-d-1}\Delta^{d+1}y=0$ for every $k>d$, and (i) gives $y_n=\sum_{k\le d}\binom nk\Delta^ky_0$, which is the polynomial $P$ of the statement evaluated at $n$: for $k>n$ the binomial coefficient $\binom nk$ vanishes, so the sums over $k\le n$ and over $k\le d$ agree.
\end{proof}

Newton's formula is the discrete Taylor series, $\binom nk=n^{\underline k}/k!$ taking the place of $x^k/k!$ and $\Delta^ky_0$ that of $f^{(k)}(0)$. Part (ii) is the counterpart of the fact that a function with $f^{(d+1)}=0$ is a polynomial of degree at most $d$, and it gives a practical test: a sequence of values of a polynomial of degree $d$ has constant $d$th differences. \Cref{tab:difference-table} shows the \term{difference table} of $y_n=n^3$, each column formed by subtracting neighbours in the column before it. The third differences are constant, equal to $3!=6$, and the fourth vanish; the top entries $0$, $1$, $6$, $6$ are the differences $\Delta^ky_0$, so that by the theorem
\[
n^3=n^{\underline1}+\frac62\,n^{\underline2}+\frac66\,n^{\underline3}=n^{\underline1}+3\,n^{\underline2}+n^{\underline3} .
\]
Summing with \cref{thm:summation} and $\Delta n^{\underline{k+1}}=(k+1)n^{\underline k}$ now gives the sum of cubes,
\[
\sum_{m=0}^{n-1}m^3=\frac{n^{\underline2}}2+n^{\underline3}+\frac{n^{\underline4}}4=\Bigl(\frac{n(n-1)}2\Bigr)^2 ,
\]
as one checks by expanding; for $n=4$ both sides are $36$. This is the method of finite calculus that \vcref{1}{sec:recurrences} described~\cite[Sect.~2.6]{graham1994concrete}, now with a systematic way to convert powers into falling factorials.

\begin{table}
\sdfloat{%
\begin{sdtable}{colspec={rrrrrr}, row{1}={font=\small\bfseries}}
\toprule
$n$ & $y_n=n^3$ & $\Delta y_n$ & $\Delta^2y_n$ & $\Delta^3y_n$ & $\Delta^4y_n$ \\
\midrule
$0$ & $0$ & $1$ & $6$ & $6$ & $0$ \\
$1$ & $1$ & $7$ & $12$ & $6$ & $0$ \\
$2$ & $8$ & $19$ & $18$ & $6$ & $0$ \\
$3$ & $27$ & $37$ & $24$ & $6$ & \\
$4$ & $64$ & $61$ & $30$ & & \\
$5$ & $125$ & $91$ & & & \\
$6$ & $216$ & & & & \\
\bottomrule
\end{sdtable}}{The difference table of $y_n=n^3$. Each entry is the entry to its left in the next row minus the entry to its left in the same row. The third differences are constant and the fourth vanish (\cref{thm:newton-forward}); the first row lists $\Delta^ky_0$.\label{tab:difference-table}}
\end{table}

\begin{note}{Stirling numbers and the calculus of finite differences}
The coefficients in $x^m=\sum_kS(m,k)\,x^{\underline k}$ are the \emph{Stirling numbers of the second kind}; $S(3,1)=1$, $S(3,2)=3$ and $S(3,3)=1$ were found above. By \cref{thm:newton-forward}, $k!\,S(m,k)=\Delta^kx^m|_{x=0}=\sum_j(-1)^{k-j}\binom kjj^m$, which by \vcref{1}{prop:surjections} is the number of surjections of an $m$-element set onto a $k$-element set; so $S(m,k)$ counts the partitions of an $m$-element set into $k$ nonempty blocks. Graham, Knuth and Patashnik~\cite[Sects.~2.6 and~6.1]{graham1994concrete} develop these numbers and the finite calculus built on them. Boole's treatise of 1860~\cite{boole1860} presented the calculus of finite differences as an algebra of the operators $E$ and $\Delta$, manipulating relations such as $E=I+\Delta$ formally, as we did in \cref{prop:difference-powers}, and Milne-Thomson's book~\cite{milnethomson1933} is a classical account of the subject, including its analytic side.
\end{note}

\begin{table}
\sdfloat{%
\begin{sdtable}{colspec={ll}, row{1}={font=\small\bfseries}}
\toprule
continuous & discrete \\
\midrule
derivative $Df=f'$ & difference $\Delta y_n=y_{n+1}-y_n$ \\
$f'=0$ $\Leftrightarrow$ $f$ constant & $\Delta y=0$ $\Leftrightarrow$ $y$ constant \\
$\int_a^bf'=f(b)-f(a)$ & $\sum_{n=a}^{b-1}\Delta y_n=y_b-y_a$ \\
integration by parts & summation by parts \\
$Dx^k=kx^{k-1}$ & $\Delta n^{\underline k}=kn^{\underline{k-1}}$ \\
$De^{\lambda t}=\lambda e^{\lambda t}$ & $Er^n=r\cdot r^n$, $\Delta r^n=(r-1)r^n$ \\
Taylor series $\sum_kf^{(k)}(0)\,t^k/k!$ & Newton's formula $\sum_k\Delta^ky_0\binom nk$ \\
$D^{d+1}f=0$ $\Leftrightarrow$ polynomial of degree $\le d$ & $\Delta^{d+1}y=0$ $\Leftrightarrow$ polynomial of degree $\le d$ \\
$p(D)y=0$, solutions $t^je^{\lambda t}$ & $p(E)y=0$, solutions $n^jr^n$ \\
Wronskian, Abel's formula & Casorati determinant (\cref{thm:casorati}) \\
\bottomrule
\end{sdtable}}{The calculus of differences beside the differential calculus. The left column refers to \vcref{1}{ch:integration,ch:series} and \cref{ch:linear-higher-order}, the right column to this chapter.\label{tab:difference-calculus}}
\end{table}

\Cref{tab:difference-calculus} collects the correspondence between the two calculi, including the solutions of linear equations with constant coefficients that \cref{sec:constant-coefficient-recurrences} establishes.

\subsection{Difference Equations}

\begin{definition}[Difference equations]
\label{def:difference-equation}
A \term{difference equation of order $k$ in normal form}, or \term{recurrence of order $k$}, is a relation
\[
y_{n+k}=f(n,y_n,y_{n+1},\ldots,y_{n+k-1})\qquad(n\in\mathbb N)
\]
for a sequence $(y_n)$, with a given function $f$; a \term{solution} is a sequence that satisfies it for every $n$. The equation is \term{linear} if it has the form
\[
y_{n+k}+a_{k-1}(n)\,y_{n+k-1}+\cdots+a_0(n)\,y_n=g_n\qquad(n\in\mathbb N)
\]
with given sequences $a_0,\ldots,a_{k-1}$ and $g$, \term{homogeneous} if $g=0$, and has \term{constant coefficients} if the $a_j$ do not depend on $n$.
\end{definition}

By \cref{prop:difference-powers}, a linear equation written with the shifts $y_{n+j}=E^jy_n$ can be rewritten with the differences $\Delta^jy_n$, and conversely: $y_{n+2}-2y_{n+1}+y_n=g_n$ is $\Delta^2y_n=g_n$. The two names, difference equation and recurrence, therefore refer to the same objects, and we use them interchangeably. An equation such as $\Delta y_n=y_n^2$ is a recurrence $y_{n+1}=y_n+y_n^2$ of the first order. As for differential equations, a linear equation of order $k$ whose leading coefficient does not vanish can be divided by it, and we always assume this normal form, with leading coefficient $1$.

\begin{proposition}[Initial value problems]
\label{prop:recurrence-ivp}
Let a recurrence of order $k$ in normal form be given. For all initial values $y_0,\ldots,y_{k-1}$ it has exactly one solution. If the equation is linear and $a_0(n)\ne0$ for all $n$, then a solution is determined by any $k$ consecutive terms $y_m,\ldots,y_{m+k-1}$, and the recurrence can be run backwards:
\[
y_n=\frac{g_n-y_{n+k}-a_{k-1}(n)\,y_{n+k-1}-\cdots-a_1(n)\,y_{n+1}}{a_0(n)} .
\]
\end{proposition}

\begin{proof}
Apply the recursion theorem (\vcref{1}{thm:recursion}) to the vectors $\mathbf s_n\coloneqq(y_n,\ldots,y_{n+k-1})$, which satisfy $\mathbf s_{n+1}=\mathbf F(n,\mathbf s_n)$ with $\mathbf F(n,\mathbf s)\coloneqq(s_2,\ldots,s_k,f(n,\mathbf s))$, as \vcref{1}{sec:recurrences} did for $k=2$. In the linear case with $a_0(n)\ne0$, the displayed formula computes $y_n$ from $y_{n+1},\ldots,y_{n+k}$, so the terms before index $m$ are determined by induction downwards.
\end{proof}

Here recurrences differ sharply from differential equations. Existence and uniqueness, which \cref{sec:existence-uniqueness} proved under a Lipschitz condition and which can fail without it, are automatic: a recurrence is its own algorithm, and its solutions never cease to exist after finitely many steps. The price is paid elsewhere. Running a recurrence backwards needs $a_0(n)\ne0$, and an equation such as $y_{n+1}=0\cdot y_n$, whose solutions are $(y_0,0,0,\ldots)$, cannot be reversed at all: it forgets its initial value after one step. Solutions of linear differential equations never forget, since their values at one time determine them at all times (\cref{thm:linear-existence}).

\section{First-Order Recurrences}
\label{sec:first-order-recurrences}

\subsection{Constant Coefficients and Fixed Points}

The simplest recurrence, $y_{n+1}=ay_n+b$ with constants $a$ and $b$, was solved in \vcref{1}{prop:first-order}: $y_n=a^ny_0+b(1+a+\cdots+a^{n-1})$. The best way to read the solution is through the fixed point. A \term{fixed point} or \term{equilibrium} of a recurrence $y_{n+1}=f(y_n)$ is a number $y^*$ with $f(y^*)=y^*$; the constant sequence $y_n=y^*$ is then a solution.

\begin{proposition}[First-order recurrences with constant coefficients]
\label{prop:first-order-constant}
Let $a,b\in\mathbb C$ with $a\ne1$, let $y^*\coloneqq b/(1-a)$, and let $(y_n)$ solve $y_{n+1}=ay_n+b$. Then
\[
y_n-y^*=a^n(y_0-y^*)\qquad(n\in\mathbb N).
\]
\begin{enumerate}
\item If $|a|<1$, then $y_n\to y^*$ for every $y_0$, and $|y_n-y^*|=|a|^n|y_0-y^*|$.
\item If $|a|>1$ and $y_0\ne y^*$, then $|y_n-y^*|\to\infty$.
\item If $|a|=1$, then $|y_n-y^*|=|y_0-y^*|$ for all $n$; for $a=-1$ the solution alternates between $y_0$ and $b-y_0$.
\end{enumerate}
For $a=1$ the solutions are $y_n=y_0+nb$. For real $a$, $b$ and $y_0\ne y^*$, the terms $y_n$ approach or leave $y^*$ monotonically if $a>0$ and alternate around it if $a<0$.
\end{proposition}

\begin{proof}
Subtracting $y^*=ay^*+b$ from $y_{n+1}=ay_n+b$ gives $y_{n+1}-y^*=a(y_n-y^*)$, and the formula follows by induction. (i)--(iii) follow from $|a^n|=|a|^n$ and \vcref{1}{thm:standard-limits}. For $a=1$ the recurrence is $\Delta y=b$, and \cref{thm:summation} gives $y_n=y_0+nb$. The sign of $a^n(y_0-y^*)$ is constant for $a>0$ and alternates for $a<0$.
\end{proof}

The substitution $z_n=y_n-y^*$, which moves the fixed point to $0$, is the discrete form of the steady state of \cref{sec:linear-first-order}, and the trichotomy $|a|<1$, $|a|>1$, $|a|=1$ is the discrete form of the trichotomy $\lambda<0$, $\lambda>0$, $\lambda=0$ for $y'=\lambda y+b$. Two differences deserve attention. A negative $a$ produces alternation, which no solution of $y'=\lambda y+b$ shows. And the case $a=0$, in which $y_n=b$ for every $n\ge1$ whatever $y_0$ was, has no continuous counterpart either: the solution reaches its fixed point in one step.

A loan shows a repelling fixed point at work. A balance $B_0$ borrowed at the monthly interest rate $r$ and repaid by a fixed monthly payment $P$ evolves as $B_{n+1}=(1+r)B_n-P$. Here $a=1+r>1$, and the fixed point $B^*=P/r$ is the balance on which the payment just covers the interest. By the proposition, $B_n=B^*+(1+r)^n(B_0-B^*)$: if $P>rB_0$, the balance falls below $B^*$ and away from it, at an accelerating rate, and reaches $0$ after $N$ months if $B^*+(1+r)^N(B_0-B^*)=0$, that is, if
\[
P=\frac{rB_0}{1-(1+r)^{-N}} .
\]
For $B_0=200\,000$, $r=0.005$ (six per cent a year) and $N=360$ months, the payment is $P=1199.10$, and the borrower pays $431\,676.38$ in all. Because the fixed point repels, small changes are amplified: the balance after $N$ months depends on the payment through $B_N=(1+r)^NB_0-P\,((1+r)^N-1)/r$, so a payment smaller by one unit leaves $((1+r)^N-1)/r=1004.52$ units unpaid at the end.

Euler's method for Newton's law of cooling, $T'=-k(T-T_{\mathrm a})$ with $k>0$ (\cref{sec:ode-basics}), gives $T_{n+1}=T_n-hk(T_n-T_{\mathrm a})$, a recurrence with $a=1-hk$ and the fixed point $T_{\mathrm a}$, the temperature of the surroundings, which is also the equilibrium of the differential equation. By the proposition, the approximations approach $T_{\mathrm a}$ monotonically if $hk<1$, with alternation if $1<hk<2$, and they diverge with alternation if $hk>2$; this is the stability condition $h<2/|\lambda|$ of \cref{sec:euler}. The exact solution, sampled at the times $t_n=nh$, satisfies the same kind of recurrence with $a=e^{-hk}$, which lies in $(0,1)$ for every step: sampling never produces alternation or growth, and \cref{sec:difference-systems} explains this in general.

\subsection{Variable Coefficients}

With variable coefficients the solution is a product and a sum, the discrete counterpart of the integrating factor of \cref{thm:linear-first-order}.

\begin{theorem}[First-order linear recurrences]
\label{thm:first-order-variable}
Let $(a_n)$ and $(b_n)$ be sequences, and for $m\le n$ let $\Phi(n,m)\coloneqq a_{n-1}a_{n-2}\cdots a_m$, an empty product being $1$, so that $\Phi(n,n)=1$. The solution of $y_{n+1}=a_ny_n+b_n$ with the initial value $y_0$ is
\begin{equation}
y_n=\Phi(n,0)\,y_0+\sum_{m=0}^{n-1}\Phi(n,m+1)\,b_m\qquad(n\in\mathbb N).
\label{eq:first-order-solution}
\end{equation}
\end{theorem}

\begin{proof}
For $n=0$ the formula gives $y_0$. If it holds for $n$, then, since $a_n\Phi(n,m)=\Phi(n+1,m)$,
\[
a_ny_n+b_n=\Phi(n+1,0)\,y_0+\sum_{m=0}^{n-1}\Phi(n+1,m+1)\,b_m+\Phi(n+1,n+1)\,b_n ,
\]
which is the formula for $n+1$.
\end{proof}

The number $\Phi(n,m)$ is the factor by which the homogeneous recurrence carries a value from the index $m$ to the index $n$, the discrete counterpart of the transition matrix $\Phi(t,s)$ of \cref{def:transition-matrix}, and the formula reads like Duhamel's formula: the input $b_m$, received at the step from $m$ to $m+1$, is carried forward to the index $n$ by the homogeneous recurrence, and the solution is the sum of all these contributions. If no $a_j$ vanishes, dividing the recurrence by $\Phi(n+1,0)$ turns it into $z_{n+1}=z_n+b_n/\Phi(n+1,0)$ for $z_n\coloneqq y_n/\Phi(n,0)$, which is solved by summation; $1/\Phi(n+1,0)$ is a \term{summing factor}, as $e^{\int p}$ was an integrating factor.

The derangement numbers of \vcref{1}{thm:derangements} satisfy $D_{n+1}=(n+1)D_n+(-1)^{n+1}$ with $D_0=1$. Here $a_n=n+1$ and $\Phi(n,m)=n!/m!$, so the theorem gives
\[
D_n=n!+\sum_{m=0}^{n-1}\frac{n!}{(m+1)!}(-1)^{m+1}=n!\sum_{j=0}^n\frac{(-1)^j}{j!},
\]
the formula found there by inclusion and exclusion. Horner's scheme (\vcref{1}{alg:from-base}) is a first-order recurrence too: for $p(x)=c_dx^d+\cdots+c_0$, the numbers $s_0\coloneqq c_d$ and $s_{j+1}\coloneqq xs_j+c_{d-1-j}$ satisfy the recurrence with $a_j=x$ and $b_j=c_{d-1-j}$, and \eqref{eq:first-order-solution} gives $s_d=x^dc_d+\sum_{m<d}x^{d-1-m}c_{d-1-m}=p(x)$. The formula also shows how an error made at step $m$ propagates: it is multiplied by $x^{d-1-m}$, which matters when $|x|>1$ and which Volume~VI analyses.

\subsection{Cobweb Diagrams}

A recurrence $y_{n+1}=f(y_n)$ with a real function $f$ can be followed in the plane. Draw the graph of $f$ and the diagonal $y=x$. From the point $(y_0,0)$ go vertically to the graph, reaching $(y_0,y_1)$, then horizontally to the diagonal, reaching $(y_1,y_1)$, then vertically to the graph again, reaching $(y_1,y_2)$, and so on. The fixed points are the intersections of the graph with the diagonal, and the resulting \term{cobweb diagram} shows at a glance whether the iterates approach a fixed point and how. \Cref{fig:cobweb-linear} shows the four cases of \cref{prop:first-order-constant} for straight lines through the fixed point $(1,1)$: a staircase that climbs towards the fixed point for $a=1/2$, a spiral that closes in on it for $a=-1/2$, and the same two shapes running away from it for $a=3/2$ and $a=-3/2$.

\begin{figure}
\sdpanel{A}{\begin{tikzpicture}
\begin{axis}[sdplot, width=0.46\textwidth, height=0.38\textwidth, xmin=0, xmax=2, ymin=0, ymax=2,
  xtick={0,1,2}, ytick={0,1,2}, xlabel={$y_n$}, ylabel={$y_{n+1}$}, clip=true]
\addplot[black!45, line width=0.4pt, domain=0:2, samples=2] {x};
\addplot[black, line width=0.7pt, domain=0:2, samples=2] {0.5*(x-1)+1};
\addplot[black, line width=0.4pt] coordinates {(0.1500,0.0000) (0.1500,0.5750) (0.5750,0.5750) (0.5750,0.7875) (0.7875,0.7875) (0.7875,0.8938) (0.8938,0.8938) (0.8938,0.9469) (0.9469,0.9469) (0.9469,0.9734) (0.9734,0.9734) (0.9734,0.9867) (0.9867,0.9867) (0.9867,0.9934) (0.9934,0.9934) (0.9934,0.9967) (0.9967,0.9967)};
\end{axis}
\end{tikzpicture}}\hfill
\sdpanel{B}{\begin{tikzpicture}
\begin{axis}[sdplot, width=0.46\textwidth, height=0.38\textwidth, xmin=0, xmax=2, ymin=0, ymax=2,
  xtick={0,1,2}, ytick={0,1,2}, xlabel={$y_n$}, ylabel={$y_{n+1}$}, clip=true]
\addplot[black!45, line width=0.4pt, domain=0:2, samples=2] {x};
\addplot[black, line width=0.7pt, domain=0:2, samples=2] {-0.5*(x-1)+1};
\addplot[black, line width=0.4pt] coordinates {(0.1500,0.0000) (0.1500,1.4250) (1.4250,1.4250) (1.4250,0.7875) (0.7875,0.7875) (0.7875,1.1062) (1.1062,1.1062) (1.1062,0.9469) (0.9469,0.9469) (0.9469,1.0266) (1.0266,1.0266) (1.0266,0.9867) (0.9867,0.9867) (0.9867,1.0066) (1.0066,1.0066) (1.0066,0.9967) (0.9967,0.9967)};
\end{axis}
\end{tikzpicture}}\sdpanelrow
\sdpanel{C}{\begin{tikzpicture}
\begin{axis}[sdplot, width=0.46\textwidth, height=0.38\textwidth, xmin=0, xmax=2, ymin=0, ymax=2,
  xtick={0,1,2}, ytick={0,1,2}, xlabel={$y_n$}, ylabel={$y_{n+1}$}, clip=true]
\addplot[black!45, line width=0.4pt, domain=0:2, samples=2] {x};
\addplot[black, line width=0.7pt, domain=0:2, samples=2] {1.5*(x-1)+1};
\addplot[black, line width=0.4pt] coordinates {(0.8500,0.0000) (0.8500,0.7750) (0.7750,0.7750) (0.7750,0.6625) (0.6625,0.6625) (0.6625,0.4937) (0.4937,0.4937) (0.4937,0.2406) (0.2406,0.2406) (0.2406,-0.1391) (-0.1391,-0.1391)};
\end{axis}
\end{tikzpicture}}\hfill
\sdpanel{D}{\begin{tikzpicture}
\begin{axis}[sdplot, width=0.46\textwidth, height=0.38\textwidth, xmin=0, xmax=2, ymin=0, ymax=2,
  xtick={0,1,2}, ytick={0,1,2}, xlabel={$y_n$}, ylabel={$y_{n+1}$}, clip=true]
\addplot[black!45, line width=0.4pt, domain=0:2, samples=2] {x};
\addplot[black, line width=0.7pt, domain=0:2, samples=2] {-1.5*(x-1)+1};
\addplot[black, line width=0.4pt] coordinates {(0.8500,0.0000) (0.8500,1.2250) (1.2250,1.2250) (1.2250,0.6625) (0.6625,0.6625) (0.6625,1.5063) (1.5063,1.5063) (1.5063,0.2406) (0.2406,0.2406) (0.2406,2.1391) (2.1391,2.1391)};
\end{axis}
\end{tikzpicture}}
\caption{Cobweb diagrams for $y_{n+1}=a(y_n-1)+1$, with the fixed point $1$; the graph of the right side is the thick line, the diagonal the grey line. (A) $a=1/2$ and (B) $a=-1/2$, both from $y_0=0.15$: the iterates approach $1$ monotonically and alternately. (C) $a=3/2$ and (D) $a=-3/2$, both from $y_0=0.85$: the iterates move away from $1$ monotonically and alternately (\cref{prop:first-order-constant}).\label{fig:cobweb-linear}}
\end{figure}

For a nonlinear $f$ the diagram behaves near a fixed point like that of the tangent line there, whose slope is $f'(y^*)$. \Vcref{1}{prop:fixed-point-contraction} turned this into a theorem for contractions of an interval, and \cref{sec:nonlinear-recurrences} returns to it. The comparison with linear recurrences also underlies the discrete Gr\"onwall inequality (\cref{lem:discrete-gronwall}): a sequence that satisfies $\varepsilon_{n+1}\le a\varepsilon_n+b$ with $a\ge0$ is bounded by the solution of the corresponding equation, since the inequality is preserved under multiplication by the nonnegative $a$.

\section{Linear Recurrences with Constant Coefficients}
\label{sec:constant-coefficient-recurrences}

\subsection{The Space of Solutions}

The solutions of a homogeneous linear differential equation of order $k$ form a space of dimension $k$ (\cref{thm:solution-space}), and the proof needed the existence theorem of \cref{ch:linear-higher-order}. For recurrences the same statement is immediate, and it holds for variable coefficients as well.

\begin{theorem}[The space of solutions]
\label{thm:recurrence-solution-space}
Let $a_0,\ldots,a_{k-1}$ and $g$ be complex sequences, $k\ge1$. The solutions of
\begin{equation}
y_{n+k}+a_{k-1}(n)\,y_{n+k-1}+\cdots+a_0(n)\,y_n=0\qquad(n\in\mathbb N)
\label{eq:linear-recurrence-homogeneous}
\end{equation}
form a subspace $\mathcal S$ of $\mathbb C^{\mathbb N}$ of dimension $k$, and the map $y\mapsto(y_0,\ldots,y_{k-1})$ is an isomorphism from $\mathcal S$ onto $\mathbb C^k$. The solutions of the equation with the right side $g_n$ are the sequences $y_{\mathrm p}+z$, where $y_{\mathrm p}$ is any one of them and $z\in\mathcal S$.
\end{theorem}

\begin{proof}
The left side of \eqref{eq:linear-recurrence-homogeneous} is a linear map $L$ of $y$, so its kernel $\mathcal S$ is a subspace, and the initial-value map is linear. By \cref{prop:recurrence-ivp} it is bijective: every vector of initial values belongs to exactly one solution. So it is an isomorphism (\cref{sec:isomorphisms}), and $\dim\mathcal S=k$. If $Ly_{\mathrm p}=g$, then $Ly=g$ if and only if $L(y-y_{\mathrm p})=0$.
\end{proof}

For the Fibonacci recurrence this is the two-dimensional space of \cref{prop:dimension-examples}\,(ii). From now on in this section the coefficients are constants. Writing
\begin{equation}
p(\lambda)\coloneqq\lambda^k+c_{k-1}\lambda^{k-1}+\cdots+c_1\lambda+c_0 ,
\label{eq:recurrence-char}
\end{equation}
the recurrence $y_{n+k}+c_{k-1}y_{n+k-1}+\cdots+c_0y_n=0$ is $p(E)y=0$, and $p$ is its \term{characteristic polynomial}. The theory now follows that of $p(D)y=0$ in \cref{sec:constant-coefficients} step by step, with the shift $E$ in the place of $D$ and the geometric sequences $r^n$ in the place of the exponentials.

\begin{lemma}[Polynomials in $E$]
\label{lem:polynomials-in-e}
Let $q$ and $s$ be polynomials with complex coefficients and $r\in\mathbb C$.
\begin{enumerate}
\item $(qs)(E)=q(E)\,s(E)=s(E)\,q(E)$.
\item $q(E)\,r^n=q(r)\,r^n$: the sequence $(r^n)$ is an eigenvector of every $q(E)$, with the eigenvalue $q(r)$.
\item \emph{Shift formula}: $q(E)\bigl(r^nu_n\bigr)=r^n\,\bigl(q(rE)u\bigr)_n$ for every sequence $(u_n)$, where $q(rE)$ is the operator of the polynomial $\mu\mapsto q(r\mu)$.
\end{enumerate}
\end{lemma}

\begin{proof}
(i) Both sides are linear in $q$ and in $s$, and for monomials both are $E^{i+j}$, since $E^iE^j=E^{i+j}$. (ii) $E^jr^n=r^{n+j}=r^j\cdot r^n$. (iii) $E^j(r^nu_n)=r^{n+j}u_{n+j}=r^n\bigl((rE)^ju\bigr)_n$, and $q(E)$ is a combination of the $E^j$.
\end{proof}

Part (iii) is the discrete form of the exponential shift of \cref{lem:polynomials-in-d}\,(iii), with multiplication instead of addition: $e^{\lambda t}$ shifts $D$ to $D+\lambda$, and $r^n$ scales $E$ to $rE$. With it we find the solutions belonging to a single root.

\begin{proposition}[Kernel of a power of $E-r$]
\label{prop:kernel-power-e}
Let $m\ge1$.
\begin{enumerate}
\item If $r\ne0$, the solutions of $(E-r)^my=0$ are exactly the sequences $y_n=P(n)\,r^n$ with a polynomial $P$ of degree less than $m$.
\item The solutions of $E^my=0$ are exactly the sequences with $y_n=0$ for all $n\ge m$.
\end{enumerate}
\end{proposition}

\begin{proof}
(i) Every sequence can be written as $y_n=r^nu_n$ with $u_n\coloneqq r^{-n}y_n$. By \cref{lem:polynomials-in-e}\,(iii) with $q(\mu)=(\mu-r)^m$, the polynomial $q(r\mu)=r^m(\mu-1)^m$, so
\[
(E-r)^m(r^nu_n)=r^n\,r^m\bigl((E-I)^mu\bigr)_n=r^{n+m}\,\Delta^mu_n .
\]
Since $r\ne0$, $y$ is a solution if and only if $\Delta^mu=0$, and by \cref{thm:newton-forward}\,(ii) this holds if and only if $u_n=P(n)$ with a polynomial of degree at most $m-1$. (ii) $(E^my)_n=y_{n+m}$.
\end{proof}

As for differential equations, we need that the sequences belonging to different roots are independent. The proof imitates that of \cref{lem:exp-polynomials}, with differences in the place of derivatives.

\begin{lemma}[Independence of exponential polynomial sequences]
\label{lem:exp-poly-sequences}
Let $r_1,\ldots,r_s$ be pairwise different nonzero complex numbers, let $P_1,\ldots,P_s$ be polynomials, and let $N_0\in\mathbb N$. If
\[
P_1(n)\,r_1^n+\cdots+P_s(n)\,r_s^n=0\qquad\text{for all }n\ge N_0 ,
\]
then $P_1=\cdots=P_s=0$.
\end{lemma}

\begin{proof}
By induction on $s$. For $s=1$, $P_1(n)=0$ for infinitely many $n$, since $r_1^n\ne0$, and so $P_1=0$ (\vcref{1}{cor:factor}). Let $s\ge2$, and let $d$ be the degree of $P_s$ ($d=0$ if $P_s=0$). Dividing by $r_s^n$ gives $\sum_{i<s}P_i(n)\rho_i^n+P_s(n)=0$ for $n\ge N_0$, with $\rho_i\coloneqq r_i/r_s$, which are pairwise different, nonzero and different from $1$. For a polynomial $Q$ and $\rho\ne0$,
\[
\Delta\bigl(\rho^nQ(n)\bigr)=\rho^n\bigl(\rho\,Q(n+1)-Q(n)\bigr),
\]
and if $\rho\ne1$ and $Q\ne0$, the polynomial $\rho Q(x+1)-Q(x)$ has the same degree as $Q$ and the leading coefficient multiplied by $\rho-1\ne0$. Apply $\Delta^{d+1}$ to the relation; a sequence that vanishes for $n\ge N_0$ has a difference that vanishes for $n\ge N_0$. The term $P_s(n)$ disappears by \cref{thm:newton-forward}\,(ii), and we obtain $\sum_{i<s}\rho_i^n\tilde P_i(n)=0$ for $n\ge N_0$, where each $\tilde P_i$ is nonzero if $P_i$ is. By the induction hypothesis all $\tilde P_i$ vanish, so $P_i=0$ for $i<s$, and the relation reduces to $P_s(n)=0$ for $n\ge N_0$, so $P_s=0$.
\end{proof}

\subsection{The General Solution}

\begin{theorem}[Linear recurrences with constant coefficients]
\label{thm:constant-coefficient-recurrences}
Let the characteristic polynomial \eqref{eq:recurrence-char} factor as
\[
p(\lambda)=\lambda^{m_0}(\lambda-r_1)^{m_1}\cdots(\lambda-r_s)^{m_s},\qquad m_0+m_1+\cdots+m_s=k,
\]
with pairwise different nonzero roots $r_1,\ldots,r_s$ and $m_0\ge0$ the multiplicity of the root $0$. Then the $k$ sequences
\[
\bigl(n^jr_i^n\bigr)_{n\ge0}\quad(1\le i\le s,\ 0\le j<m_i)\qquad\text{and}\qquad\mathbf e_l=(\delta_{n,l})_{n\ge0}\quad(0\le l<m_0)
\]
form a basis of the space of solutions of $p(E)y=0$, where $\delta_{n,l}$ is $1$ for $n=l$ and $0$ otherwise. Thus the solutions are exactly the sequences
\[
y_n=\sum_{i=1}^sP_i(n)\,r_i^n+\sum_{l<m_0}d_l\,\delta_{n,l}
\]
with polynomials $P_i$ of degree less than $m_i$ and numbers $d_l$, and in particular $y_n=\sum_iP_i(n)r_i^n$ for all $n\ge m_0$. The polynomials and the numbers are determined by $y$.
\end{theorem}

\begin{proof}
Each of the $k$ sequences is a solution. For the root $r_i$ write $p=q_i\cdot(\lambda-r_i)^{m_i}$; by \cref{lem:polynomials-in-e}\,(i), $p(E)y=q_i(E)\bigl((E-r_i)^{m_i}y\bigr)$, which vanishes for $y_n=n^jr_i^n$ with $j<m_i$ by \cref{prop:kernel-power-e}\,(i). Similarly $p=q_0\cdot\lambda^{m_0}$, and $E^{m_0}\mathbf e_l=0$ for $l<m_0$ by \cref{prop:kernel-power-e}\,(ii). The sequences are independent: if a combination $\sum_iP_i(n)r_i^n+\sum_ld_l\delta_{n,l}$ vanishes for all $n$, then for $n\ge m_0$ the second sum is zero, so all $P_i$ vanish by \cref{lem:exp-poly-sequences} with $N_0=m_0$, and then $d_l=0$ by evaluating at $n=l$. Their number is $k$, the dimension of the space of solutions (\cref{thm:recurrence-solution-space}), so they form a basis (\cref{thm:basis-extension}).
\end{proof}

If $c_0\ne0$, the root $0$ does not occur, and the theorem is \cref{cor:linear-recurrence}, which \cref{ch:canonical-forms} proved by a different route, through the Jordan form of the companion matrix and a count of dimensions; the present proof verifies directly that each sequence is a solution, as \cref{thm:constant-coefficients} did for differential equations. The theorem also covers the case $c_0=0$, which \cref{cor:linear-recurrence} excluded. Then the recurrence $y_{n+k}+\cdots+c_{m_0}y_{n+m_0}=0$ does not involve $y_n,\ldots,y_{n+m_0-1}$; it is a recurrence of order $k-m_0$ for the shifted sequence $(y_{n+m_0})_{n\ge0}$, and the first $m_0$ terms are free. The recurrence $y_{n+2}=y_{n+1}$, with $p(\lambda)=\lambda^2-\lambda$, has the solutions $y_0$ arbitrary and $y_n=y_1$ for $n\ge1$, spanned by the constant sequence $1^n$ and by $\mathbf e_0$. A root $0$ of $p$ is what makes a recurrence forget: it has no counterpart for differential equations, whose characteristic polynomials may well have the root $0$ but whose corresponding solutions $t^j$ do not disappear.

The roles of the roots are those of \cref{thm:constant-coefficients}. A simple root $r$ contributes the geometric sequence $r^n$; a root of multiplicity $m$ contributes $r^n,nr^n,\ldots,n^{m-1}r^n$, the factors $n^j$ being, as \cref{ch:canonical-forms} put it, the shadow of a Jordan block of the companion matrix. The recurrence $y_{n+2}-4y_{n+1}+4y_n=0$ of \vcref{1}{sec:recurrences} has $p(\lambda)=(\lambda-2)^2$ and the solutions $(\alpha+\beta n)2^n$. The recurrence $y_{n+3}-3y_{n+2}+3y_{n+1}-y_n=0$ is $\Delta^3y=0$, with $p(\lambda)=(\lambda-1)^3$, and its solutions are the quadratic polynomials in $n$, in accordance with \cref{thm:newton-forward}.

\subsection{Real Solutions}

For real coefficients, nonreal roots come in conjugate pairs of equal multiplicity, by the argument of \cref{cor:real-fundamental}, and the corresponding complex sequences can be replaced by real ones. Write a nonreal root in polar form, $r=\rho e^{i\theta}$ with $\rho>0$ and $0<\theta<\pi$ (\vcref{1}{sec:complex}).

\begin{corollary}[Real solutions]
\label{cor:real-recurrence-solutions}
Let $p$ have real coefficients. If $\rho e^{\pm i\theta}$, with $\rho>0$ and $0<\theta<\pi$, is a pair of roots of multiplicity $m$, then in the basis of \cref{thm:constant-coefficient-recurrences} the $2m$ sequences $n^j\rho^ne^{\pm in\theta}$ may be replaced by the $2m$ real sequences
\[
n^j\rho^n\cos n\theta,\qquad n^j\rho^n\sin n\theta\qquad(0\le j<m),
\]
and the resulting real sequences form a basis of the real solutions.
\end{corollary}

\begin{proof}
By de Moivre's theorem (\vcref{1}{thm:de-moivre}), $(\rho e^{\pm i\theta})^n=\rho^n(\cos n\theta\pm i\sin n\theta)$, so the two complex sequences and the two real sequences of the statement span the same complex space, for each $j$. The real sequences are solutions, since the real and imaginary parts of a solution of an equation with real coefficients are solutions, and they span the space of complex solutions, which has dimension $k$; so they are independent. Every real solution is a complex combination of them, and taking real parts shows that it is a real combination; so they form a basis of the real solutions as well.
\end{proof}

A pair of roots $\rho e^{\pm i\theta}$ produces oscillations with the angle $\theta$ per step and an amplitude that grows or decays like $\rho^n$: the modulus of the root gives the rate of growth, its argument the frequency, as the real and imaginary parts of $\lambda$ did for $p(D)y=0$. \Cref{fig:real-solutions} shows the three cases. A real negative root $r<0$ produces the alternating sequence $r^n=(-1)^n|r|^n$, which has no counterpart among the real solutions of differential equations with constant coefficients, where a real root $\lambda$ gives the positive function $e^{\lambda t}$. The recurrence $y_{n+2}-y_{n+1}+y_n=0$ of \vcref{1}{sec:recurrences}, with the roots $e^{\pm i\pi/3}$, has the real solutions $\cos(n\pi/3)$ and $\sin(n\pi/3)$, periodic with period $6$, and the solution found there, $(2/\sqrt3)\sin(n\pi/3)$, is a combination of them.

\begin{figure}
\sdfloat{%
\begin{tikzpicture}
\begin{axis}[sdplot, width=0.62\textwidth, height=0.40\textwidth, xmin=-0.5, xmax=40.5, ymin=-3.4, ymax=3.4,
  xtick={0,10,20,30,40}, ytick={-3,-2,-1,0,1,2,3}, xlabel={$n$}, legend pos=north west, legend columns=3]
\addplot[black!35, line width=0.3pt, samples at={0,...,40}, mark=*, mark size=1.1pt, mark options={black}] {0.9^x*cos(deg(x*pi/6))};
\addlegendentry{$\rho=0.9$}
\addplot[black!35, line width=0.3pt, samples at={0,...,40}, mark=o, mark size=1.2pt, mark options={black}] {cos(deg(x*pi/6))};
\addlegendentry{$\rho=1$}
\addplot[black!35, line width=0.3pt, samples at={0,...,40}, mark=square*, mark size=0.9pt, mark options={black!60}] {1.03^x*cos(deg(x*pi/6))};
\addlegendentry{$\rho=1.03$}
\end{axis}
\end{tikzpicture}}{The real solutions $y_n=\rho^n\cos(n\pi/6)$ of $y_{n+2}-2\rho\cos(\pi/6)\,y_{n+1}+\rho^2y_n=0$, whose characteristic roots are $\rho e^{\pm i\pi/6}$, for $\rho=0.9$ (filled circles), $1$ (open circles) and $1.03$ (squares), $n=0,\ldots,40$; consecutive terms are joined by thin lines. The oscillation repeats every $12$ steps, and its amplitude decays, stays constant or grows according as the modulus $\rho$ of the roots is less than, equal to or greater than $1$ (\cref{cor:real-recurrence-solutions}).\label{fig:real-solutions}}
\end{figure}

\subsection{The Casorati Determinant}

The Wronskian decides whether $k$ solutions of a differential equation form a fundamental set. For recurrences its role is played by the determinant of consecutive terms.

\begin{definition}[Casorati determinant]
\label{def:casorati}
For sequences $y^{(1)},\ldots,y^{(k)}$, the \term{Casorati matrix} and the \term{Casorati determinant} are
\[
K(n)\coloneqq\begin{bmatrix}y^{(1)}_n&y^{(2)}_n&\cdots&y^{(k)}_n\\y^{(1)}_{n+1}&y^{(2)}_{n+1}&\cdots&y^{(k)}_{n+1}\\\vdots&\vdots&&\vdots\\y^{(1)}_{n+k-1}&y^{(2)}_{n+k-1}&\cdots&y^{(k)}_{n+k-1}\end{bmatrix},\qquad C(n)\coloneqq\det K(n).
\]
\end{definition}

\begin{theorem}[Casorati determinant of solutions]
\label{thm:casorati}
Let $y^{(1)},\ldots,y^{(k)}$ be solutions of the recurrence \eqref{eq:linear-recurrence-homogeneous} with variable coefficients, and let $C(n)$ be their Casorati determinant.
\begin{enumerate}
\item $C(n+1)=(-1)^ka_0(n)\,C(n)$ for all $n$, and therefore
\[
C(n)=(-1)^{kn}\,a_0(0)\,a_0(1)\cdots a_0(n-1)\,C(0) .
\]
\item The solutions are linearly independent if and only if $C(0)\ne0$.
\item If $a_0(n)\ne0$ for all $n$, then $C(n)$ either vanishes for all $n$ or for no $n$.
\end{enumerate}
\end{theorem}

\begin{proof}
(i) The $j$th column of $K(n+1)$ is obtained from that of $K(n)$ by dropping the first entry and appending $y^{(j)}_{n+k}=-\sum_{i<k}a_i(n)y^{(j)}_{n+i}$. So $K(n+1)=A(n)K(n)$, where $A(n)$ has the rows $\mathbf e_2^{\mathsf T},\ldots,\mathbf e_k^{\mathsf T}$ and $(-a_0(n),\ldots,-a_{k-1}(n))$; it is the transposed companion matrix $C_q^{\mathsf T}$ of $q(\lambda)=\lambda^k+a_{k-1}(n)\lambda^{k-1}+\cdots+a_0(n)$ (\cref{prop:companion}). Its characteristic polynomial is $q$, so $(-1)^k\det A(n)=\det(-A(n))=q(0)=a_0(n)$, and $\det A(n)=(-1)^ka_0(n)$. The product rule for determinants (\cref{thm:det-product}) gives $C(n+1)=(-1)^ka_0(n)C(n)$, and induction gives the formula. (ii) The columns of $K(0)$ are the vectors of initial values of the solutions, and by \cref{thm:recurrence-solution-space} the solutions are independent if and only if these vectors are, that is, if and only if $\det K(0)\ne0$ (\cref{thm:det-invertible}). (iii) follows from the formula in (i).
\end{proof}

Part (i) is the discrete form of Abel's formula (\cref{thm:abel}): the factor $(-1)^ka_0(n)$ per step takes the place of $\exp(-\int a_{k-1})$, and it involves the coefficient of the lowest term instead of that of the second-highest derivative. Part (iii) differs from \cref{thm:wronskian-solutions} in its hypothesis. When $a_0(n_1)=0$ for some $n_1$, independent solutions can have a Casorati determinant that vanishes from $n_1+1$ on; for the recurrence $y_{n+2}=y_{n+1}$, the independent solutions $1^n$ and $\mathbf e_0$ have the Casorati matrices with the rows $(1,1)$, $(1,0)$ at $n=0$ and $(1,0)$, $(1,0)$ for $n\ge1$, so that $C(0)=-1$ and $C(n)=0$ for $n\ge1$, the trace in the determinant of the information that the recurrence forgets.

The Fibonacci numbers and the \term{Lucas numbers} $L_n$, defined by $L_0=2$, $L_1=1$ and the Fibonacci recurrence, are two solutions of $y_{n+2}-y_{n+1}-y_n=0$, with $a_0=-1$; by \cref{thm:constant-coefficient-recurrences} and the values at $n=0,1$,
\[
F_n=\frac{\phi^n-\hat\phi^n}{\sqrt5},\qquad L_n=\phi^n+\hat\phi^n ,
\]
so that $L_n$ is the integer nearest to $\phi^n$ for $n\ge2$, since $|\hat\phi|^n<1/2$ then. The theorem gives $C(n)=(-1)^n C(0)$ for any two solutions. For $F$ and $L$, $C(0)=F_0L_1-F_1L_0=-2$, so $F_nL_{n+1}-F_{n+1}L_n=2(-1)^{n+1}$. For the two solutions $F_n$ and $F_{n+1}$, $C(0)=F_0F_2-F_1^2=-1$, which gives \term{Cassini's identity} $F_nF_{n+2}-F_{n+1}^2=(-1)^{n+1}$, or $F_{n-1}F_{n+1}-F_n^2=(-1)^n$; it is also the determinant of the matrix power $F^n$ of \cref{sec:matrix-multiplication}, whose determinant is $(\det F)^n=(-1)^n$.

\subsection{Chebyshev Polynomials and Tridiagonal Determinants}

Many recurrences contain a parameter, and then their solutions are functions of it. The \term{Chebyshev polynomials} of \cref{tab:orthogonal-polynomials} are defined by
\begin{equation}
T_0(x)=1,\qquad T_1(x)=x,\qquad T_{n+1}(x)=2x\,T_n(x)-T_{n-1}(x)\quad(n\ge1),
\label{eq:chebyshev-recurrence}
\end{equation}
and the \term{Chebyshev polynomials of the second kind} $U_n$ by the same recurrence with $U_0(x)=1$ and $U_1(x)=2x$. By induction, $T_n$ and $U_n$ are polynomials in $x$ of degree $n$, with the leading coefficients $2^{n-1}$ ($n\ge1$) and $2^n$. For a fixed $x$, the recurrence in $n$ has constant coefficients and the characteristic polynomial $\lambda^2-2x\lambda+1$.

\begin{proposition}[Closed forms of the Chebyshev polynomials]
\label{prop:chebyshev-closed}
For $0\le\theta\le\pi$ and $n\in\mathbb N$,
\[
T_n(\cos\theta)=\cos n\theta,\qquad U_n(\cos\theta)=\frac{\sin(n+1)\theta}{\sin\theta}\quad(0<\theta<\pi),
\]
and for real $x$ with $|x|\ge1$, with $\lambda_\pm\coloneqq x\pm\sqrt{x^2-1}$,
\[
T_n(x)=\frac{\lambda_+^n+\lambda_-^n}2 .
\]
In particular $T_n(x)=\cos(n\arccos x)$ and $|T_n(x)|\le1$ for $|x|\le1$, while $T_n(x)\ge\lambda_+^n/2$ grows exponentially for $x>1$.
\end{proposition}

\begin{proof}
For $x=\cos\theta$ with $0<\theta<\pi$, the characteristic roots are $e^{\pm i\theta}$, which are different, and by \cref{cor:real-recurrence-solutions} the solutions are $A\cos n\theta+B\sin n\theta$. The initial values $1$ and $\cos\theta$ give $A=1$ and $B=0$; the initial values $1$ and $2\cos\theta$ give $A=1$ and $B\sin\theta=\cos\theta$, and then $\cos n\theta+\cos\theta\sin n\theta/\sin\theta=\sin(n+1)\theta/\sin\theta$ by the addition formula (\vcref{1}{thm:addition}). For $\theta=0$ and $\theta=\pi$, $x=\pm1$ is a double root, the solutions are $(A+Bn)(\pm1)^n$, and the initial values give $T_n(\pm1)=(\pm1)^n=\cos n\theta$. For $|x|>1$ the roots $\lambda_\pm$ are real and different, with $\lambda_++\lambda_-=2x$ and $\lambda_+\lambda_-=1$, and $T_n=(\lambda_+^n+\lambda_-^n)/2$ has the right initial values. For $x>1$, $\lambda_+>1>\lambda_->0$, which gives the lower bound.
\end{proof}

The proposition explains the shape of the Chebyshev polynomials: on $[-1,1]$ they are cosines in disguise, oscillating between $-1$ and $1$, and outside it they grow like $\lambda_+^n$. \Cref{sec:recurrence-computation} compares the three ways of evaluating $T_n(x)$ that we now have: the recurrence, the formula $\cos(n\arccos x)$ and the coefficients of the polynomial.

The determinants of \cref{prop:tridiagonal-det} satisfy a recurrence of the same kind. For a tridiagonal matrix with constant entries $a$ on the diagonal, $b$ above and $c$ below it, the determinants $D_n$ of order $n$ satisfy
\[
D_n=a\,D_{n-1}-bc\,D_{n-2}\quad(n\ge2),\qquad D_0=1,\quad D_1=a ,
\]
and with $D_{-1}\coloneqq0$ the recurrence holds also for $n=1$.

\begin{proposition}[Tridiagonal Toeplitz matrices]
\label{prop:tridiagonal-closed}
Let $a,b,c\in\mathbb C$ with $bc\ne0$, let $\lambda_1,\lambda_2$ be the roots of $\lambda^2-a\lambda+bc$, and let $D_n$ be the determinant of the $n\times n$ tridiagonal matrix $T(a,b,c)$ with $a$ on the diagonal, $b$ above and $c$ below it. Then
\[
D_n=\frac{\lambda_1^{n+1}-\lambda_2^{n+1}}{\lambda_1-\lambda_2}\quad\text{if }\lambda_1\ne\lambda_2,\qquad D_n=(n+1)\Bigl(\frac a2\Bigr)^n\quad\text{if }\lambda_1=\lambda_2 .
\]
Moreover, if $\beta$ is a square root of $bc$, the eigenvalues of $T(a,b,c)$ are the $n$ different numbers
\[
\mu_j=a+2\beta\cos\frac{j\pi}{n+1}\qquad(j=1,\ldots,n).
\]
\end{proposition}

\begin{proof}
The sequence $(D_n)_{n\ge-1}$ satisfies the recurrence with the characteristic polynomial $\lambda^2-a\lambda+bc$. If the roots are different, the sequence $(\lambda_1^{n+1}-\lambda_2^{n+1})/(\lambda_1-\lambda_2)$ is a solution by \cref{thm:constant-coefficient-recurrences} (shifted by one index) and has the values $0$ and $1$ at $n=-1$ and $n=0$, so it is $D_n$. If $\lambda_1=\lambda_2=a/2$, the solution $(n+1)(a/2)^n$ has the same values. For the eigenvalues, $\det(\mu I-T(a,b,c))$ is the determinant of $T(\mu-a,-b,-c)$, whose off-diagonal product is again $bc$. For $\mu=\mu_j$ and $\theta_j\coloneqq j\pi/(n+1)$, the roots of $\lambda^2-(\mu_j-a)\lambda+bc=\lambda^2-2\beta\cos\theta_j\,\lambda+\beta^2$ are $\beta e^{\pm i\theta_j}$, which are different because $0<\theta_j<\pi$, and the first formula gives
\[
\det(\mu_jI-T)=\frac{\beta^{n+1}\bigl(e^{i(n+1)\theta_j}-e^{-i(n+1)\theta_j}\bigr)}{\beta\bigl(e^{i\theta_j}-e^{-i\theta_j}\bigr)}=\beta^n\frac{\sin(n+1)\theta_j}{\sin\theta_j}=0 ,
\]
since $(n+1)\theta_j=j\pi$. The numbers $\mu_j$ are different, because the cosine is injective on $(0,\pi)$ and $\beta\ne0$. So the characteristic polynomial, of degree $n$, has the $n$ different roots $\mu_1,\ldots,\mu_n$, and these are all the eigenvalues.
\end{proof}

For the second difference matrix $T_n$ of \cref{sec:special-matrices}, $a=2$ and $b=c=-1$, so $bc=1$ and $\lambda^2-2\lambda+1$ has the double root $1$: $D_n=n+1$, the value $\det T_n=n+1$ found in \cref{ch:determinants}. With $\beta=-1$ the eigenvalues are $2-2\cos(j\pi/(n+1))$, those of \cref{thm:second-difference-eigen}, now obtained from a recurrence instead of a guess. For $a=1$, $b=1$ and $c=-1$, $D_n=F_{n+1}$, the Fibonacci numbers of \cref{prop:tridiagonal-det}, and the formula is Binet's.

\subsection{A Boundary Value Problem: The Gambler's Ruin}

Recurrences also arise with conditions at both ends. A gambler with $i$ units bets one unit at a time and wins each bet with probability $p$, loses with probability $q\coloneqq1-p$, independently of the past, and stops when she is ruined or has reached $N$ units. Let $P_i$ be the probability that she reaches $N$ before $0$, starting with $i$ units. Conditioning on the outcome of the first bet gives $P_i=pP_{i+1}+qP_{i-1}$ for $0<i<N$, together with $P_0=0$ and $P_N=1$; Volume~IV justifies this step in the general theory of Markov chains, and Feller~\cite[Chap.~XIV]{feller1968} treats the problem in detail.

\begin{proposition}[Gambler's ruin]
\label{prop:gamblers-ruin}
Let $0<p<1$, $q=1-p$ and $N\ge1$. The system $P_i=pP_{i+1}+qP_{i-1}$ for $0<i<N$, $P_0=0$, $P_N=1$, has exactly one solution,
\[
P_i=\frac{1-(q/p)^i}{1-(q/p)^N}\quad\text{if }p\ne q,\qquad P_i=\frac iN\quad\text{if }p=q=\frac12 .
\]
\end{proposition}

\begin{proof}
The equations say $pP_{i+1}-P_i+qP_{i-1}=0$, a recurrence of order $2$ with the characteristic polynomial $p\lambda^2-\lambda+q=(\lambda-1)(p\lambda-q)$, whose roots are $1$ and $q/p$. By \cref{thm:constant-coefficient-recurrences}, applied to the sequence $(P_{i})_{0\le i\le N}$, which is determined by $P_0$ and $P_1$ and may be continued beyond $N$ by the recurrence, the solutions are $A+B(q/p)^i$ if $p\ne q$ and $A+Bi$ if $p=q$. If $p\ne q$, the conditions $P_0=0$ and $P_N=1$ give $A+B=0$ and $A+B(q/p)^N=1$, a linear system with the determinant $(q/p)^N-1\ne0$; if $p=q$, they give $A=0$ and $A+BN=1$. Each system has exactly one solution, which gives the formulas.
\end{proof}

Unlike an initial value problem, a boundary value problem for a recurrence need not have a solution: the two conditions at different ends lead to a linear system for the constants, whose determinant may vanish, as it did for the differential equations of \cref{prop:boundary-value}. Here it does not. The result is sobering for gamblers (\cref{fig:gamblers-ruin}). In a fair game the probability of doubling $50$ units before losing them is $1/2$, but with $p=0.49$ it is $1/(1+(51/49)^{50})=0.119$, and at roulette, betting on red with $p=18/38$, it is $0.0051$.

\begin{figure}
\sdfloat{%
\begin{tikzpicture}
\begin{axis}[sdplot, width=0.56\textwidth, height=0.38\textwidth, xmin=0, xmax=100, ymin=0, ymax=1,
  xtick={0,25,50,75,100}, ytick={0,0.5,1}, xlabel={$i$}, ylabel={$P_i$}, legend pos=north west]
\addplot[black, line width=0.7pt, domain=0:100, samples=2] {x/100};
\addlegendentry{$p=0.5$}
\addplot[black, dashed, line width=0.7pt, domain=0:100, samples=101] {(1-(51/49)^x)/(1-(51/49)^100)};
\addlegendentry{$p=0.49$}
\addplot[black, densely dotted, line width=0.8pt, domain=0:100, samples=101] {(1-(53/47)^x)/(1-(53/47)^100)};
\addlegendentry{$p=0.47$}
\end{axis}
\end{tikzpicture}}{The probability $P_i$ of reaching $N=100$ units before ruin, starting with $i$ units, for the probabilities $p=0.5$, $0.49$ and $0.47$ of winning each bet (\cref{prop:gamblers-ruin}), drawn as curves through the values at integers $i$. A small disadvantage in each bet makes success unlikely unless the gambler starts close to the goal.\label{fig:gamblers-ruin}}
\end{figure}

\subsection{Undetermined Coefficients}

By \cref{thm:recurrence-solution-space}, the inhomogeneous equation $p(E)y=g$ is solved once one particular solution is known. For inputs that are polynomials times geometric sequences, a particular solution of the same kind can be found by solving a linear system, as for differential equations (\cref{thm:undetermined}); sines and cosines are included as real and imaginary parts of $e^{in\theta}$.

\begin{theorem}[Undetermined coefficients for recurrences]
\label{thm:discrete-undetermined}
Let $p$ be the characteristic polynomial \eqref{eq:recurrence-char}, let $\mu\in\mathbb C$ with $\mu\ne0$, let $s\ge0$ be the multiplicity of $\mu$ as a root of $p$, with $s=0$ if $p(\mu)\ne0$, and let $q$ be a polynomial of degree $d$. Then the recurrence
\[
p(E)y=q(n)\,\mu^n
\]
has exactly one solution of the form $y_n=n^sr(n)\,\mu^n$ with a polynomial $r$ of degree at most $d$, and this $r$ has the degree $d$.
\end{theorem}

\begin{proof}
Factor the characteristic polynomial as $p(\lambda)=(\lambda-\mu)^s\,p_1(\lambda)$, where $p_1(\mu)\ne0$. By the shift formula (\cref{lem:polynomials-in-e}\,(iii)), for every sequence $u$,
\[
p(E)(\mu^nu_n)=\mu^n\bigl(p(\mu E)u\bigr)_n=\mu^n\mu^s\bigl(p_1(\mu E)\,\Delta^su\bigr)_n ,
\]
since $p(\mu\lambda)=(\mu\lambda-\mu)^sp_1(\mu\lambda)=\mu^s(\lambda-1)^sp_1(\mu\lambda)$ and $E-I=\Delta$. So $y=\mu^nu$ solves the equation if and only if $M(\Delta^su)=\mu^{-s}q$, where $M\coloneqq p_1(\mu E)$. We consider sequences $u_n=U(n)$ given by polynomials $U$, on which $E$ acts as $U(x)\mapsto U(x+1)$ and $\Delta$ lowers the degree by one, as in the proof of \cref{thm:newton-forward}.

First, $M$ maps $\mathcal P_d$ bijectively onto itself and preserves degrees. Indeed, $p_1(\mu E)=p_1(\mu I+\mu\Delta)$, and expanding the polynomial $p_1(\mu+\mu\lambda)$ in powers of $\lambda$ gives $M=p_1(\mu)I+\sum_{j\ge1}\beta_j\Delta^j$ with constants $\beta_j$; so $MU$ has the degree of $U$ and the leading coefficient multiplied by $p_1(\mu)\ne0$. Thus $M$ is injective on $\mathcal P_d$, hence bijective (\cref{cor:injective-surjective}). Second, $\Delta^s$ maps the space $x^s\mathcal P_d$ of the polynomials $x^sr(x)$ with $r\in\mathcal P_d$ bijectively onto $\mathcal P_d$. It maps $\mathcal P_{s+d}$ onto $\mathcal P_d$ with the kernel $\mathcal P_{s-1}$ (\cref{thm:newton-forward}\,(ii)), and $x^s\mathcal P_d$ meets $\mathcal P_{s-1}$ only in $0$; so $\Delta^s$ is injective on $x^s\mathcal P_d$, which has the dimension $d+1$ of $\mathcal P_d$, and therefore bijective; and it maps a polynomial of degree $s+e$ to one of degree $e$.

Consequently $U\mapsto M(\Delta^sU)$ is a bijection from $x^s\mathcal P_d$ onto $\mathcal P_d$, and there is exactly one $U=x^sr$ with $r\in\mathcal P_d$ and $M(\Delta^sU)=\mu^{-s}q$. Since $M$ preserves degrees and $\Delta^s$ lowers them by $s$, the polynomial $U$ has the degree $s+d$, and $r$ the degree $d$. This $U$ gives the solution $y_n=n^sr(n)\mu^n$, and it is the only solution of this form.
\end{proof}

The factor $n^s$ is needed exactly when $\mu$ is a characteristic root, the case of resonance, in which the trial sequence $r(n)\mu^n$ would contain solutions of the homogeneous equation. Three examples show the method. For the numbers of moves of the Tower of Hanoi, $T_{n+1}-2T_n=1$: here $\mu=1$ is not a root of $\lambda-2$, the trial sequence is a constant $A$, and $A-2A=1$ gives $A=-1$; the general solution is $c2^n-1$, and $T_0=0$ gives $T_n=2^n-1$, as in \vcref{1}{sec:recurrences}. For $y_{n+2}-4y_{n+1}+4y_n=2^n$, $\mu=2$ is a double root, $s=2$, and the trial sequence $An^22^n$ gives
\[
A2^n\bigl(4(n+2)^2-8(n+1)^2+4n^2\bigr)=8A\,2^n=2^n ,
\]
so $A=1/8$, and $n^22^{n-3}$ is a particular solution. For $y_{n+1}-y_n=n^2$, whose solution with $y_0=0$ is the sum $\sum_{m<n}m^2$, $\mu=1$ is a simple root and $q(n)=n^2$, so the trial sequence is $n(an^2+bn+c)$; its difference is $a(3n^2+3n+1)+b(2n+1)+c$, and comparing coefficients gives $a=1/3$, $b=-1/2$, $c=1/6$, that is, $\sum_{m<n}m^2=(n-1)n(2n-1)/6$.

Recurrences of divide-and-conquer algorithms (\vcref{1}{sec:recurrences}) become recurrences with constant coefficients after a change of the index. For $T(n)=2T(n/2)+n$ with $n=2^m$, the numbers $S_m\coloneqq T(2^m)$ satisfy $S_{m+1}=2S_m+2^{m+1}$. Here $\mu=2$ is the root of $\lambda-2$, so the trial sequence is $Am2^m$, and $A(m+1)2^{m+1}-2Am2^m=2^{m+1}$ gives $A=1$: $S_m=(S_0+m)2^m$, that is, $T(n)=n\log_2n+T(1)\,n$. For $T(n)=3T(n/2)+n$, the recurrence of Karatsuba's multiplication, $S_{m+1}=3S_m+2^{m+1}$; now $\mu=2$ is not a root of $\lambda-3$, the trial $A2^m$ gives $2A-3A=2$, and $S_m=(S_0+2)3^m-2^{m+1}$, that is, $T(n)=(T(1)+2)\,n^{\log_23}-2n$. Volume~V develops the general theory of such recurrences, the master theorem, with the algorithms that produce them.

\section{Generating Functions}
\label{sec:recurrence-gf}

\subsection{Recurrences and Rational Functions}

\Vcref{1}{sec:gf} solved the Fibonacci recurrence by its generating function, the formal power series $\sum_nF_nx^n$, and found that it is the rational function $x/(1-x-x^2)$. The same happens for every linear recurrence with constant coefficients, and the converse holds as well. The denominator is the characteristic polynomial written backwards. For $p$ as in \eqref{eq:recurrence-char} let
\[
\tilde p(x)\coloneqq x^kp(1/x)=1+c_{k-1}x+c_{k-2}x^2+\cdots+c_0x^k ,
\]
the \term{reversed polynomial}, whose constant term is $1$, so that it is invertible in $\mathbb C[[x]]$ (\vcref{1}{prop:fps-ring}). If $p(\lambda)=\lambda^{m_0}\prod_i(\lambda-r_i)^{m_i}$ as in \cref{thm:constant-coefficient-recurrences}, then
\begin{equation}
\tilde p(x)=x^k\,x^{-m_0}\prod_i\Bigl(\frac1x-r_i\Bigr)^{m_i}=\prod_{i=1}^s(1-r_ix)^{m_i},
\label{eq:reversed-factorisation}
\end{equation}
since $m_0+\sum_im_i=k$. The nonzero roots $r_i$ of $p$ thus appear as the reciprocals $1/r_i$ of the roots of $\tilde p$, and the root $0$ of $p$ lowers the degree of $\tilde p$ to $k-m_0$.

\begin{theorem}[Rational generating functions]
\label{thm:rational-gf}
A sequence $(y_n)$ satisfies $p(E)y=0$ if and only if its generating function $Y(x)=\sum_ny_nx^n$ satisfies
\[
\tilde p(x)\,Y(x)=P(x)
\]
for a polynomial $P$ of degree less than $k$; then $Y=P/\tilde p$ in $\mathbb C[[x]]$, and the coefficients of $P$ are $\sum_{j=0}^m\tilde p_jy_{m-j}$ for $m<k$, where $\tilde p_j$ is the coefficient of $x^j$ in $\tilde p$. Consequently the sequences with a rational generating function $P/Q$, where $Q(0)\ne0$ and $\deg P<\deg Q$, are exactly the solutions of linear recurrences with constant coefficients and $c_0\ne0$.
\end{theorem}

\begin{proof}
The coefficient of $x^j$ in $\tilde p$ is $\tilde p_j=c_{k-j}$, with $c_k\coloneqq1$. By the definition of the product of formal power series, the coefficient of $x^{n+k}$ in $\tilde pY$ is
\[
\sum_{j=0}^k\tilde p_j\,y_{n+k-j}=\sum_{j=0}^kc_{k-j}\,y_{n+k-j}=\sum_{i=0}^kc_iy_{n+i}=(p(E)y)_n .
\]
So $p(E)y=0$ if and only if all coefficients of $\tilde pY$ with index at least $k$ vanish, that is, if and only if $\tilde pY$ is a polynomial of degree less than $k$; its coefficients of index $m<k$ are $\sum_{j\le m}\tilde p_jy_{m-j}$. Multiplying by the inverse of $\tilde p$ gives $Y=P/\tilde p$. For the last statement, if $Y=P/Q$ with $Q(0)\ne0$ and $\deg P<\deg Q\eqqcolon k$, divide $P$ and $Q$ by $Q(0)$; then $Q$ is the reversed polynomial $\tilde p$ of $p(\lambda)\coloneqq\lambda^kQ(1/\lambda)$, which has degree $k$, leading coefficient $1$ and $c_0\ne0$, and the first part applies. Conversely, if $c_0\ne0$, then $\tilde p$ has degree $k$.
\end{proof}

The proof makes precise the computation of \vcref{1}{sec:gf}: multiplying the generating function by $\tilde p$ applies the recurrence to all coefficients at once, and only the first $k$ coefficients, which the recurrence does not constrain, survive. Inhomogeneous recurrences are handled in the same way. For $y_{n+1}=2y_n+1$, multiplying by $x^{n+1}$ and summing over $n\ge0$ gives $Y(x)-y_0=2xY(x)+x/(1-x)$, so
\[
Y(x)=\frac{y_0}{1-2x}+\frac x{(1-x)(1-2x)}=\frac{y_0}{1-2x}+\frac1{1-2x}-\frac1{1-x},
\]
and by \veqref{1}{eq:geometric-series}, $y_n=(y_0+1)2^n-1$. The input $1$ has the generating function $1/(1-x)$, whose denominator is the reversed characteristic polynomial of the recurrence $\Delta y=0$ that annihilates it; in general, $q(n)\mu^n$ with $\deg q=d$ has a generating function with the denominator $(1-\mu x)^{d+1}$, and the generating function of the solution has the denominator $\tilde p(x)(1-\mu x)^{d+1}$, in which $1-\mu x$ occurs with the multiplicity $s+d+1$. This explains the factor $n^s$ of \cref{thm:discrete-undetermined}, as the annihilator method explained the factor $t^s$ for differential equations (\cref{sec:undetermined}).

\subsection{Partial Fractions}

The generating function of a geometric sequence is $1/(1-rx)$, and by \vcref{1}{prop:gf-rules}\,(iii) with $x$ replaced by $rx$,
\begin{equation}
\frac1{(1-rx)^j}=\sum_{n\ge0}\binom{n+j-1}{j-1}r^nx^n\qquad(j\ge1).
\label{eq:gf-power-of-geometric}
\end{equation}
The coefficient $\binom{n+j-1}{j-1}=(n+j-1)^{\underline{j-1}}/(j-1)!$ is a polynomial in $n$ of degree $j-1$. Hence the generating functions of the solutions found in \cref{thm:constant-coefficient-recurrences} are sums of terms $A/(1-r_ix)^j$, and the theorem implies the existence of partial fraction decompositions with complex coefficients, which \vcref{1}{thm:partial-fractions} established only over the real numbers.

\begin{corollary}[Partial fractions over the complex numbers]
\label{cor:complex-partial-fractions}
Let $\tilde p(x)=\prod_{i=1}^s(1-r_ix)^{m_i}$ with pairwise different nonzero $r_i$, and let $P$ be a polynomial of degree less than $k\coloneqq m_0+\sum_im_i$, for some $m_0\ge0$. Then there are unique complex numbers $A_{ij}$ and a unique polynomial $R$ of degree less than $m_0$ with
\[
\frac{P(x)}{\tilde p(x)}=R(x)+\sum_{i=1}^s\sum_{j=1}^{m_i}\frac{A_{ij}}{(1-r_ix)^j} .
\]
\end{corollary}

\begin{proof}
Let $p(\lambda)\coloneqq\lambda^{m_0}\prod_i(\lambda-r_i)^{m_i}$, a monic polynomial of degree $k$ with the reversed polynomial $\tilde p$, by \eqref{eq:reversed-factorisation}. By \cref{thm:rational-gf}, $y\mapsto\tilde pY$ maps the space $\mathcal S$ of solutions of $p(E)y=0$ into $\mathcal P_{k-1}$; it is linear and injective, since $Y=(\tilde pY)/\tilde p$, and both spaces have dimension $k$ (\cref{thm:recurrence-solution-space}), so it is bijective. The polynomials $\binom{x+j-1}{j-1}$, $1\le j\le m_i$, have the degrees $0,\ldots,m_i-1$ and so form a basis of $\mathcal P_{m_i-1}$. By \cref{thm:constant-coefficient-recurrences}, the sequences $\binom{n+j-1}{j-1}r_i^n$ and $\mathbf e_l$ ($l<m_0$) therefore form a basis of $\mathcal S$, and their generating functions are $1/(1-r_ix)^j$ and $x^l$. The solution $y$ with $\tilde pY=P$ is a unique combination of the basis, and its generating function $P/\tilde p$ is the same combination of these generating functions.
\end{proof}

Read in the other direction, the corollary gives a second proof of the general solution, the one that \vcref{1}{prop:linear-recurrence-k} sketched: the generating function of a solution is $P/\tilde p$, its partial fractions are found by algebra, and \eqref{eq:gf-power-of-geometric} turns each term into a sequence $\binom{n+j-1}{j-1}r_i^n$~\cite[Sect.~7.3]{graham1994concrete}. With an independent proof of partial fractions, such as the one over the reals in \vcref{1}{thm:partial-fractions} extended to complex polynomials, this route is self-contained; here we obtained the decomposition from the recurrence instead.

\subsection{Growth and the Radius of Convergence}

By \vcref{1}{prop:formal-analytic}, a generating function with a positive radius of convergence is also a function, and \vcref{1}{sec:power-series} computed the radius $1/\phi$ of the generating function of the Fibonacci numbers from the ratios $F_n/F_{n+1}$. For every solution of a linear recurrence the radius is determined by the roots, through the following lemma on oscillating sums.

\begin{lemma}[Sums of unimodular powers]
\label{lem:unimodular-means}
Let $\omega_1,\ldots,\omega_s$ be pairwise different complex numbers of modulus $1$ and $a_1,\ldots,a_s\in\mathbb C$. If $a_1\omega_1^n+\cdots+a_s\omega_s^n\to0$ as $n\to\infty$, then $a_1=\cdots=a_s=0$.
\end{lemma}

\begin{proof}
Let $z_n\coloneqq\sum_ia_i\omega_i^n$. Then
\[
\frac1N\sum_{n=0}^{N-1}|z_n|^2=\sum_{i,j}a_i\bar a_j\,\frac1N\sum_{n=0}^{N-1}(\omega_i\bar\omega_j)^n .
\]
For $i=j$ the inner mean is $1$, since $\omega_i\bar\omega_i=|\omega_i|^2=1$. For $i\ne j$, $\zeta\coloneqq\omega_i\bar\omega_j$ has modulus $1$ and is different from $1$, and $|\sum_{n<N}\zeta^n|=|1-\zeta^N|/|1-\zeta|\le2/|1-\zeta|$, so the mean tends to $0$. Hence the left side tends to $\sum_i|a_i|^2$. On the other hand, if $z_n\to0$, then for every $\varepsilon>0$ there is $M$ with $|z_n|^2<\varepsilon$ for $n\ge M$, and the left side is at most $(M\max_{n<M}|z_n|^2)/N+\varepsilon$, which is less than $2\varepsilon$ for large $N$. So $\sum_i|a_i|^2=0$.
\end{proof}

\begin{proposition}[Growth of solutions]
\label{prop:recurrence-growth}
Let $c_0\ne0$, and let $y_n=\sum_iP_i(n)r_i^n$ be a nonzero solution of $p(E)y=0$, written as in \cref{thm:constant-coefficient-recurrences}. Let $\rho$ be the largest modulus $|r_i|$ of a root with $P_i\ne0$, and $d$ the largest degree of such a $P_i$ with $|r_i|=\rho$. Then there are constants $0<c\le C$ such that
\[
|y_n|\le C\,n^d\rho^n\quad\text{for all }n\ge1,\qquad|y_n|\ge c\,n^d\rho^n\quad\text{for infinitely many }n .
\]
Consequently $\limsup_n|y_n|^{1/n}=\rho$, and the generating function of $(y_n)$ has the radius of convergence $1/\rho$.
\end{proposition}

\begin{proof}
For the upper bound, there are constants $C_i$ with $|P_i(n)|\le C_in^{\deg P_i}$ for all $n\ge1$, and each term $|P_i(n)r_i^n|/(n^d\rho^n)$ is bounded: it is at most $C_i$ if $|r_i|=\rho$, since then $\deg P_i\le d$, and it tends to zero if $|r_i|<\rho$, since a power of $|r_i|/\rho<1$ beats every polynomial (\vcref{1}{thm:standard-limits}). For the lower bound, let $a_i$ be the coefficient of $n^d$ in $P_i$ for the roots with $|r_i|=\rho$, and $\omega_i\coloneqq r_i/\rho$. Then
\[
\frac{y_n}{n^d\rho^n}=\sum_{|r_i|=\rho}a_i\,\omega_i^n+\varepsilon_n
\]
with $\varepsilon_n\to0$, by the same comparison of the remaining terms, of lower degree or smaller modulus. At least one $a_i$ is nonzero, so by \cref{lem:unimodular-means} the sum does not tend to zero: there is $c>0$ with $|\sum a_i\omega_i^n|\ge2c$ for infinitely many $n$, and for those that are large enough $|y_n|\ge cn^d\rho^n$. Finally $(cn^d\rho^n)^{1/n}$ and $(Cn^d\rho^n)^{1/n}$ tend to $\rho$, since $n^{1/n}\to1$, which gives the limit superior, and \vcref{1}{thm:cauchy-hadamard} gives the radius.
\end{proof}

The radius of convergence is thus the distance from $0$ to the nearest pole $1/r_i$ of the rational function $P/\tilde p$ that has not been cancelled by the numerator, the discrete counterpart of the statement of \cref{thm:rational-coefficients} that the radius of a power series solution is the distance to the nearest singular point. The proposition also makes precise how fast solutions grow: like $n^d\rho^n$, but only along a subsequence, since roots of equal modulus can cancel for some $n$. The solution $(2/\sqrt3)\sin(n\pi/3)$ of \vcref{1}{sec:recurrences} vanishes for every third $n$, and the lemma shows that this is the most that cancellation can achieve: the sum does not tend to zero.

\subsection{Sequences Without a Linear Recurrence}

\Cref{thm:rational-gf} also tells which sequences satisfy no linear recurrence with constant coefficients: those whose generating function is not rational. The Catalan numbers are an example. Their generating function $C(x)=\sum_nC_nx^n$ satisfies $C=1+xC^2$ \veqref{1}{eq:catalan-gf}, and \vcref{1}{sec:gf} solved this quadratic equation with the binomial series: $1-2xC(x)=N_{1/2}(-4x)$, the series of $\sqrt{1-4x}$.

\begin{proposition}[Catalan numbers]
\label{prop:catalan-no-recurrence}
The Catalan numbers satisfy no linear recurrence with constant coefficients. They satisfy the first-order recurrence with variable coefficients
\[
C_{n+1}=\frac{2(2n+1)}{n+2}\,C_n\qquad(n\in\mathbb N),
\]
and $C_n\sim4^n/(\sqrt\pi\,n^{3/2})$ as $n\to\infty$.
\end{proposition}

\begin{proof}
Suppose that $(C_n)$ satisfies such a recurrence, of some order $k$. Its generating function is $P/\tilde p$ by \cref{thm:rational-gf}, so $N_{1/2}(-4x)=1-2xP/\tilde p=U/V$ with polynomials $U\coloneqq\tilde p-2xP$ and $V\coloneqq\tilde p\ne0$. By \vcref{1}{thm:newton}, $N_{1/2}(-4x)^2=1-4x$, so $U^2=(1-4x)V^2$ in $\mathbb C[[x]]$, and therefore as polynomials. The degree of $U^2$ is even and that of $(1-4x)V^2$ is odd, a contradiction. The recurrence follows from $C_n=(2n)!/(n!\,(n+1)!)$ (\vcref{1}{thm:catalan}):
\[
\frac{C_{n+1}}{C_n}=\frac{(2n+2)(2n+1)}{(n+1)(n+2)}=\frac{2(2n+1)}{n+2} .
\]
Finally $C_n=\binom{2n}n/(n+1)$ and $\binom{2n}n\sim4^n/\sqrt{\pi n}$ (\vcref{1}{cor:central-binomial}), so $C_n\sim4^n/(\sqrt{\pi n}\,(n+1))\sim4^n/(\sqrt\pi\,n^{3/2})$.
\end{proof}

The growth $4^n/n^{3/2}$, with a power of $n$ that is not an integer, is itself incompatible with \cref{prop:recurrence-growth}, where the power of $n$ is always an integer. Applied to the first-order recurrence, \cref{thm:first-order-variable} turns the product of the ratios back into the closed form: $C_n=\prod_{j<n}2(2j+1)/(j+2)=2^n\cdot1\cdot3\cdots(2n-1)/(n+1)!$. Many sequences of combinatorics and of special functions satisfy linear recurrences whose coefficients are polynomials in $n$, as the Catalan numbers do, rather than constants; their generating functions satisfy linear differential equations, and they include the coefficients of the series solutions of \cref{ch:series-solutions}.

\subsection{Exponential Generating Functions}

A second way to pack a sequence into a series divides by factorials. The \term{exponential generating function} of $(y_n)$ is $\hat Y(x)\coloneqq\sum_ny_nx^n/n!$. For it, the shift corresponds to differentiation: by \vcref{1}{thm:termwise-diff}, the derivative of $\hat Y$ is $\sum_ny_{n+1}x^n/n!$, the exponential generating function of $Ey$. A recurrence for the sequence therefore becomes a differential equation for $\hat Y$ with the same characteristic polynomial.

\begin{proposition}[Exponential generating functions]
\label{prop:egf-recurrence}
Let $(y_n)$ solve $p(E)y=0$. Then its exponential generating function $\hat Y$ converges on all of $\mathbb R$ and solves the differential equation $p(D)\hat Y=0$ with the initial values $\hat Y^{(j)}(0)=y_j$, $j<k$. Conversely, for every solution $Y$ of $p(D)Y=0$ on an interval containing $0$, the numbers $y_n\coloneqq Y^{(n)}(0)$ solve $p(E)y=0$.
\end{proposition}

\begin{proof}
Let $M\coloneqq\max(1,|c_0|+\cdots+|c_{k-1}|)$ and $B\coloneqq\max_{j<k}|y_j|$. By induction, $|y_n|\le BM^n$ for all $n$: this holds for $n<k$, and if it holds below $n+k$, then $|y_{n+k}|\le\sum_j|c_j|BM^{n+j}\le BM^{n+k-1}\sum_j|c_j|\le BM^{n+k}$. So $\sum_n|y_n||x|^n/n!\le Be^{M|x|}$ converges for every $x$. By \vcref{1}{thm:termwise-diff}, $\hat Y^{(j)}(x)=\sum_ny_{n+j}x^n/n!$, and so $p(D)\hat Y(x)=\sum_n(p(E)y)_nx^n/n!=0$; at $x=0$, $\hat Y^{(j)}(0)=y_j$. Conversely, a solution of $p(D)Y=0$ is infinitely differentiable (\cref{prop:linear-structure}), and differentiating the equation $n$ times and setting $t=0$ gives $\sum_jc_jY^{(n+j)}(0)+Y^{(n+k)}(0)=0$.
\end{proof}

The proposition is a bridge between \cref{ch:linear-higher-order} and this chapter. By \cref{thm:constant-coefficients}, the solutions of $p(D)Y=0$ are combinations of $t^je^{r_it}$, and the $n$th derivative of $t^je^{rt}$ at $0$ is $n^{\underline j}r^{n-j}$, by the Leibniz rule; so the proposition yields the solutions $n^{\underline j}r^{n-j}$ of $p(E)y=0$, whose span for $j<m$ is that of $n^jr^n$, $j<m$, a third route to \cref{thm:constant-coefficient-recurrences} for $c_0\ne0$. For the Fibonacci numbers, $\hat F$ solves $Y''=Y'+Y$ with $Y(0)=0$ and $Y'(0)=1$, so $\hat F(x)=(e^{\phi x}-e^{\hat\phi x})/\sqrt5$, and Binet's formula appears once more.

The \term{$z$-transform} $\sum_ny_nz^{-n}$, the generating function evaluated at $x=1/z$, is the form in which generating functions are used in signal processing and control theory. For a solution of $p(E)y=0$ it is $z^kP(1/z)/p(z)$ by \cref{thm:rational-gf}, with poles at the characteristic roots themselves, and all solutions of a linear recurrence with constant coefficients tend to zero exactly when these poles lie inside the unit circle, as the next section shows. It plays for recurrences the role that the Laplace transform plays for differential equations; we do not develop it, and refer to Jury's book~\cite{jury1964} and to Elaydi~\cite[Chap.~6]{elaydi2005}.

\section{Systems and Stability}
\label{sec:difference-systems}

\subsection{Linear Systems of Difference Equations}

A system of linear difference equations with constant coefficients is a recurrence $\mathbf x_{n+1}=A\mathbf x_n$ for vectors $\mathbf x_n\in\mathbb C^d$ with a matrix $A\in\mathbb C^{d\times d}$, and its solution is $\mathbf x_n=A^n\mathbf x_0$. A scalar recurrence of order $k$ is the special case of the companion matrix: as in the proof of \cref{thm:casorati}, the vectors $\mathbf s_n=(y_n,\ldots,y_{n+k-1})$ satisfy $\mathbf s_{n+1}=C_p^{\mathsf T}\mathbf s_n$, with the transposed companion matrix of the characteristic polynomial (\cref{prop:companion}). Conversely, the behaviour of the powers $A^n$ was the subject of \cref{sec:matrix-powers,sec:powers-again}: for a diagonalisable $A=S\Lambda S^{-1}$ with the eigenvalues $\mu_i$ and the eigenvectors $\mathbf v_i$, the solution is $\mathbf x_n=\sum_ic_i\mu_i^n\mathbf v_i$, where $\mathbf x_0=\sum_ic_i\mathbf v_i$, and in general every entry of $A^n$ is a combination of the sequences $n^j\mu_i^n$ (\cref{thm:power-entries}). This is the counterpart of the solution $e^{tA}\mathbf y_0$ of $\mathbf y'=A\mathbf y$ by the eigenvalue method of \cref{sec:eigenvalue-method}, and the question of stability is answered in the same way.

\begin{definition}[Stability of a linear difference system]
\label{def:difference-stability}
Let $A\in\mathbb C^{d\times d}$. The equilibrium $\mathbf 0$ of $\mathbf x_{n+1}=A\mathbf x_n$ is \term{stable} if for every $\varepsilon>0$ there is $\delta>0$ such that $\|\mathbf x_0\|<\delta$ implies $\|\mathbf x_n\|<\varepsilon$ for all $n\in\mathbb N$, \term{asymptotically stable} if it is stable and every solution tends to $\mathbf 0$, and \term{unstable} otherwise.
\end{definition}

The norm may be any of the norms $\|\cdot\|_1$, $\|\cdot\|_2$, $\|\cdot\|_\infty$; since each of them is bounded by a constant multiple of each other one (\cref{prop:standard-norms}, which holds verbatim for complex vectors), the definition does not depend on the choice, and we mostly use $\|\cdot\|_\infty$ with the induced matrix norm. As for differential equations (\cref{prop:stability-exponential}), and with the same proof, with $A^n$ in the place of $e^{tA}$, the equilibrium $\mathbf 0$ is stable if and only if the powers $A^n$ are bounded, and asymptotically stable if and only if $A^n\to0$.

\begin{theorem}[Stability of linear difference systems]
\label{thm:difference-stability}
Let $A\in\mathbb C^{d\times d}$, with the spectral radius $\rho(A)$.
\begin{enumerate}
\item $A^n\to0$, and $\mathbf 0$ is asymptotically stable, if and only if $\rho(A)<1$, that is, every eigenvalue has modulus less than $1$.
\item $(A^n)$ is bounded, and $\mathbf 0$ is stable, if and only if $\rho(A)\le1$ and every eigenvalue of modulus $1$ is semisimple.
\item If an eigenvalue has modulus greater than $1$, or an eigenvalue of modulus $1$ is not semisimple, then $\mathbf x_{n+1}=A\mathbf x_n$ has unbounded solutions.
\end{enumerate}
\end{theorem}

\begin{proof}
(i) and (ii) are \cref{thm:powers-size}\,(i) and (ii), together with the equivalences just stated; an eigenvalue is semisimple exactly when its index is $1$. (iii) If $A\mathbf v=\mu\mathbf v$ with $|\mu|>1$ and $\mathbf v\ne\mathbf 0$, the solution $A^n\mathbf v=\mu^n\mathbf v$ is unbounded. If $|\mu|=1$ and $\mu$ is not semisimple, there is a Jordan chain $\mathbf v_1,\mathbf v_2$ with $(A-\mu I)\mathbf v_1=\mathbf 0$ and $(A-\mu I)\mathbf v_2=\mathbf v_1$, $\mathbf v_1\ne\mathbf 0$. Then $A\mathbf v_2=\mu\mathbf v_2+\mathbf v_1$, and by induction $A^n\mathbf v_2=\mu^n\mathbf v_2+n\mu^{n-1}\mathbf v_1$, whose norm is at least $n\|\mathbf v_1\|-\|\mathbf v_2\|$.
\end{proof}

For a scalar recurrence the condition can be read off the characteristic polynomial, and it is the condition under which the numerical solution of a recurrence does not amplify perturbations of its initial values, the stability that \cref{ch:eigenvalues,ch:canonical-forms} referred forward to this chapter.

\begin{corollary}[Stability of a scalar recurrence]
\label{cor:scalar-recurrence-stability}
Let $p$ be the characteristic polynomial of $p(E)y=0$.
\begin{enumerate}
\item All solutions tend to zero if and only if all roots of $p$ lie in the open unit disc $|\lambda|<1$.
\item All solutions are bounded if and only if all roots lie in the closed unit disc and the roots of modulus $1$ are simple.
\end{enumerate}
\end{corollary}

\begin{proof}
Every solution is a combination of the basis of \cref{thm:constant-coefficient-recurrences}, and each basis sequence is itself a solution. So all solutions tend to zero, or are bounded, if and only if all basis sequences do. The sequences $\mathbf e_l$ vanish from some index on. A sequence $n^jr^n$ with $r\ne0$ tends to zero if and only if $|r|<1$, by \vcref{1}{thm:standard-limits}, and it is bounded if and only if $|r|<1$, or $|r|=1$ and $j=0$; for a root of modulus $1$ the exponents $j$ run up to the multiplicity minus one.
\end{proof}

The corollary is the theorem applied to the companion matrix, whose eigenvalues are the roots of $p$ and which is nonderogatory (\cref{prop:nonderogatory}), so that an eigenvalue is semisimple exactly when it is a simple root. \Cref{tab:continuous-discrete} sets the discrete theory beside the continuous one of \cref{ch:linear-ode-systems}; the remaining entries are established in this section.

\begin{table}
\sdfloat{%
\begin{sdtable}{colspec={lll}, row{1}={font=\small\bfseries}}
\toprule
 & $\mathbf y'=A\mathbf y$ & $\mathbf x_{n+1}=A\mathbf x_n$ \\
\midrule
solution & $e^{tA}\mathbf y_0$ & $A^n\mathbf x_0$ \\
modes & $t^je^{\lambda t}\mathbf v$ & $n^j\mu^n\mathbf v$ \\
asymptotically stable & $\operatorname{Re}\lambda<0$ for all $\lambda$ & $|\mu|<1$ for all $\mu$ \\
rate of decay & spectral abscissa $\alpha(A)$ & spectral radius $\rho(A)$ \\
stable & also $\operatorname{Re}\lambda=0$, semisimple & also $|\mu|=1$, semisimple \\
test on coefficients & Routh--Hurwitz & Schur--Cohn, Jury \\
Lyapunov equation & $A^{\mathsf H}P+PA=-Q$ & $A^{\mathsf H}PA-P=-Q$ \\
variation of constants & $\int_{t_0}^te^{(t-s)A}\mathbf b(s)\,ds$ & $\sum_{m=0}^{n-1}A^{n-1-m}\mathbf b_m$ \\
determinant of solutions & $\det Y'=\operatorname{tr}A\,\det Y$ & $\det X_{n+1}=\det A\,\det X_n$ \\
\bottomrule
\end{sdtable}}{Linear systems of differential and difference equations with a constant matrix $A$ (\cref{thm:linear-stability,thm:lyapunov-equation,thm:duhamel,thm:difference-stability,thm:discrete-lyapunov,thm:discrete-duhamel}). The two columns are linked by $\mu=e^{h\lambda}$, which maps the left half-plane onto the unit disc without $0$ (\cref{prop:sampling-stability}).\label{tab:continuous-discrete}}
\end{table}

The last row refers to a matrix $X_n$ whose columns are solutions, so that $X_{n+1}=AX_n$, the counterpart of the fundamental matrices of \cref{sec:fundamental-matrices}; for the companion matrix it is the Casorati matrix of \cref{thm:casorati}.

When $\rho(A)<1$, the powers decay at least geometrically, at any rate exceeding $\rho(A)$; the constant in front measures the transient growth that \cref{fig:jordan-powers} showed for Jordan blocks.

\begin{lemma}[Geometric bound for powers]
\label{lem:power-bound}
Let $A\in\mathbb C^{d\times d}$ and $\sigma>\rho(A)$. There is a constant $C$ with $\|A^n\|_\infty\le C\sigma^n$ for all $n\in\mathbb N$.
\end{lemma}

\begin{proof}
By \cref{prop:norm-near-radius} with $\varepsilon=\sigma-\rho(A)$, there is an invertible $M$ with $\|M^{-1}AM\|_\infty\le\sigma$. Then $A^n=M(M^{-1}AM)^nM^{-1}$, and $\|A^n\|_\infty\le\|M\|_\infty\sigma^n\|M^{-1}\|_\infty$ by the submultiplicativity of the operator norm (\cref{prop:operator-norms}).
\end{proof}

\subsection{Lyapunov's Equation for Difference Systems}

Lyapunov's method of \cref{thm:lyapunov-equation} carries over to difference systems, and here its proof is even simpler. For $V(\mathbf x)=\mathbf x^{\mathsf H}P\mathbf x$ with a Hermitian positive definite $P$, along a solution of $\mathbf x_{n+1}=A\mathbf x_n$,
\[
V(\mathbf x_{n+1})-V(\mathbf x_n)=\mathbf x_n^{\mathsf H}\bigl(A^{\mathsf H}PA-P\bigr)\mathbf x_n ,
\]
so $V$ decreases along every solution if $A^{\mathsf H}PA-P$ is negative definite.

\begin{theorem}[Lyapunov's theorem for difference systems]
\label{thm:discrete-lyapunov}
Let $A\in\mathbb C^{d\times d}$.
\begin{enumerate}
\item If $\rho(A)<1$, then for every Hermitian $Q$ the \term{discrete Lyapunov equation}
\begin{equation}
A^{\mathsf H}PA-P=-Q
\label{eq:discrete-lyapunov}
\end{equation}
has exactly one solution, $P=\sum_{n=0}^\infty(A^{\mathsf H})^nQA^n$, the series converging entrywise. It is Hermitian, and positive definite if $Q$ is.
\item Conversely, if \eqref{eq:discrete-lyapunov} holds with Hermitian positive definite $P$ and $Q$, then $\rho(A)<1$, and every solution of $\mathbf x_{n+1}=A\mathbf x_n$ satisfies
\[
\|\mathbf x_n\|\le\sqrt{\kappa_2(P)}\,(1-\gamma)^{n/2}\,\|\mathbf x_0\|\qquad(n\in\mathbb N),\qquad\gamma\coloneqq\frac{\lambda_1(Q)}{\lambda_d(P)}\le1 .
\]
\end{enumerate}
\end{theorem}

\begin{proof}
(i) Choose $\sigma$ with $\rho(A)<\sigma<1$. By \cref{lem:power-bound}, $\|A^n\|_\infty\le C\sigma^n$, and since $\|A^{\mathsf H}\|_\infty=\|A\|_1\le d\,\|A\|_\infty$ (\cref{thm:norm-formulas}), the entries of $(A^{\mathsf H})^nQA^n$ are bounded by a constant times $\sigma^{2n}$; so the series converges entrywise, and absolutely. For its partial sums $P_N\coloneqq\sum_{n<N}(A^{\mathsf H})^nQA^n$,
\[
A^{\mathsf H}P_NA-P_N=\sum_{n=1}^{N}(A^{\mathsf H})^nQA^n-\sum_{n=0}^{N-1}(A^{\mathsf H})^nQA^n=(A^{\mathsf H})^NQA^N-Q ,
\]
which tends to $-Q$ as $N\to\infty$; so $P=\lim P_N$ solves \eqref{eq:discrete-lyapunov}. If $P'$ is another solution, $X\coloneqq P-P'$ satisfies $A^{\mathsf H}XA=X$, hence $X=(A^{\mathsf H})^nXA^n\to0$, and $X=0$. Each term of the series is Hermitian, and so is $P$. If $Q$ is positive definite and $\mathbf x\ne\mathbf 0$, then $\mathbf x^{\mathsf H}P\mathbf x\ge\mathbf x^{\mathsf H}Q\mathbf x>0$, since the other terms $(A^n\mathbf x)^{\mathsf H}Q(A^n\mathbf x)$ are nonnegative.

(ii) Let $V_n\coloneqq\mathbf x_n^{\mathsf H}P\mathbf x_n$. By \cref{thm:rayleigh}, $\lambda_1(P)\|\mathbf x_n\|^2\le V_n\le\lambda_d(P)\|\mathbf x_n\|^2$, and by the computation before the theorem
\[
V_{n+1}-V_n=-\mathbf x_n^{\mathsf H}Q\mathbf x_n\le-\lambda_1(Q)\|\mathbf x_n\|^2\le-\gamma V_n .
\]
So $0\le V_{n+1}\le(1-\gamma)V_n$, which for $\mathbf x_n\ne\mathbf 0$ forces $\gamma\le1$, and $V_n\le(1-\gamma)^nV_0$. Hence $\lambda_1(P)\|\mathbf x_n\|^2\le(1-\gamma)^n\lambda_d(P)\|\mathbf x_0\|^2$, which is the estimate, since $\kappa_2(P)=\lambda_d(P)/\lambda_1(P)$ (\cref{prop:singular-eigen}). If $\gamma=1$, all solutions vanish from $n=1$ on; otherwise they tend to $\mathbf 0$ geometrically. In both cases $\rho(A)<1$ by \cref{thm:difference-stability}.
\end{proof}

Where the continuous theorem needed the Sylvester equation and an integral, the discrete one needs only a geometric series of matrices, the counterpart of the integral $\int_0^\infty e^{tA^{\mathsf H}}Qe^{tA}\,dt$. For
\[
A=\begin{bmatrix}1/2&1\\0&1/2\end{bmatrix},\qquad A^n=\begin{bmatrix}2^{-n}&n2^{1-n}\\0&2^{-n}\end{bmatrix},
\]
the Jordan block $J_2(1/2)$, and $Q=I$, the series can be summed with $\sum_nx^n=1/(1-x)$, $\sum_nnx^n=x/(1-x)^2$ and $\sum_nn^2x^n=x(1+x)/(1-x)^3$ at $x=1/4$, which give $4/3$, $4/9$ and $20/27$:
\[
P=\sum_{n\ge0}(A^{\mathsf T})^nA^n=\begin{bmatrix}4/3&8/9\\8/9&116/27\end{bmatrix},
\]
and one checks $A^{\mathsf T}PA-P=-I$ directly. The eigenvalues of $P$ are $1.0871$ and $4.5425$, so $\kappa_2(P)=4.178$ and $\gamma=1/4.5425=0.2201$, and part (ii) guarantees $\|\mathbf x_n\|\le2.044\cdot0.8831^n\|\mathbf x_0\|$. As in \cref{ch:linear-ode-systems}, the certificate is valid but not sharp: the powers decay like $n2^{-n}$.

\subsection{Sampling and Discretising Differential Equations}

The solutions of $\mathbf y'=A\mathbf y$, observed at the times $t_n=nh$, satisfy $\mathbf y(t_{n+1})=e^{hA}\mathbf y(t_n)$ by \cref{cor:exp-laws}: sampling a linear differential system gives a linear difference system with the matrix $e^{hA}$, whose eigenvalues are the numbers $\mu=e^{h\lambda}$ for the eigenvalues $\lambda$ of $A$ (\cref{prop:exp-basic}). Since $|e^{h\lambda}|=e^{h\operatorname{Re}\lambda}$, the map $\lambda\mapsto e^{h\lambda}$ takes the open left half-plane onto the punctured open unit disc, the imaginary axis onto the unit circle, and the right half-plane onto the exterior of the circle. Sampling therefore preserves stability.

\begin{proposition}[Sampling preserves stability]
\label{prop:sampling-stability}
Let $A\in\mathbb C^{d\times d}$ and $h>0$. Then $\rho(e^{hA})=e^{h\alpha(A)}$, so that $e^{hA}$ satisfies $\rho(e^{hA})<1$ if and only if $\alpha(A)<0$. Moreover the powers $e^{nhA}$, $n\in\mathbb N$, are bounded if and only if $e^{tA}$ is bounded for $t\ge0$. Thus the sampled system $\mathbf x_{n+1}=e^{hA}\mathbf x_n$ is stable or asymptotically stable exactly when $\mathbf y'=A\mathbf y$ is.
\end{proposition}

\begin{proof}
The eigenvalues of $e^{hA}$ are the numbers $e^{h\lambda}$, of modulus $e^{h\operatorname{Re}\lambda}$, and the largest of these is $e^{h\alpha(A)}$. If $e^{tA}$ is bounded for $t\ge0$, so are the powers $(e^{hA})^n=e^{nhA}$. Conversely, if $\|e^{nhA}\|\le M$ for all $n$, write $t\ge0$ as $t=nh+\tau$ with $0\le\tau<h$; then $e^{tA}=e^{\tau A}e^{nhA}$ and $\|e^{tA}\|\le\|e^{\tau A}\|M\le e^{h\|A\|}M$, by \cref{prop:exp-basic}\,(i). The last statement follows from \cref{thm:linear-stability,thm:difference-stability}.
\end{proof}

A numerical method does not compute $e^{hA}$; it replaces it by an approximation, and the approximation may lose the property. Euler's method replaces $e^{hA}$ by $I+hA$, the Taylor polynomial of degree $1$, with the eigenvalues $1+h\lambda$; for small $h\lambda$ these are close to $e^{h\lambda}$, but $1+h\lambda$ lies in the unit disc only if $h\lambda$ lies in the disc $|1+z|<1$ of \cref{fig:stability-regions}.

\begin{proposition}[Step size condition for Euler's method]
\label{prop:euler-step-condition}
Let $A\in\mathbb C^{d\times d}$ and $h>0$.
\begin{enumerate}
\item The iteration $\mathbf x_{n+1}=(I+hA)\mathbf x_n$ of Euler's method is asymptotically stable if and only if every eigenvalue $\lambda$ of $A$ satisfies $\operatorname{Re}\lambda<0$ and
\[
h<\frac{2|\operatorname{Re}\lambda|}{|\lambda|^2} .
\]
\item If $A$ has an eigenvalue $i\omega$ with $\omega\ne0$, Euler's iteration is unstable for every $h>0$.
\item If $\alpha(A)<0$, the backward Euler iteration $\mathbf x_{n+1}=(I-hA)^{-1}\mathbf x_n$ is asymptotically stable for every $h>0$.
\end{enumerate}
\end{proposition}

\begin{proof}
(i) By \cref{thm:spectral-mapping} and \cref{thm:difference-stability}\,(i), Euler's iteration is asymptotically stable if and only if $|1+h\lambda|<1$ for every eigenvalue. Now $|1+h\lambda|^2=1+2h\operatorname{Re}\lambda+h^2|\lambda|^2$, which is less than $1$ if and only if $h|\lambda|^2<-2\operatorname{Re}\lambda$; this forces $\operatorname{Re}\lambda<0$ and is then the stated bound. (ii) For $\lambda=i\omega$, $|1+ih\omega|^2=1+h^2\omega^2>1$, and \cref{thm:difference-stability}\,(iii) applies. (iii) is \cref{prop:euler-stability}\,(ii).
\end{proof}

For a real negative eigenvalue the condition is $h<2/|\lambda|$, the condition of \cref{sec:euler}, and for the stiff system of \cref{tab:stiff}, with the eigenvalues $-1$ and $-1000$, it is $h<0.002$. For the oscillator $y''+\omega^2y=0$, with the eigenvalues $\pm i\omega$, the growth factor $|1+ih\omega|^2=1+h^2\omega^2$ is the factor by which Euler's method multiplied the energy in \cref{prop:discrete-energy}. The proposition thus collects the stability results for Euler's method of \cref{ch:first-order-equations,ch:linear-higher-order,ch:linear-ode-systems} under one statement about the eigenvalues of a difference system. \Cref{sec:ode-to-difference} turns to methods whose difference equations have a higher order than the differential equation, where a new phenomenon appears.

\subsection{Criteria from the Coefficients}

Whether all roots of a polynomial lie in the unit disc can be decided from its coefficients, as the Routh--Hurwitz criterion decides whether they lie in the left half-plane. For degrees $2$ and $3$ the discrete criteria follow from the continuous ones by a M\"obius transformation that maps the disc onto the half-plane.

\begin{lemma}[From the disc to the half-plane]
\label{lem:cayley-transform}
For $\lambda\ne1$ let $w\coloneqq(\lambda+1)/(\lambda-1)$. Then
\[
\operatorname{Re}w=\frac{|\lambda|^2-1}{|\lambda-1|^2},
\]
so that $|\lambda|<1$ if and only if $\operatorname{Re}w<0$. If $p$ is a polynomial of degree $k$ with $p(1)\ne0$ and the roots $\lambda_1,\ldots,\lambda_k$ (repeated according to multiplicity), then $q(w)\coloneqq(w-1)^kp\bigl((w+1)/(w-1)\bigr)$ is a polynomial of degree $k$ with the leading coefficient $p(1)$ and the roots $w_j=(\lambda_j+1)/(\lambda_j-1)$. In particular, all roots of $p$ lie in the open unit disc if and only if all roots of $q$ lie in the open left half-plane.
\end{lemma}

\begin{proof}
$\operatorname{Re}w=\operatorname{Re}\bigl((\lambda+1)(\bar\lambda-1)\bigr)/|\lambda-1|^2$, and $(\lambda+1)(\bar\lambda-1)=|\lambda|^2-1+(\bar\lambda-\lambda)$, whose last term is purely imaginary. If $p(\lambda)=a\prod_j(\lambda-\lambda_j)$, then
\[
q(w)=a\prod_{j=1}^k\bigl((w+1)-\lambda_j(w-1)\bigr)=a\prod_{j=1}^k\bigl((1-\lambda_j)\,w+(1+\lambda_j)\bigr).
\]
No $\lambda_j$ is $1$, since $p(1)\ne0$, so each factor has degree $1$ and the root $(1+\lambda_j)/(\lambda_j-1)=w_j$, and the leading coefficient is $a\prod_j(1-\lambda_j)=p(1)$. The last statement follows from the first, applied to each $\lambda_j$.
\end{proof}

\begin{theorem}[Polynomials of degree $2$ and $3$ with roots in the unit disc]
\label{thm:jury-small}
Let the polynomials below have real coefficients.
\begin{enumerate}
\item All roots of $p(\lambda)=\lambda^2+a_1\lambda+a_0$ lie in the open unit disc if and only if
\[
|a_0|<1\qquad\text{and}\qquad|a_1|<1+a_0 .
\]
\item All roots of $p(\lambda)=\lambda^3+a_2\lambda^2+a_1\lambda+a_0$ lie in the open unit disc if and only if, with $b_2\coloneqq3+a_2-a_1-3a_0$ and $b_1\coloneqq3-a_2-a_1+3a_0$,
\[
p(1)>0,\qquad p(-1)<0,\qquad b_2>0,\qquad b_2b_1>-p(1)\,p(-1).
\]
\end{enumerate}
\end{theorem}

\begin{proof}
If all roots $\lambda_j$ of a real monic polynomial of degree $k$ lie in the open unit disc, then $p(1)=\prod_j(1-\lambda_j)>0$ and $(-1)^kp(-1)=\prod_j(1+\lambda_j)>0$: the real roots give positive factors, and a pair of conjugate roots gives $|1\mp\lambda|^2>0$. (i) The condition $|a_1|<1+a_0$ means $p(1)=1+a_1+a_0>0$ and $p(-1)=1-a_1+a_0>0$. By \cref{lem:cayley-transform}, if $p(1)\ne0$,
\[
q(w)=(w+1)^2+a_1(w+1)(w-1)+a_0(w-1)^2=p(1)\,w^2+2(1-a_0)\,w+p(-1).
\]
If the roots of $p$ lie in the disc, then $|a_0|=|\lambda_1\lambda_2|<1$ and $p(\pm1)>0$, as shown above. Conversely, if $|a_0|<1$ and $p(\pm1)>0$, then $q/p(1)$ is a monic quadratic with the positive coefficients $2(1-a_0)/p(1)$ and $p(-1)/p(1)$, so its roots have negative real parts by \cref{thm:hurwitz-small}\,(i), and the roots of $p$ lie in the disc by the lemma.

(ii) Expanding $(w+1)^3$, $(w+1)^2(w-1)$, $(w+1)(w-1)^2$ and $(w-1)^3$ gives
\[
q(w)=p(1)\,w^3+b_2\,w^2+b_1\,w-p(-1).
\]
If the roots of $p$ lie in the disc, then $p(1)>0$ and $-p(-1)>0$, the lemma applies, and \cref{thm:hurwitz-small}\,(ii), applied to $q/p(1)$, gives $b_2/p(1)>0$, $-p(-1)/p(1)>0$ and $(b_2/p(1))(b_1/p(1))>-p(-1)/p(1)$, which is the condition. Conversely, if the condition holds, then $p(1)>0$, the lemma applies, and the same theorem shows that the roots of $q$ lie in the left half-plane, so those of $p$ lie in the disc.
\end{proof}

For $p(\lambda)=\lambda^2-\lambda+c$ the conditions of (i) are $|c|<1$ and $1<1+c$, that is, $0<c<1$; at $c=1$ the roots $e^{\pm i\pi/3}$ of the periodic recurrence of \vcref{1}{sec:recurrences} sit on the circle. For degree $2$ the condition $|a_0|<1$ is needed in addition to $p(\pm1)>0$: the polynomial $\lambda^2+2$ has $p(\pm1)=3$ but the roots $\pm i\sqrt2$. In general, the following test reduces the degree by one at each step.

\begin{theorem}[Schur--Cohn test]
\label{thm:schur-cohn}
Let $p(\lambda)=a_k\lambda^k+\cdots+a_1\lambda+a_0$ with complex coefficients, $a_k\ne0$, $k\ge1$, and let $p^*(\lambda)\coloneqq\bar a_0\lambda^k+\bar a_1\lambda^{k-1}+\cdots+\bar a_k$ be its conjugate reversed polynomial. If $|a_0|\ge|a_k|$, then $p$ has a root outside the open unit disc. If $|a_0|<|a_k|$, then
\[
(Tp)(\lambda)\coloneqq\frac{\bar a_k\,p(\lambda)-a_0\,p^*(\lambda)}{\lambda}
\]
is a polynomial of degree $k-1$, and all roots of $p$ lie in the open unit disc if and only if all roots of $Tp$ do.
\end{theorem}

\begin{proof}[sketch]
The product of the roots of $p$ has the modulus $|a_0/a_k|$, which gives the first statement. If $|a_0|<|a_k|$, the constant term of $\bar a_kp-a_0p^*$ is $\bar a_ka_0-a_0\bar a_k=0$ and its leading coefficient is $|a_k|^2-|a_0|^2\ne0$, so $Tp$ is a polynomial of degree $k-1$. On the unit circle, $\bar\lambda=1/\lambda$ gives $p^*(\lambda)=\lambda^k\overline{p(\lambda)}$, so $|p^*(\lambda)|=|p(\lambda)|$ and $|a_0p^*(\lambda)|<|\bar a_kp(\lambda)|$ wherever $p(\lambda)\ne0$. By Rouch\'e's theorem of complex analysis (Volume~III), $\bar a_kp-a_0p^*=\lambda\,Tp$ then has as many roots in the open disc as $p$, if $p$ has no root on the circle; removing the root $0$ shows that $Tp$ has one root fewer there. If $p$ has a root $\zeta$ on the circle, then $p^*(\zeta)=0$ as well, and $\zeta$ is a root of $Tp$. Henrici~\cite{henrici1974} gives the complete proof and the algorithm.
\end{proof}

Applying the test $k$ times decides the question with a number of arithmetic operations proportional to $k^2$. Jury arranged the same computation in a table, and his criteria for degree $3$ are equivalent to those of \cref{thm:jury-small}\,(ii)~\cite{jury1964}. The test is the discrete counterpart of the Routh--Hurwitz criterion of \cref{thm:routh-hurwitz}.

\subsection{Inhomogeneous Systems}

\begin{theorem}[Variation of constants for difference systems]
\label{thm:discrete-duhamel}
Let $A\in\mathbb C^{d\times d}$ and let $(\mathbf b_n)$ be a sequence of vectors.
\begin{enumerate}
\item The solution of $\mathbf x_{n+1}=A\mathbf x_n+\mathbf b_n$ with the initial value $\mathbf x_0$ is
\begin{equation}
\mathbf x_n=A^n\mathbf x_0+\sum_{m=0}^{n-1}A^{n-1-m}\,\mathbf b_m .
\label{eq:discrete-duhamel}
\end{equation}
\item If $\rho(A)<\sigma<1$, there is a constant $C$ with $\|A^n\|_\infty\le C\sigma^n$ for all $n$, and if moreover $\|\mathbf b_n\|_\infty\le\beta$ for all $n$, then
\[
\|\mathbf x_n\|_\infty\le C\sigma^n\|\mathbf x_0\|_\infty+\frac{C\beta}{1-\sigma}\qquad(n\in\mathbb N).
\]
\item If $\rho(A)<1$ and $\mathbf b_n\to\mathbf b_\infty$, then $\mathbf x_n\to(I-A)^{-1}\mathbf b_\infty$.
\end{enumerate}
\end{theorem}

\begin{proof}
(i) is \cref{thm:first-order-variable} with matrices: for $n=0$ the formula gives $\mathbf x_0$, and multiplying it by $A$ and adding $\mathbf b_n$ gives it for $n+1$. (ii) The constant exists by \cref{lem:power-bound}. Then \eqref{eq:discrete-duhamel} gives $\|\mathbf x_n\|_\infty\le C\sigma^n\|\mathbf x_0\|_\infty+C\beta\sum_{j<n}\sigma^j$, and the geometric sum is less than $1/(1-\sigma)$. (iii) $I-A$ is invertible, since $1$ is not an eigenvalue, and $\mathbf w_n\coloneqq\mathbf x_n-(I-A)^{-1}\mathbf b_\infty$ satisfies $\mathbf w_{n+1}=A\mathbf w_n+(\mathbf b_n-\mathbf b_\infty)$. Given $\varepsilon>0$, choose $N$ with $\|\mathbf b_n-\mathbf b_\infty\|_\infty\le\varepsilon$ for $n\ge N$. Applying (ii) to the solution started at the index $N$ gives $\|\mathbf w_n\|_\infty\le C\sigma^{n-N}\|\mathbf w_N\|_\infty+C\varepsilon/(1-\sigma)$ for $n\ge N$, which is less than $2C\varepsilon/(1-\sigma)$ for large $n$.
\end{proof}

The formula has the reading of Duhamel's formula (\cref{thm:duhamel}): the input $\mathbf b_m$ enters at the step from $m$ to $m+1$ and is then carried forward by $n-1-m$ applications of the homogeneous system. Part (ii) says that a stable system responds to bounded inputs with bounded outputs, and part (iii) that the solutions approach the equilibrium $(I-A)^{-1}\mathbf b_\infty$ of $\mathbf x=A\mathbf x+\mathbf b_\infty$, as in \cref{thm:bounded-input}. The same formula describes the propagation of errors. If each step of a computation of $\mathbf x_{n+1}=A\mathbf x_n$ commits an error $\mathbf b_n$, the computed sequence differs from the exact one by $\sum_mA^{n-1-m}\mathbf b_m$, and part (ii) bounds the accumulated error by $C\beta/(1-\sigma)$ when the system is asymptotically stable; for an unstable system the early errors are amplified by the growing powers. Volume~VI analyses the accumulation of rounding errors in this way.

\subsection{Iterative Methods, Markov Chains and Populations}

A linear system $B\mathbf x=\mathbf c$ can be solved by rewriting it as a fixed point equation $\mathbf x=G\mathbf x+\mathbf f$ and iterating, as the note on iterative methods in \cref{sec:special-systems} described for the methods of Jacobi and Gauss and Seidel. With the diagonal part $D$ of $B$, assumed invertible, and the rest $R\coloneqq B-D$, the equation $B\mathbf x=\mathbf c$ is $\mathbf x=-D^{-1}R\mathbf x+D^{-1}\mathbf c$, and Jacobi's method is the iteration with $G=-D^{-1}R$.

\begin{proposition}[Convergence of stationary iterations]
\label{prop:stationary-iteration}
Let $G\in\mathbb C^{d\times d}$ with $1$ not an eigenvalue, $\mathbf f\in\mathbb C^d$, and $\mathbf x^*\coloneqq(I-G)^{-1}\mathbf f$, the solution of $\mathbf x=G\mathbf x+\mathbf f$. The iterates $\mathbf x_{n+1}=G\mathbf x_n+\mathbf f$ converge to $\mathbf x^*$ for every $\mathbf x_0$ if and only if $\rho(G)<1$. In that case, for every $\sigma$ with $\rho(G)<\sigma<1$ there is a constant $C$ with $\|\mathbf x_n-\mathbf x^*\|_\infty\le C\sigma^n\|\mathbf x_0-\mathbf x^*\|_\infty$.
\end{proposition}

\begin{proof}
The errors $\mathbf e_n\coloneqq\mathbf x_n-\mathbf x^*$ satisfy $\mathbf e_{n+1}=G\mathbf e_n$, so $\mathbf e_n=G^n\mathbf e_0$, and $\mathbf e_n\to\mathbf 0$ for every $\mathbf e_0$ if and only if $G^n\to0$, that is, $\rho(G)<1$ (\cref{thm:difference-stability}). The bound is \cref{lem:power-bound}.
\end{proof}

Started at $\mathbf x_0=\mathbf 0$, the iterates are the partial sums $\sum_{j<n}G^j\mathbf f$ of the Neumann series, and the proposition repeats \cref{cor:neumann-spectral}. The rate of convergence is the spectral radius, and it can be close to $1$. For the second difference matrix $T_d$, $D=2I$ and Jacobi's matrix is $G=I-T_d/2$, whose eigenvalues are $\cos(j\pi/(d+1))$ by \cref{thm:second-difference-eigen}; so $\rho(G)=\cos(\pi/(d+1))\approx1-\pi^2/(2(d+1)^2)$. Reducing the component of the error along the slowest eigenvector by the factor $10^{-6}$ takes $\ln(10^6)/(-\ln\cos(\pi/(d+1)))$ steps: $334$ for $d=10$, $28\,554$ for $d=100$ and about $2.8\times10^6$ for $d=1000$, roughly $2.8(d+1)^2$. This is why the simple iterations are rarely used alone, and why Volume~X turns to the methods mentioned in that note.

A Markov chain $\mathbf p_{n+1}=M\mathbf p_n$ (\cref{sec:markov}) is a difference system whose matrix has the eigenvalue $1$, so it is never asymptotically stable. If $M$ has positive entries, \cref{thm:perron-positive} shows that the eigenvalue $1$ is simple and the others have modulus at most $q<1$; then the chain is stable, and on the invariant subspace of the vectors with component sum $0$, which contains the differences $\mathbf p_n-\boldsymbol\pi$, it is asymptotically stable, so that $\mathbf p_n\to\boldsymbol\pi$. Volume~IV develops the probability theory of such chains.

Populations structured by age give another nonnegative system. Divide a female population into $k$ age classes of equal length and let $\mathbf x_n\in\mathbb R^k$ hold the numbers in each class at the time $n$. In one period, a member of class $i$ has on average $f_i\ge0$ daughters, who enter class $1$, and survives into class $i+1$ with probability $s_i>0$. Then $\mathbf x_{n+1}=L\mathbf x_n$ with the \term{Leslie matrix}~\cite{leslie1945}
\[
L=\begin{bmatrix}f_1&f_2&\cdots&f_{k-1}&f_k\\s_1&0&\cdots&0&0\\0&s_2&\cdots&0&0\\\vdots&&\ddots&&\vdots\\0&0&\cdots&s_{k-1}&0\end{bmatrix}.
\]

\begin{proposition}[Growth rate of a population]
\label{prop:leslie}
Let $f_1,\ldots,f_k\ge0$ with $f_k>0$, let $s_1,\ldots,s_{k-1}>0$, and let $\ell_1\coloneqq1$ and $\ell_i\coloneqq s_1\cdots s_{i-1}$ for $i\ge2$, the probability of surviving into class $i$. Then $0$ is not an eigenvalue of $L$, and $\lambda\ne0$ is an eigenvalue if and only if
\[
\sum_{i=1}^kf_i\,\ell_i\,\lambda^{-i}=1 .
\]
This equation has exactly one positive root $r^*$, and every eigenvalue satisfies $|\lambda|\le r^*$; so $r^*=\rho(L)$.
\end{proposition}

\begin{proof}
For a vector $\mathbf v$, $(L\mathbf v)_1=\sum_if_iv_i$ and $(L\mathbf v)_i=s_{i-1}v_{i-1}$ for $i\ge2$. If $L\mathbf v=\mathbf 0$, then $v_1=\cdots=v_{k-1}=0$, and the first row gives $f_kv_k=0$, so $\mathbf v=\mathbf 0$. If $L\mathbf v=\lambda\mathbf v$ with $\lambda\ne0$ and $\mathbf v\ne\mathbf 0$, then $v_i=s_{i-1}v_{i-1}/\lambda$, so $v_i=\ell_i\lambda^{1-i}v_1$ and $v_1\ne0$; the first row becomes $\sum_if_i\ell_i\lambda^{1-i}v_1=\lambda v_1$, which is the equation. Conversely, for a nonzero root $\lambda$ of the equation, the vector with $v_i=\ell_i\lambda^{1-i}$ is an eigenvector. On $(0,\infty)$ the function $\varphi(t)\coloneqq\sum_if_i\ell_it^{-i}$ is continuous and strictly decreasing, since the term $f_k\ell_kt^{-k}$ is and the others do not increase, and it tends to $\infty$ as $t\to0^+$ and to $0$ as $t\to\infty$; by the intermediate value theorem it takes the value $1$ exactly once. For an eigenvalue $\lambda$, $1=|\sum_if_i\ell_i\lambda^{-i}|\le\varphi(|\lambda|)$, so $|\lambda|\le r^*$, since $\varphi$ is decreasing.
\end{proof}

The number $r^*$ is the long-term growth factor of the population per period, and the eigenvector $(\ell_i(r^*)^{1-i})_i$ its stable age distribution. For three classes with the fecundities $(0,1.2,1.0)$ and the survival probabilities $s_1=0.6$, $s_2=0.5$, the equation is $0.72\lambda^{-2}+0.3\lambda^{-3}=1$, that is, $\lambda^3-0.72\lambda-0.3=0$. Its positive root is $r^*=1.00867$, and the other two eigenvalues are $-0.5043\pm0.2075i$, of modulus $0.5454$. The eigenvalues are distinct, so $L$ is diagonalisable, and by \cref{sec:matrix-powers} $L^n\mathbf x_0/(r^*)^n$ tends to a multiple of the eigenvector, normalised to the proportions $(0.5292,0.3148,0.1560)$. Starting from $100$ newborns, the computed proportions differ from these by at most $7\times10^{-6}$ after $20$ periods and $1.1\times10^{-11}$ after $40$, and from then on the population grows by $0.87$ per cent per period. The rate of approach is the ratio $0.5454/1.00867=0.54$ of the second largest modulus to $r^*$.

\section{Three-Term Recurrences}
\label{sec:three-term}

\subsection{Minimal and Dominant Solutions}

Many functions of analysis satisfy a recurrence of the second order with variable coefficients,
\begin{equation}
y_{n+1}+a_ny_n+b_ny_{n-1}=0\qquad(n\ge1),
\label{eq:three-term}
\end{equation}
a \term{three-term recurrence}. The Bessel functions satisfy it in the order, $J_{n+1}(x)-(2n/x)J_n(x)+J_{n-1}(x)=0$ for fixed $x>0$ (\cref{prop:bessel-recurrences}), so that $a_n=-2n/x$ and $b_n=1$; the Legendre polynomials satisfy it in the degree, $(n+1)P_{n+1}(x)=(2n+1)xP_n(x)-nP_{n-1}(x)$ (\cref{thm:legendre-recurrence}), with $a_n=-(2n+1)x/(n+1)$ and $b_n=n/(n+1)$; and so do the Chebyshev polynomials and the determinants of tridiagonal matrices. We assume throughout that $b_n\ne0$ for all $n$, so that the recurrence can be run in both directions (\cref{prop:recurrence-ivp}). Its solutions form a space of dimension $2$ (\cref{thm:recurrence-solution-space}), and for two solutions $f$ and $g$ the Casorati determinant $C_n\coloneqq f_ng_{n+1}-f_{n+1}g_n$ satisfies
\begin{equation}
C_n=b_n\,C_{n-1}\qquad(n\ge1),
\label{eq:three-term-casorati}
\end{equation}
by \cref{thm:casorati} or by substituting $f_{n+1}$ and $g_{n+1}$ from the recurrence. For the Legendre recurrence $C_n=C_0/(n+1)$, and for the Bessel recurrence $C_n$ is constant.

\Cref{sec:series-computation} found that the Legendre recurrence computes $P_n(0.9)$ to full accuracy, and asked whether the Bessel recurrence can be used in the same way to compute $J_n(x)$ from $J_0(x)$ and $J_1(x)$. The answer depends on how the solution to be computed compares with the other solutions.

\begin{definition}[Minimal and dominant solutions]
\label{def:minimal-solution}
A nonzero solution $f$ of \eqref{eq:three-term} is \term{minimal} if there is a solution $g$ with $g_n\ne0$ for all large $n$ and $f_n/g_n\to0$ as $n\to\infty$. A solution that is not a multiple of a minimal solution is called \term{dominant}.
\end{definition}

\begin{proposition}[Uniqueness of the minimal solution]
\label{prop:minimal-unique}
Let $f$ be a minimal solution of \eqref{eq:three-term}.
\begin{enumerate}
\item Every solution $y$ that is not a multiple of $f$ satisfies $y_n\ne0$ for all large $n$ and $f_n/y_n\to0$.
\item Every minimal solution is a multiple of $f$.
\end{enumerate}
\end{proposition}

\begin{proof}
(i) Let $g$ be as in the definition. It is not a multiple of $f$, since $f_n/g_n$ would then be a nonzero constant, so $f$ and $g$ form a basis, and $y=\alpha f+\beta g$ with $\beta\ne0$. Then $y_n/g_n=\alpha f_n/g_n+\beta\to\beta$, so $y_n\ne0$ for large $n$, and $f_n/y_n=(f_n/g_n)/(y_n/g_n)\to0$. (ii) If a minimal solution $f'$ were not a multiple of $f$, then (i) applied to $f$ would give $f_n/f'_n\to0$, and (i) applied to $f'$ would give $f'_n/f_n\to0$, with all terms defined for large $n$; but the product of the two quotients is $1$.
\end{proof}

A minimal solution, if it exists, is thus unique up to a factor and negligible compared with every other solution, and the other solutions are all comparable with each other. For a recurrence with constant coefficients whose characteristic roots have different moduli, $|r_1|>|r_2|$, the minimal solution is $r_2^n$ and the dominant solutions are those with a nonzero component along $r_1^n$. The Fibonacci recurrence has the minimal solution $\hat\phi^n$, and the Fibonacci numbers are dominant. When $r_1\ne r_2$ and $|r_1|=|r_2|$, there is no minimal solution, since by \cref{lem:unimodular-means} no quotient of two independent solutions tends to zero; a double root $r$, on the other hand, has the minimal solution $r^n$, negligible compared with $nr^n$.

\subsection{The Direction of a Recurrence}

Suppose a solution $y$ of \eqref{eq:three-term} is computed by running the recurrence forward from slightly perturbed initial values. Even if every later step were exact, the perturbation would propagate as a solution of the recurrence, and its fate depends on the kind of solution being computed.

\begin{theorem}[Forward recurrence for minimal and dominant solutions]
\label{thm:recurrence-direction}
Let $f$ be a minimal and $g$ a dominant solution of \eqref{eq:three-term}, and let $\tilde y$ be the solution with the initial values $\tilde y_0=y_0+\delta_0$ and $\tilde y_1=y_1+\delta_1$, where $y$ is a solution. Then $\tilde y-y=\alpha f+\beta g$ with
\[
\beta=\frac{f_0\delta_1-f_1\delta_0}{f_0g_1-f_1g_0} .
\]
\begin{enumerate}
\item If $y=f$ is the minimal solution and $\beta\ne0$, then $(\tilde y_n-f_n)/g_n\to\beta$ while $f_n/g_n\to0$: the error eventually exceeds every fixed multiple of $|f_n|$, and the relative error $|\tilde y_n-f_n|/|f_n|$ is unbounded.
\item If $y=\gamma f+\delta g$ is dominant, $\delta\ne0$, then $(\tilde y_n-y_n)/y_n\to\beta/\delta$: the relative error converges to the relative perturbation of the dominant component.
\end{enumerate}
\end{theorem}

\begin{proof}
The difference $e\coloneqq\tilde y-y$ is a solution with $e_0=\delta_0$ and $e_1=\delta_1$. Writing $e=\alpha f+\beta g$, the conditions $\alpha f_0+\beta g_0=\delta_0$ and $\alpha f_1+\beta g_1=\delta_1$ form a linear system whose determinant $f_0g_1-f_1g_0$ is the Casorati determinant of $f$ and $g$, which is not zero because $f$ and $g$ are independent (\cref{thm:casorati}); Cramer's rule (\cref{thm:cramer}) gives $\beta$. (i) $e_n/g_n=\alpha f_n/g_n+\beta\to\beta$ and $f_n/g_n\to0$. So $|e_n|\ge|\beta||g_n|/2$ for large $n$, while $|f_n|\le\varepsilon|g_n|$ for every $\varepsilon>0$ and large $n$; hence $|e_n|/|f_n|\ge|\beta|/(2\varepsilon)$ whenever $f_n\ne0$. (ii) $e_n/y_n=(\alpha f_n/g_n+\beta)/(\gamma f_n/g_n+\delta)\to\beta/\delta$.
\end{proof}

In floating-point arithmetic the initial values carry rounding errors of relative size about $u$, so $\beta$ is typically of the order $u|f_0f_1|/|C_0|$, where $C_0=f_0g_1-f_1g_0$, and it is not zero. Every later step commits a new rounding error, which acts as a perturbation of the initial values of the remaining computation, and by superposition the total error is a sum of solutions started at the different steps, as in \cref{thm:discrete-duhamel}. The conclusion is the same: computed forward, a dominant solution is obtained with a relative error that does not grow beyond the accumulated relative perturbations, while a minimal solution is swamped by a multiple of the dominant solutions, however small its initial coefficient. Reversing the direction reverses the roles. Running \eqref{eq:three-term} backwards, $y_{n-1}=-(y_{n+1}+a_ny_n)/b_n$, the minimal solution becomes the one that grows relative to the others, and the dominant solutions decay relative to it. The following idea of J.~C.~P. Miller exploits this.

\subsection{Miller's Backward Recurrence}

To compute a minimal solution, start far out, at an index $N+1$, with arbitrary values that are not those of the solution, say $0$ and $1$, run the recurrence backwards to $n=0$, and fix the unknown factor at the end by a normalising condition. The errors of the arbitrary start are components along dominant solutions, which decay relative to the minimal one on the way down.

\begin{theorem}[Convergence of Miller's algorithm]
\label{thm:miller}
Let \eqref{eq:three-term} have a minimal solution $f$ with $f_0\ne0$, and for $N\ge1$ let $y^{(N)}$ be the solution with $y^{(N)}_{N+1}=0$ and $y^{(N)}_N=1$. Then for every fixed $n$,
\[
\frac{y^{(N)}_n}{y^{(N)}_0}\longrightarrow\frac{f_n}{f_0}\qquad(N\to\infty).
\]
\end{theorem}

\begin{proof}
Let $g$ be a solution that is not a multiple of $f$. By \cref{prop:minimal-unique}, $g_n\ne0$ for large $n$ and $\varepsilon_N\coloneqq f_{N+1}/g_{N+1}\to0$. For large $N$, the solution $z^{(N)}\coloneqq f-\varepsilon_Ng$ satisfies $z^{(N)}_{N+1}=0$, and it is not zero, since $f$ and $g$ are independent. The solutions with $y_{N+1}=0$ form a space of dimension $1$, since the initial-value map at the indices $N$ and $N+1$ is an isomorphism (\cref{prop:recurrence-ivp}); so $y^{(N)}$ is a multiple of $z^{(N)}$, and
\[
\frac{y^{(N)}_n}{y^{(N)}_0}=\frac{f_n-\varepsilon_Ng_n}{f_0-\varepsilon_Ng_0}\longrightarrow\frac{f_n}{f_0},
\]
where the denominator is nonzero for large $N$ because $f_0\ne0$.
\end{proof}

The proof shows the size of the error: the computed ratios differ from the true ones by terms of the order $\varepsilon_N=f_{N+1}/g_{N+1}$, the ratio of the minimal to a dominant solution at the starting index, which decides how large $N$ must be. Rounding errors committed during the backward recurrence are perturbations of the same kind and are damped in the same way. In practice the normalisation uses whatever relation is known: a value $f_0$ computed by another method, or an identity such as $\sum_n\lambda_nf_n=s$. For the Bessel functions, the generating function of the note in \cref{sec:bessel} at $t=1$ gives $1=\sum_{n\in\mathbb Z}J_n(x)$, and with $J_{-n}=(-1)^nJ_n$ (\cref{prop:bessel-negative-integer})
\begin{equation}
J_0(x)+2\sum_{k\ge1}J_{2k}(x)=1 ,
\label{eq:bessel-normalisation}
\end{equation}
which needs no other value of a Bessel function. \Cref{alg:miller} states the method with such a normalisation. Miller used it for the tables of Bessel functions of the British Association, published in 1952~\cite[Sect.~3.6]{olver2010nist}, and Gautschi~\cite{gautschi1967} analyses its convergence, including the normalisation by a sum.

\begin{algorithm}
\begin{algorithmic}[1]
\Require the coefficients $a_n$, $b_n\ne0$ of \eqref{eq:three-term}; the number $n_{\max}$ of wanted terms; a starting index $N>n_{\max}$; weights $\lambda_0,\ldots,\lambda_N$ and the value $s$ of $\sum_n\lambda_nf_n$
\Ensure approximations $\hat f_0,\ldots,\hat f_{n_{\max}}$ of the minimal solution $f$
\Name{MillerBackward}{a,b,n_{\max},N,\lambda,s}
\State $y_{N+1}\gets0$; $y_N\gets1$
\For{$n=N,N-1,\ldots,1$}
  \State $y_{n-1}\gets-(y_{n+1}+a_ny_n)/b_n$ \Comment{the recurrence run backwards}
\EndFor
\State $\sigma\gets\sum_{n=0}^N\lambda_ny_n$
\State \Return $\hat f_n\gets s\,y_n/\sigma$ for $n=0,\ldots,n_{\max}$
\end{algorithmic}
\caption{Miller's backward recurrence for the minimal solution of a three-term recurrence. The start $y_N=1$ may need to be replaced by a small number, or the $y_n$ rescaled during the loop, to avoid overflow. For the Bessel functions $J_n(x)$, $a_n=-2n/x$, $b_n=1$, $\lambda_0=1$, $\lambda_{2k}=2$, $\lambda_{2k+1}=0$ and $s=1$, by \eqref{eq:bessel-normalisation}.\label{alg:miller}}
\end{algorithm}

\subsection{Bessel Functions of Integral Order}

\begin{proposition}[The Bessel functions are minimal]
\label{prop:bessel-minimal}
Let $x>0$. As a sequence in $n$, $J_n(x)$ is a minimal solution of $y_{n+1}-(2n/x)y_n+y_{n-1}=0$. More precisely,
\[
\frac{n!\,J_n(x)}{(x/2)^n}\to1\qquad(n\to\infty),
\]
while the solution $g$ with $g_m=0$ and $g_{m+1}=1$, where $m$ is an integer with $m\ge x$, satisfies $g_{n+1}\ge(2n/x-1)\,g_n\ge g_n>0$ for $n>m$. Every solution that is not a multiple of $(J_n(x))$, in particular $(Y_n(x))$, is dominant.
\end{proposition}

\begin{proof}
By the series of \cref{thm:bessel-j}, $J_n(x)=(x/2)^nS_n/n!$ with
\[
S_n\coloneqq\sum_{k\ge0}\frac{(-1)^k(x^2/4)^k}{k!\,(n+1)(n+2)\cdots(n+k)},
\]
and
\[
|S_n-1|\le\sum_{k\ge1}\frac1{k!}\Bigl(\frac{x^2}{4(n+1)}\Bigr)^k=e^{x^2/(4(n+1))}-1 ,
\]
which tends to $0$. For $g$, induction on $n>m$: $g_{m+1}=1>0=g_m$, and if $g_n\ge g_{n-1}\ge0$ and $g_n>0$, then $g_{n+1}=(2n/x)g_n-g_{n-1}\ge(2n/x-1)g_n$, and $2n/x-1\ge1$ because $n>m\ge x$; so $g_{n+1}\ge g_n>0$. Consequently, for $n>m$ with $S_n\ne0$,
\[
\frac{J_{n+1}(x)/g_{n+1}}{J_n(x)/g_n}\le\frac{x}{2(n+1)}\,\frac{|S_{n+1}|}{|S_n|}\,\frac1{2n/x-1}\qquad\text{in absolute value},
\]
which tends to $0$; so $J_n(x)/g_n\to0$, and $J$ is minimal. The last statement is \cref{prop:minimal-unique}. The functions $Y_n$ satisfy the same recurrence and the relation $Y_0'=-Y_1$ (\cref{sec:bessel}), so the Casorati determinant of $(J_n(x))$ and $(Y_n(x))$, which is constant by \eqref{eq:three-term-casorati}, equals its value at $n=0$,
\[
J_0(x)Y_1(x)-J_1(x)Y_0(x)=-J_0(x)Y_0'(x)+J_0'(x)Y_0(x)=-W(J_0,Y_0)(x)=-\frac2{\pi x}\ne0
\]
by \cref{thm:bessel-y-limit}; hence $(Y_n(x))$ is not a multiple of $(J_n(x))$. The identity $J_{n+1}(x)Y_n(x)-J_n(x)Y_{n+1}(x)=2/(\pi x)$ is also \cite[Eq.~10.5.2]{olver2010nist}.
\end{proof}

The proof also shows where the character of the recurrence changes. For $n<x$ the coefficient $2n/x$ lies between $0$ and $2$, and the recurrence behaves locally like one with constant coefficients and characteristic roots $e^{\pm i\theta}$ on the unit circle, $\cos\theta=n/x$: all solutions oscillate with comparable amplitudes, as $J_n(x)$ and $Y_n(x)$ do. For $n>x$ the local roots are real, one larger and one smaller than $1$ in modulus, and the minimal solution separates from the others. \Cref{fig:bessel-minimal} shows $|J_n(x)|$ and $|Y_n(x)|$ as functions of $n$ for $x=1$ and $x=10$. The forward recurrence, started from the correctly rounded values of $J_0(x)$ and $J_1(x)$, follows $J_n(x)$ while $n$ is below $x$ and for some way beyond, and then turns up and grows like a dominant solution, roughly like $u|Y_n(x)|$: at $x=1$ this happens after $n=8$, at $x=10$ after $n=24$. So the forward recurrence answers the question of \cref{prop:bessel-recurrences} in the negative for $n>x$, while Miller's backward recurrence succeeds; \cref{sec:recurrence-computation} measures both.

\begin{figure}
\sdpanel{A}{\begin{tikzpicture}
\begin{axis}[sdplot, ymode=log, width=0.47\textwidth, height=0.42\textwidth, xmin=0, xmax=40, ymin=1e-62, ymax=1e60,
  xtick={0,10,20,30,40}, ytick={1e-60,1e-40,1e-20,1,1e20,1e40,1e60}, xlabel={$n$}, log basis y=10]
\addplot[black, line width=0.7pt] coordinates {(0,7.6520e-01) (1,4.4005e-01) (2,1.1490e-01) (3,1.9563e-02) (4,2.4766e-03) (5,2.4976e-04) (6,2.0938e-05) (7,1.5023e-06) (8,9.4223e-08) (9,5.2493e-09) (10,2.6306e-10) (11,1.1980e-11) (12,4.9997e-13) (13,1.9256e-14) (14,6.8854e-16) (15,2.2975e-17) (16,7.1864e-19) (17,2.1154e-20) (18,5.8803e-22) (19,1.5485e-23) (20,3.8735e-25) (21,9.2276e-27) (22,2.0982e-28) (23,4.5634e-30) (24,9.5111e-32) (25,1.9030e-33) (26,3.6608e-35) (27,6.7816e-37) (28,1.2114e-38) (29,2.0892e-40) (30,3.4829e-42) (31,5.6189e-44) (32,8.7817e-46) (33,1.3309e-47) (34,1.9576e-49) (35,2.7971e-51) (36,3.8855e-53) (37,5.2517e-55) (38,6.9112e-57) (39,8.8620e-59) (40,1.1079e-60)};
\addplot[black, dashed, line width=0.7pt] coordinates {(0,8.8257e-02) (1,7.8121e-01) (2,1.6507e+00) (3,5.8215e+00) (4,3.3278e+01) (5,2.6041e+02) (6,2.5708e+03) (7,3.0589e+04) (8,4.2567e+05) (9,6.7802e+06) (10,1.2162e+08) (11,2.4256e+09) (12,5.3241e+10) (13,1.2754e+12) (14,3.3106e+13) (15,9.2570e+14) (16,2.7738e+16) (17,8.8668e+17) (18,3.0120e+19) (19,1.0834e+21) (20,4.1140e+22) (21,1.6445e+24) (22,6.9028e+25) (23,3.0356e+27) (24,1.3957e+29) (25,6.6962e+30) (26,3.3467e+32) (27,1.7396e+34) (28,9.3906e+35) (29,5.2570e+37) (30,3.0481e+39) (31,1.8284e+41) (32,1.1333e+43) (33,7.2511e+44) (34,4.7846e+46) (35,3.2528e+48) (36,2.2765e+50) (37,1.6387e+52) (38,1.2124e+54) (39,9.2129e+55) (40,7.1849e+57)};
\addplot[only marks, mark=*, mark size=0.9pt, black] coordinates {(0,7.6520e-01) (1,4.4005e-01) (2,1.1490e-01) (3,1.9563e-02) (4,2.4766e-03) (5,2.4976e-04) (6,2.0938e-05) (7,1.5023e-06) (8,9.4198e-08) (9,4.8491e-09) (10,6.9148e-09) (11,1.4315e-07) (12,3.1423e-06) (13,7.5272e-05) (14,1.9539e-03) (15,5.4635e-02) (16,1.6371e+00) (17,5.2332e+01) (18,1.7776e+03) (19,6.3943e+04) (20,2.4281e+06) (21,9.7058e+07) (22,4.0740e+09) (23,1.7916e+11) (24,8.2373e+12) (25,3.9521e+14) (26,1.9752e+16) (27,1.0267e+18) (28,5.5423e+19) (29,3.1027e+21) (30,1.7990e+23) (31,1.0791e+25) (32,6.6886e+26) (33,4.2796e+28) (34,2.8239e+30) (35,1.9198e+32) (36,1.3436e+34) (37,9.6718e+35) (38,7.1558e+37) (39,5.4375e+39) (40,4.2405e+41)};
\end{axis}
\end{tikzpicture}}\hfill
\sdpanel{B}{\begin{tikzpicture}
\begin{axis}[sdplot, ymode=log, width=0.47\textwidth, height=0.42\textwidth, xmin=0, xmax=40, ymin=1e-22, ymax=1e20,
  xtick={0,10,20,30,40}, ytick={1e-20,1e-10,1,1e10,1e20}, xlabel={$n$}, log basis y=10]
\addplot[black, line width=0.7pt] coordinates {(0,2.4594e-01) (1,4.3473e-02) (2,2.5463e-01) (3,5.8379e-02) (4,2.1960e-01) (5,2.3406e-01) (6,1.4459e-02) (7,2.1671e-01) (8,3.1785e-01) (9,2.9186e-01) (10,2.0749e-01) (11,1.2312e-01) (12,6.3370e-02) (13,2.8972e-02) (14,1.1957e-02) (15,4.5080e-03) (16,1.5668e-03) (17,5.0565e-04) (18,1.5244e-04) (19,4.3146e-05) (20,1.1513e-05) (21,2.9072e-06) (22,6.9687e-07) (23,1.5902e-07) (24,3.4633e-08) (25,7.2146e-09) (26,1.4405e-09) (27,2.7620e-10) (28,5.0938e-11) (29,9.0498e-12) (30,1.5511e-12) (31,2.5681e-13) (32,4.1123e-14) (33,6.3759e-15) (34,9.5818e-16) (35,1.3971e-16) (36,1.9782e-17) (37,2.7225e-18) (38,3.6447e-19) (39,4.7501e-20) (40,6.0309e-21)};
\addplot[black, dashed, line width=0.7pt] coordinates {(0,5.5671e-02) (1,2.4902e-01) (2,5.8681e-03) (3,2.5136e-01) (4,1.4495e-01) (5,1.3540e-01) (6,2.8035e-01) (7,2.0102e-01) (8,1.0755e-03) (9,1.9930e-01) (10,3.5981e-01) (11,5.2033e-01) (12,7.8491e-01) (13,1.3635e+00) (14,2.7601e+00) (15,6.3647e+00) (16,1.6334e+01) (17,4.5905e+01) (18,1.3974e+02) (19,4.5716e+02) (20,1.5975e+03) (21,5.9328e+03) (22,2.3320e+04) (23,9.6676e+04) (24,4.2139e+05) (25,1.9260e+06) (26,9.2086e+06) (27,4.5959e+07) (28,2.3897e+08) (29,1.2923e+09) (30,7.2561e+09) (31,4.2245e+10) (32,2.5466e+11) (33,1.5876e+12) (34,1.0223e+13) (35,6.7931e+13) (36,4.6530e+14) (37,3.2822e+15) (38,2.3823e+16) (39,1.7777e+17) (40,1.3628e+18)};
\addplot[only marks, mark=*, mark size=0.9pt, black] coordinates {(0,2.4594e-01) (1,4.3473e-02) (2,2.5463e-01) (3,5.8379e-02) (4,2.1960e-01) (5,2.3406e-01) (6,1.4459e-02) (7,2.1671e-01) (8,3.1785e-01) (9,2.9186e-01) (10,2.0749e-01) (11,1.2312e-01) (12,6.3370e-02) (13,2.8972e-02) (14,1.1957e-02) (15,4.5080e-03) (16,1.5668e-03) (17,5.0565e-04) (18,1.5244e-04) (19,4.3146e-05) (20,1.1513e-05) (21,2.9072e-06) (22,6.9688e-07) (23,1.5906e-07) (24,3.4792e-08) (25,7.9443e-09) (26,4.9292e-09) (27,1.7688e-08) (28,9.0584e-08) (29,4.8958e-07) (30,2.7490e-06) (31,1.6004e-05) (32,9.6479e-05) (33,6.0146e-04) (34,3.8731e-03) (35,2.5736e-02) (36,1.7628e-01) (37,1.2435e+00) (38,9.0254e+00) (39,6.7349e+01) (40,5.1630e+02)};
\end{axis}
\end{tikzpicture}}
\caption{The minimal solution $|J_n(x)|$ (solid) and the dominant solution $|Y_n(x)|$ (dashed) of the Bessel recurrence, as functions of the order $n$, for (A) $x=1$ and (B) $x=10$; the values at integers are joined by lines. The dots are the moduli of the values computed in binary64 by the forward recurrence from the correctly rounded $J_0(x)$ and $J_1(x)$. They follow $J_n(x)$ up to $n=8$ in (A) and $n=24$ in (B) and then grow like a dominant solution (\cref{prop:bessel-minimal,thm:recurrence-direction}).\label{fig:bessel-minimal}}
\end{figure}

\subsection{Recurrences with Convergent Coefficients}

For many recurrences the coefficients converge as $n\to\infty$, and then the solutions behave asymptotically like those of the limiting recurrence with constant coefficients.

\begin{theorem}[Poincar\'e and Perron]
\label{thm:poincare-perron}
Suppose that $a_n\to a$ and $b_n\to b$ as $n\to\infty$, and that the roots $\lambda_1,\lambda_2$ of $\lambda^2+a\lambda+b$ satisfy $|\lambda_1|>|\lambda_2|$.
\begin{enumerate}
\item (Poincar\'e) For every solution of \eqref{eq:three-term} with $y_n\ne0$ for all large $n$, the ratio $y_{n+1}/y_n$ tends to $\lambda_1$ or to $\lambda_2$.
\item (Perron) If $b\ne0$, the solutions with $y_{n+1}/y_n\to\lambda_2$ are exactly the nonzero multiples of one solution $f$, which is minimal, and every other nonzero solution satisfies $y_{n+1}/y_n\to\lambda_1$.
\end{enumerate}
\end{theorem}

\begin{proof}[sketch]
Write the recurrence as a system for $(y_n,y_{n+1})$ whose matrix tends to the companion matrix of $\lambda^2+a\lambda+b$, and change coordinates so that the limit matrix becomes $\operatorname{diag}(\lambda_1,\lambda_2)$. In the new coordinates $(u_n,v_n)$ the system is diagonal up to perturbations that tend to zero. For Poincar\'e's statement one shows that the ratio $|v_n|/|u_n|$ either eventually stays below every $\varepsilon>0$ or eventually exceeds every bound, because the factor $|\lambda_2|/|\lambda_1|<1$ wins over the small perturbations; in the first case $y_{n+1}/y_n\to\lambda_1$, in the second $\to\lambda_2$. Perron's statement requires the construction of a solution in the second case, by a fixed point argument for the backward recurrence or by the convergence of an associated continued fraction. The ratio of a solution of the second kind to one of the first kind tends to zero geometrically, which makes it minimal. Elaydi~\cite[Chap.~8]{elaydi2005} gives complete proofs for equations of any order, and Gautschi~\cite{gautschi1967} discusses the second-order case and its use in computation; the first statement is from Poincar\'e's memoir of 1885~\cite{poincare1885}.
\end{proof}

The Legendre recurrence has $a_n\to-2x$ and $b_n\to1$, with the limiting polynomial $\lambda^2-2x\lambda+1$ and the roots $x\pm\sqrt{x^2-1}$. For real $x>1$ the roots are real with $x+\sqrt{x^2-1}>1>x-\sqrt{x^2-1}>0$, and the theorem applies.

\begin{proposition}[Legendre polynomials outside the interval]
\label{prop:legendre-outside}
For $x\ge1$ the sequence $P_n(x)$ is nondecreasing in $n$, with $P_n(x)\ge1$, and for $x>1$
\[
P_{n+1}(x)\ge\frac{(2n+1)x-n}{n+1}\,P_n(x),
\]
so that $P_n(x)$ grows at least geometrically. For $x>1$, $(P_n(x))$ is a dominant solution of the Legendre recurrence, and $P_{n+1}(x)/P_n(x)\to x+\sqrt{x^2-1}$.
\end{proposition}

\begin{proof}
$P_0=1$ and $P_1=x\ge1$. If $P_n\ge P_{n-1}\ge1$, the recurrence gives
\[
P_{n+1}=\frac{(2n+1)xP_n-nP_{n-1}}{n+1}\ge\frac{(2n+1)x-n}{n+1}\,P_n\ge\frac{n+1}{n+1}\,P_n=P_n ,
\]
using $P_{n-1}\le P_n$ and $x\ge1$. For $x>1$ the factor $((2n+1)x-n)/(n+1)$ tends to $2x-1>1$, so $P_n$ grows at least geometrically. By \cref{thm:poincare-perron}\,(i), $P_{n+1}/P_n$ tends to one of the two roots, and since it is at least $1>x-\sqrt{x^2-1}$, it tends to the larger. A minimal solution would have the ratio tending to the smaller root, by part (ii), so $P$ is not minimal, and it is dominant.
\end{proof}

So for $x>1$ the forward recurrence is the right way to compute $P_n(x)$, by \cref{thm:recurrence-direction}\,(ii), and the minimal solution, a Legendre function of the second kind~\cite[Chap.~XV]{whittaker1927}, must be computed backwards. For $-1<x<1$ the limiting roots $e^{\pm i\theta}$, $x=\cos\theta$, have the same modulus, and the theorem says nothing. Both $P_n(x)$ and the solution $Q_n(x)$ of \cref{sec:legendre}, which is defined on $(-1,1)$, oscillate in $n$ with amplitudes of the order $n^{-1/2}$, as the asymptotic formulas for large degree show~\cite[Chap.~VIII]{szego1975}, \cite[Sect.~14.15]{olver2010nist}; this agrees with their Casorati determinant $C_0/(n+1)$. Neither solution is negligible compared with the other, and an error committed in the forward recurrence is carried along with an amplitude of the same order as the solution. This is why the recurrence computed $P_n(0.9)$ to nearly full accuracy in \cref{tab:legendre-evaluation}. The same applies to the Chebyshev recurrence on $(-1,1)$, which has constant coefficients and the different roots $e^{\pm i\theta}$ of modulus $1$, and \cref{sec:recurrence-computation} confirms it.

\subsection{An Integral Computed in Two Directions}

The phenomenon occurs already in first-order recurrences. The integrals
\[
I_n\coloneqq\int_0^1x^ne^{x-1}\,dx\qquad(n\in\mathbb N)
\]
satisfy $I_0=1-e^{-1}$ and, by integration by parts with $u=x^n$ and $dv=e^{x-1}dx$,
\begin{equation}
I_n=1-n\,I_{n-1}\qquad(n\ge1).
\label{eq:integral-recurrence}
\end{equation}
Since $e^{-1}\le e^{x-1}\le1$ on $[0,1]$, the values satisfy $e^{-1}/(n+1)\le I_n\le1/(n+1)$: they are positive and decrease slowly to $0$. The homogeneous recurrence $z_n=-nz_{n-1}$ has the solutions $c(-1)^nn!$, which grow factorially, and $I_n$ is negligible compared with them, the analogue of a minimal solution.

\begin{proposition}[Two directions for $I_n$]
\label{prop:integral-two-directions}
\leavevmode
\begin{enumerate}
\item If \eqref{eq:integral-recurrence} is run forward in exact arithmetic from $\tilde I_0=I_0+\delta$, then $\tilde I_n-I_n=(-1)^nn!\,\delta$.
\item If it is run backward, $\tilde I_{n-1}=(1-\tilde I_n)/n$, in exact arithmetic from $\tilde I_N\coloneqq0$, then $|\tilde I_n-I_n|\le n!/(N+1)!$ for $0\le n\le N$.
\end{enumerate}
\end{proposition}

\begin{proof}
The differences $e_n\coloneqq\tilde I_n-I_n$ satisfy $e_n=-ne_{n-1}$ in both cases. Forward, $e_n=(-1)^nn!\,e_0$. Backward, $e_{n-1}=-e_n/n$, so $e_n=(-1)^{N-n}(n!/N!)\,e_N$, and $|e_N|=I_N\le1/(N+1)$.
\end{proof}

In binary64 the initial value $I_0$ carries a rounding error of about $10^{-17}$, which the forward recurrence multiplies by $n!$: since $18!=6.4\times10^{15}$, nothing of $I_{18}$ survives. The backward recurrence started at $N=30$ with the absurd value $0$ gives $I_{20}$ with an error below $20!/31!=3.0\times10^{-16}$. \Cref{sec:recurrence-computation} shows the computed values.

\subsection{Continued Fractions}

A minimal solution is connected with a continued fraction. If $f$ is a solution of \eqref{eq:three-term} without zeros, its ratios $r_n\coloneqq f_n/f_{n-1}$ satisfy $r_{n+1}+a_n+b_n/r_n=0$, that is, $r_n=-b_n/(a_n+r_{n+1})$, and repeating the substitution gives
\[
\frac{f_1}{f_0}=\cfrac{-b_1}{a_1-\cfrac{b_2}{a_2-\cfrac{b_3}{a_3-\cdots}}} .
\]
Stopping after the $N$th level, with $r_{N+1}$ replaced by $0$, gives the $N$th \term{approximant} of the continued fraction, and this is exactly the ratio $y^{(N)}_1/y^{(N)}_0$ of \cref{thm:miller}, since the backward recurrence started with $y_{N+1}=0$ computes the ratios $r_n=y_n/y_{n-1}$ from the same formula. So \cref{thm:miller} shows that the approximants converge to $f_1/f_0$ whenever a minimal solution exists, and the converse holds as well.

\begin{theorem}[Pincherle]
\label{thm:pincherle}
Let $b_n\ne0$ for all $n$. The recurrence \eqref{eq:three-term} has a minimal solution if and only if the continued fraction above converges, that is, its approximants converge, possibly to $\infty$. In that case its value is $f_1/f_0$ for the minimal solution $f$, read as $\infty$ if $f_0=0$.
\end{theorem}

\begin{proof}[sketch]
If a minimal solution exists, the approximants are the ratios $y^{(N)}_1/y^{(N)}_0$, and the proof of \cref{thm:miller} shows that they converge to $f_1/f_0$; the case $f_0=0$ is handled in the same way, with the reciprocals. Conversely, if the approximants converge, one defines the ratios of a candidate solution by the tails of the continued fraction and shows, through the determinant formula for consecutive approximants, which is the Casorati relation \eqref{eq:three-term-casorati}, that this solution is minimal. Gautschi~\cite{gautschi1967} gives the proof and the applications to special functions; Volume~XI uses continued fractions to compute them.
\end{proof}

For the Bessel recurrence, $a_n=-2n/x$ and $b_n=1$, and the theorem gives the continued fraction
\[
\frac{J_1(x)}{J_0(x)}=\cfrac1{\dfrac2x-\cfrac1{\dfrac4x-\cfrac1{\dfrac6x-\cdots}}} .
\]
At $x=1$ the approximant of the fourth level, $1/(2-1/(4-1/(6-1/8)))=0.5750799$, differs from $J_1(1)/J_0(1)=0.5750809$ by $1.0\times10^{-6}$.

\subsection{Orthogonal Polynomials and Jacobi Matrices}

The note on three-term recurrences in \cref{subsec:legendre} showed that the monic orthogonal polynomials of an inner product $\int_a^bwpq$ satisfy a recurrence
\begin{equation}
p_{n+1}(x)=(x-\alpha_n)\,p_n(x)-\beta_n\,p_{n-1}(x),\qquad p_0=1,\quad p_{-1}=0,
\label{eq:orthogonal-recurrence}
\end{equation}
and one checks from the same argument that $\beta_n=\langle p_n,p_n\rangle/\langle p_{n-1},p_{n-1}\rangle>0$. The converse is also true: every such recurrence with $\beta_n>0$ produces orthogonal polynomials.

\begin{theorem}[Favard]
\label{thm:favard}
Let $\alpha_0,\alpha_1,\ldots$ be real and $\beta_1,\beta_2,\ldots$ positive, and let $p_n$ be defined by \eqref{eq:orthogonal-recurrence}. Then there is exactly one linear functional $\mathcal L$ on the real polynomials with $\mathcal L(1)=1$ and $\mathcal L(p_mp_n)=0$ for $m\ne n$. It satisfies $\mathcal L(p_n^2)=\beta_1\cdots\beta_n$ and $\mathcal L(q^2)>0$ for every real polynomial $q\ne0$, so that $\langle p,q\rangle\coloneqq\mathcal L(pq)$ is an inner product on the real polynomials, and $p_0,p_1,\ldots$ are its monic orthogonal polynomials.
\end{theorem}

\begin{proof}
By induction, $p_n$ is monic of degree $n$, so $p_0,p_1,\ldots$ form a basis of the real polynomials, and a linear functional is determined by its values on them. If $\mathcal L$ has the stated properties, then $\mathcal L(p_0)=1$ and $\mathcal L(p_n)=\mathcal L(p_np_0)=0$ for $n\ge1$; this proves uniqueness, and we define $\mathcal L$ by these values. We show that $\mathcal L(x^kp_n)=0$ for $0\le k<n$, by induction on $k$. For $k=0$ it is the definition. For $1\le k<n$, the recurrence gives $xp_n=p_{n+1}+\alpha_np_n+\beta_np_{n-1}$, so
\[
\mathcal L(x^kp_n)=\mathcal L(x^{k-1}p_{n+1})+\alpha_n\mathcal L(x^{k-1}p_n)+\beta_n\mathcal L(x^{k-1}p_{n-1}),
\]
and all three terms vanish by the induction hypothesis, since $k-1<n-1<n<n+1$. Hence $\mathcal L(qp_n)=0$ for every polynomial $q$ of degree less than $n$, and in particular $\mathcal L(p_mp_n)=0$ for $m<n$. Since $p_n=x^n+(\text{lower terms})$, the same formula with $k=n$ gives
\[
\mathcal L(p_n^2)=\mathcal L(x^np_n)=\mathcal L(x^{n-1}p_{n+1})+\alpha_n\mathcal L(x^{n-1}p_n)+\beta_n\mathcal L(x^{n-1}p_{n-1})=\beta_n\mathcal L(p_{n-1}^2),
\]
so $\mathcal L(p_n^2)=\beta_1\cdots\beta_n>0$. Finally, a real polynomial $q\ne0$ is $\sum_nc_np_n$ with real $c_n$ not all zero, and $\mathcal L(q^2)=\sum_nc_n^2\mathcal L(p_n^2)>0$.
\end{proof}

The theorem gives the inner product as a functional, not yet as an integral. That $\mathcal L(q)=\int q\,d\mu$ for a positive measure $\mu$ on the real line is the deeper part of Favard's theorem of 1935~\cite{favard1935}; it requires the measure theory of Volume~III, and Chihara~\cite[Chaps.~I and~II]{chihara1978} gives the proof. The recurrence also turns the zeros of orthogonal polynomials into eigenvalues.

\begin{proposition}[Jacobi matrices]
\label{prop:jacobi-matrix}
Let $p_n$ be defined by \eqref{eq:orthogonal-recurrence} with real $\alpha_n$ and positive $\beta_n$, and let $J_n$ be the real symmetric tridiagonal $n\times n$ matrix with the diagonal $\alpha_0,\ldots,\alpha_{n-1}$ and the entries $\sqrt{\beta_1},\ldots,\sqrt{\beta_{n-1}}$ next to it, the \term{Jacobi matrix}. Then $p_n(x)=\det(xI-J_n)$, and the zeros of $p_n$ are real and simple: they are the $n$ different eigenvalues of $J_n$.
\end{proposition}

\begin{proof}
The matrix $xI-J_n$ is tridiagonal with the diagonal entries $x-\alpha_{k-1}$ and the entries $-\sqrt{\beta_k}$ above and below the diagonal, so by \cref{prop:tridiagonal-det} its leading principal minors satisfy $D_k=(x-\alpha_{k-1})D_{k-1}-\beta_{k-1}D_{k-2}$, with $D_0=1$ and $D_1=x-\alpha_0$. These are the recurrence and the initial values of the $p_k$, so $D_n=p_n$. The eigenvalues of the real symmetric matrix $J_n$ are real (\cref{prop:hermitian-eigen}). If $J_n\mathbf v=\mu\mathbf v$, the first row gives $v_2=(\mu-\alpha_0)v_1/\sqrt{\beta_1}$, and row $i$ gives $v_{i+1}$ in terms of $v_{i-1}$ and $v_i$, since $\sqrt{\beta_i}\ne0$; so $\mathbf v$ is determined by $v_1$, and every eigenspace has dimension $1$. A real symmetric matrix is diagonalisable (\cref{thm:spectral}), so every eigenvalue is simple, and the $n$ eigenvalues are different.
\end{proof}

For the Legendre polynomials, normalised to be monic, the recurrence of \cref{thm:legendre-recurrence} becomes \eqref{eq:orthogonal-recurrence} with $\alpha_n=0$ and $\beta_n=n^2/(4n^2-1)$, as one checks with the leading coefficients $(2n)!/(2^n(n!)^2)$ of \cref{thm:rodrigues}. For $n=2$, $J_2$ has the zero diagonal and the off-diagonal entries $\sqrt{1/3}$, and its eigenvalues $\pm1/\sqrt3$ are the zeros of $P_2$, the nodes of the two-point Gauss rule of \cref{subsec:legendre}. Computing the zeros of orthogonal polynomials as eigenvalues of Jacobi matrices, together with the weights of the Gauss rules from the eigenvectors, is the method of Golub and Welsch~\cite{golubwelsch1969}, and Volume~XII uses it; the eigenvalues of symmetric tridiagonal matrices are the subject of Volume~IX. \Cref{prop:eigen-count} counted them by the signs of the pivots, which are ratios of consecutive minors $D_k$, a Sturm sequence for the polynomials $p_k$.

\subsection{Summing Chebyshev Series by Clenshaw's Algorithm}

A recurrence can also be run backwards for a different purpose: to evaluate a sum of solutions without computing the solutions. Sums $\sum_{k=0}^Nc_kT_k(x)$ of Chebyshev polynomials, which Volume~XI uses to approximate functions, are evaluated in this way.

\begin{proposition}[Clenshaw's algorithm]
\label{prop:clenshaw}
Let $c_0,\ldots,c_N$ and $x$ be real, let $b_{N+1}=b_{N+2}=0$, and let
\[
b_k=c_k+2x\,b_{k+1}-b_{k+2}\qquad(k=N,N-1,\ldots,0).
\]
Then $\sum_{k=0}^Nc_kT_k(x)=b_0-x\,b_1$.
\end{proposition}

\begin{proof}
Substitute $c_k=b_k-2xb_{k+1}+b_{k+2}$ and collect the coefficients of each $b_j$:
\[
\sum_{k=0}^Nc_kT_k=\sum_{j=0}^{N}b_j\bigl(T_j-2xT_{j-1}+T_{j-2}\bigr),
\]
where $T_{-1}$ and $T_{-2}$ are read as $0$, and the terms with $b_{N+1}$ and $b_{N+2}$ vanish. For $j\ge2$ the bracket is $0$ by the recurrence \eqref{eq:chebyshev-recurrence}; for $j=1$ it is $T_1-2xT_0=-x$, and for $j=0$ it is $T_0=1$.
\end{proof}

\begin{algorithm}
\begin{algorithmic}[1]
\Require the coefficients $c_0,\ldots,c_N$ of a Chebyshev sum; the point $x$
\Ensure $s=\sum_{k=0}^Nc_kT_k(x)$
\Name{Clenshaw}{c,x}
\State $b_{N+1}\gets0$; $b_{N+2}\gets0$
\For{$k=N,N-1,\ldots,0$}
  \State $b_k\gets c_k+2x\,b_{k+1}-b_{k+2}$ \Comment{the recurrence of $T_k$, run backwards with input $c_k$}
\EndFor
\State \Return $s\gets b_0-x\,b_1$
\end{algorithmic}
\caption{Clenshaw's algorithm for a sum of Chebyshev polynomials (\cref{prop:clenshaw}). Each coefficient costs one multiplication by the precomputed $2x$ and two additions, and the values $T_k(x)$ are never formed.\label{alg:clenshaw}}
\end{algorithm}

The numbers $b_k$ satisfy the Chebyshev recurrence backwards, with the coefficients $c_k$ as inputs, and the algorithm is the counterpart of Horner's scheme for the basis of Chebyshev polynomials; Clenshaw introduced it in 1955~\cite{clenshaw1955}, and it works in the same way for every family defined by a three-term recurrence. For the sum with $c_k=1/(k+1)^2$ at the binary64 number nearest to $0.9$, the algorithm in binary64 gave the relative errors $6.5\times10^{-18}$, $8.1\times10^{-17}$ and $5.5\times10^{-17}$ for $N=10$, $40$ and $160$, measured against the exact rational value of the sum. Volume~XI analyses its accuracy and uses it.

\section{From Differential to Difference Equations}
\label{sec:ode-to-difference}

\subsection{Multistep Methods and Parasitic Solutions}

A numerical method for $y'=f(t,y)$ replaces the differential equation by a difference equation for the approximations $y_n$ of $y(t_n)$, $t_n=t_0+nh$. Euler's method and the methods of Heun and of the midpoint in \cref{sec:euler} compute $y_{n+1}$ from $y_n$ alone, and for the test equation $y'=\lambda y$ they give first-order recurrences $y_{n+1}=R(h\lambda)y_n$, which \cref{sec:difference-systems} analysed. Methods that use several previous values give recurrences of higher order, and these have more solutions than the differential equation. The simplest is the \term{explicit midpoint rule}, or \term{leapfrog method},
\begin{equation}
y_{n+1}=y_{n-1}+2h\,f(t_n,y_n)\qquad(n\ge1),
\label{eq:leapfrog}
\end{equation}
which replaces $y'(t_n)$ by the central difference quotient $D_h^0y(t_n)=(y(t_{n+1})-y(t_{n-1}))/(2h)$, whose error is of order $h^2$ (\vcref{1}{prop:difference-taylor}\,(ii)). It needs two starting values, $y_0$ and an approximation $y_1$ of $y(t_1)$ computed by another method. For the test equation it is the recurrence $y_{n+1}-2h\lambda\,y_n-y_{n-1}=0$ of order $2$, whose two characteristic roots must be compared with the single factor $e^{h\lambda}$ of the exact solution.

\begin{proposition}[Roots of the leapfrog method]
\label{prop:leapfrog-roots}
Let $\lambda<0$ and $h>0$, let $\sigma\coloneqq-h\lambda>0$ and $s\coloneqq\operatorname{arsinh}\sigma$. The characteristic roots of $y_{n+1}-2h\lambda y_n-y_{n-1}=0$ are
\[
\mu_1=e^{-s}\in(0,1)\qquad\text{and}\qquad\mu_2=-e^{s}<-1 ,
\]
and the solution with the starting values $y_0$ and $y_1$ is $y_n=A\mu_1^n+B\mu_2^n$ with
\[
B=\frac{y_1-\mu_1y_0}{\mu_2-\mu_1},\qquad A=y_0-B .
\]
If $y_0\ne0$ and $y_1=e^{h\lambda}y_0$ is the exact value, then $B\ne0$, and $|y_n|\to\infty$ as $n\to\infty$.
\end{proposition}

\begin{proof}
The characteristic polynomial is $\mu^2+2\sigma\mu-1$, with the roots $-\sigma\pm\sqrt{\sigma^2+1}$. Since $\sinh s=\sigma$ and $\cosh s=\sqrt{1+\sinh^2s}=\sqrt{1+\sigma^2}$ (\vcref{1}{sec:hyperbolic}), the roots are $\cosh s-\sinh s=e^{-s}$ and $-(\cosh s+\sinh s)=-e^s$, and $s>0$. They are different, so $y_n=A\mu_1^n+B\mu_2^n$ (\cref{thm:constant-coefficient-recurrences}), and the conditions at $n=0,1$ give $A$ and $B$. Since $\sinh t>t$ for $t>0$, we have $\sigma=\sinh s>s$, so $\mu_1=e^{-s}>e^{-\sigma}=e^{h\lambda}$; hence $y_1=e^{h\lambda}y_0\ne\mu_1y_0$ for $y_0\ne0$, and $B\ne0$. Then $|B\mu_2^n|=|B|e^{ns}\to\infty$ while $A\mu_1^n\to0$.
\end{proof}

The root $\mu_1$ is the \term{principal root}: by Taylor's theorem, $s=\sigma-\sigma^3/6+O(\sigma^5)$, so $\mu_1^n=e^{-ns}$ agrees with $e^{n h\lambda}=e^{-n\sigma}$ up to a relative error of the order $n\sigma^3=t_n|\lambda|^3h^2$, as befits a method of order $2$. The root $\mu_2\approx-e^{h|\lambda|}$ is \term{parasitic}: it belongs to the recurrence, not to the differential equation, and the corresponding solution $B\mu_2^n\approx B(-1)^ne^{|\lambda|t_n}$ alternates in sign and grows like the solution of $y'=-\lambda y$, the equation with the opposite sign. Its coefficient is small: $y_1-\mu_1y_0=(e^{-\sigma}-e^{-s})y_0\approx-\sigma^3y_0/6$ and $\mu_2-\mu_1\approx-2$, so $B\approx\sigma^3y_0/12$. On a fixed interval $[0,T]$, with $n\le T/h$, the parasitic term is at most $|B|e^{|\lambda|T}=O(h^3)$, and the method converges as $h\to0$. But for a fixed step it eventually dominates. \Cref{sec:recurrence-computation} shows this for $y'=-y$ and $h=0.1$, where $B=7.5\times10^{-5}$: the parasitic component is $2.0\times10^{-4}$ at $t=1$, overtakes the decaying solution before $t=5$, and is $1.6$ at $t=10$, where the solution is $4.5\times10^{-5}$.

For an oscillation, $\lambda=i\omega$ with $0<h\omega<1$, the roots $ih\omega\pm\sqrt{1-h^2\omega^2}$ both have modulus $1$, and the leapfrog method neither damps nor amplifies, which makes it useful for undamped oscillations and waves. The instability for decaying solutions is called \term{weak}: the method converges, but the parasitic root lies on or near the unit circle and on the wrong side of it.

\subsection{The Root Condition}

A general \term{linear multistep method} with $k$ steps computes $y_{n+k}$ from
\[
\sum_{j=0}^k\alpha_jy_{n+j}=h\sum_{j=0}^k\beta_jf(t_{n+j},y_{n+j}),\qquad\alpha_k=1,
\]
and its properties are expressed through the polynomials $\rho(\zeta)\coloneqq\sum_j\alpha_j\zeta^j$ and $\sigma(\zeta)\coloneqq\sum_j\beta_j\zeta^j$. For the leapfrog method, $\rho(\zeta)=\zeta^2-1$ and $\sigma(\zeta)=2\zeta$. For the trivial equation $y'=0$ the method reduces to the recurrence $\rho(E)y=0$, and its solutions, which are the errors propagated from perturbed starting values, must stay bounded if the method is to converge. By \cref{cor:scalar-recurrence-stability} this is the \term{root condition}: the roots of $\rho$ lie in the closed unit disc, and those on the unit circle are simple. A method that satisfies it is called \term{zero-stable}. The method is \term{consistent} if it is exact for $y=1$ and $y=t$, which means $\rho(1)=0$ and $\rho'(1)=\sigma(1)$.

\begin{theorem}[Dahlquist's equivalence theorem]
\label{thm:dahlquist}
A linear multistep method is convergent, in the sense that for every initial value problem $y'=f(t,y)$, $y(t_0)=y_0$, with $f$ continuous and Lipschitz with respect to $y$, and all starting values that tend to the exact ones as $h\to0$, the approximations converge to the solution uniformly on bounded intervals, if and only if it is consistent and zero-stable.
\end{theorem}

\begin{proof}[sketch]
Necessity: applied to $y'=0$ with perturbed starting values, a root $\zeta$ of $\rho$ with $|\zeta|>1$ produces errors $\varepsilon\zeta^n$, and a double root on the circle errors $\varepsilon n\zeta^n$; with $\varepsilon$ proportional to $h$ and $n=T/h$ they do not tend to zero. Consistency follows from the equations $y'=0$ and $y'=1$. Sufficiency: the errors $e_n$ satisfy $\rho(E)e_n=h\,(\ldots)+\tau_n$, with the differences of the values of $f$, bounded by the Lipschitz constant times the errors, and the local errors $\tau_n$, which are $o(h)$ by consistency. Writing this as a system with the companion matrix of $\rho$, whose powers are bounded by the root condition (\cref{thm:difference-stability}), the variation of constants formula \eqref{eq:discrete-duhamel} and the discrete Gr\"onwall inequality (\cref{lem:discrete-gronwall}) bound the errors by a constant times the starting errors plus $\sum_n|\tau_n|$, which tends to zero. Dahlquist proved the theorem in 1956~\cite{dahlquist1956}; Hairer, N{\o}rsett and Wanner~\cite[Chap.~III]{hairer1993} give the complete proof.
\end{proof}

The leapfrog method has $\rho(\zeta)=(\zeta-1)(\zeta+1)$, with two simple roots on the circle: it is zero-stable, which is why it converges, and its weak instability comes from the root $-1$, which becomes $\mu_2\approx-e^{h|\lambda|}$ when $\lambda\ne0$. Consistency alone does not suffice. The two-step method
\[
y_{n+2}+4y_{n+1}-5y_n=h\bigl(4f(t_{n+1},y_{n+1})+2f(t_n,y_n)\bigr)
\]
is consistent, and even of order $3$, as one checks by inserting the Taylor expansions of $e^t$, but $\rho(\zeta)=\zeta^2+4\zeta-5=(\zeta-1)(\zeta+5)$ violates the root condition: every perturbation is multiplied by about $-5$ per step, and halving the step doubles the number of these multiplications. \Cref{sec:recurrence-computation} shows the result. Dahlquist also proved that the root condition limits the order: a zero-stable explicit method with $k$ steps has order at most $k$~\cite[Chap.~III]{hairer1993}, so the method above, explicit with two steps and order $3$, could not have been zero-stable. The theory of multistep methods is developed by Hairer, N{\o}rsett and Wanner~\cite{hairer1993} and Henrici~\cite{henrici1962discrete}; the numerical solution of differential equations is not among the subjects of the later volumes of this series.

\subsection{Boundary Value Problems}

Discretising a boundary value problem gives a recurrence with conditions at both ends, as in the gambler's ruin. The problem $-u''=f$ on $[0,1]$ with $u(0)=u(1)=0$ becomes, with $h=1/(N+1)$, the grid points $x_n=nh$ and the second difference quotient of \vcref{1}{prop:difference-taylor}\,(iii),
\begin{equation}
-y_{n-1}+2y_n-y_{n+1}=h^2f(x_n)\quad(1\le n\le N),\qquad y_0=y_{N+1}=0 ,
\label{eq:discrete-bvp}
\end{equation}
that is, the linear system $T_N\mathbf y=h^2\mathbf f$ of \cref{sec:special-matrices}, which has exactly one solution because $T_N$ is invertible. The recurrence $\Delta^2y_{n-1}=-h^2f(x_n)$ also shows how the solution is built: two summations, as $-u''=f$ is solved by two integrations, and the inverse of $T_N$ is the discrete Green's function.

\begin{proposition}[Convergence of the discrete boundary value problem]
\label{prop:discrete-bvp}
Let $u\in C^4[0,1]$ solve $-u''=f$ with $u(0)=u(1)=0$, let $M_4\coloneqq\max|u^{(4)}|$, and let $\mathbf y$ solve \eqref{eq:discrete-bvp}. Then
\[
\max_{1\le n\le N}|y_n-u(x_n)|\le\frac{M_4h^2}{96} .
\]
\end{proposition}

\begin{proof}
By \vcref{1}{prop:difference-taylor}\,(iii), the exact values satisfy $-u(x_{n-1})+2u(x_n)-u(x_{n+1})=h^2f(x_n)+h^2\tau_n$ with $\tau_n=-h^2u^{(4)}(\xi_n)/12$, so $|\tau_n|\le M_4h^2/12$, and $u(x_0)=u(x_{N+1})=0$. Subtracting \eqref{eq:discrete-bvp}, the errors $e_n\coloneqq u(x_n)-y_n$ satisfy $T_N\mathbf e=h^2\boldsymbol\tau$, so $\mathbf e=h^2T_N^{-1}\boldsymbol\tau$. By \cref{prop:second-difference-inverse}, $T_N^{-1}$ has the positive entries $\min(i,j)(N+1-\max(i,j))/(N+1)$, and the sum of row $i$ is
\[
\frac{1}{N+1}\Bigl(\sum_{j\le i}j\,(N+1-i)+\sum_{j>i}i\,(N+1-j)\Bigr)=\frac{i\,(N+1-i)}2\le\frac{(N+1)^2}8 .
\]
So $\|T_N^{-1}\|_\infty\le(N+1)^2/8=1/(8h^2)$ (\cref{thm:norm-formulas}), and $\|\mathbf e\|_\infty\le h^2\cdot(1/(8h^2))\cdot M_4h^2/12=M_4h^2/96$.
\end{proof}

The proof has the structure of every convergence proof for discretisations: a local error of order $h^2$, made by the exact solution in the difference equation, and a stability bound, here $\|(T_N/h^2)^{-1}\|_\infty\le1/8$ independently of $N$, which keeps the local errors from being amplified. For $f(x)=\sin\pi x$ the solution is $u=\sin(\pi x)/\pi^2$, and the sine vector is an eigenvector of $T_N$ (\cref{thm:second-difference-eigen}), so $y_n=h^2\sin(\pi x_n)/(4\sin^2(\pi h/2))$ exactly. With $z\coloneqq\pi h/2$, $y_n-u(x_n)=(\sin(\pi x_n)/\pi^2)\bigl((z/\sin z)^2-1\bigr)$, and $(z/\sin z)^2=1+z^2/3+O(z^4)$, so the largest error is $h^2/12+O(h^4)$, compared with the bound $\pi^2h^2/96=0.103h^2$ of the proposition, in which $M_4=\pi^2$. Volume~XII returns to difference quotients in its treatment of numerical differentiation.

\subsection{Where the Analogy Ends}

The convergence proof of Euler's method (\cref{thm:euler-convergence}) can be read in the language of this chapter. The errors satisfy the first-order recurrence $e_{n+1}=e_n+h\bigl(f(t_n,y_n)-f(t_n,y(t_n))\bigr)-\tau_{n+1}$, and for the test equation $f=\lambda y$ this is $e_{n+1}=(1+h\lambda)e_n-\tau_{n+1}$, whose solution is $e_n=(1+h\lambda)^ne_0-\sum_{m<n}(1+h\lambda)^{n-1-m}\tau_{m+1}$ by \cref{thm:first-order-variable}. With $|1+h\lambda|\le e^{h|\lambda|}$ and $|\tau_m|\le M_2h^2/2$, summing the geometric series $\sum_{m<n}e^{h|\lambda|(n-1-m)}$ gives the bound of the theorem with $L=|\lambda|$: $n$ local errors of order $h^2$, each amplified by at most $e^{|\lambda|T}$.

The correspondence between differential and difference equations through $\mu=e^{h\lambda}$ is not a bijection, and the recurrences that fall outside it are instructive. A real number $e^{h\lambda}$ is positive, so the exact discretisation of a real scalar equation never produces a negative characteristic root; a recurrence $y_{n+1}=ay_n$ with $a<0$, whose solutions alternate, has no real differential equation of the first order behind it, and it arises only as an approximation, such as Euler's method with $h\lambda<-1$. A matrix $e^{hA}$ is always invertible (\cref{prop:exp-basic}\,(iii)), so no exact discretisation has a characteristic root $0$, and the forgetting recurrences of \cref{sec:constant-coefficient-recurrences} have no continuous counterpart either. And the step from the differential equation to a recurrence of higher order, as in the leapfrog method, adds roots that the differential equation does not have. In each case the recurrence carries behaviour of its own, and the next section shows how much richer this behaviour becomes when the recurrence is nonlinear.

\section{Nonlinear Recurrences: A Glimpse}
\label{sec:nonlinear-recurrences}

\subsection{Fixed Points and Their Stability}

For a nonlinear recurrence $x_{n+1}=f(x_n)$ there is no general formula for the solutions, and as for nonlinear differential equations (\cref{sec:autonomous}) the first questions concern equilibria. A fixed point $x^*=f(x^*)$ is \term{stable} if for every $\varepsilon>0$ there is $\delta>0$ such that $|x_0-x^*|<\delta$ implies $|x_n-x^*|<\varepsilon$ for all $n$, \term{asymptotically stable} if moreover $x_n\to x^*$ for all $x_0$ near $x^*$, and \term{unstable} if it is not stable, as in \cref{def:stability}. The linearisation decides the question unless $|f'(x^*)|=1$.

\begin{theorem}[Linearised stability of fixed points]
\label{thm:map-linearised}
Let $f$ be continuously differentiable on an open interval containing the fixed point $x^*$.
\begin{enumerate}
\item If $|f'(x^*)|<1$, then $x^*$ is asymptotically stable; for every $q$ with $|f'(x^*)|<q<1$ there is $\delta>0$ such that $|x_n-x^*|\le q^n|x_0-x^*|$ whenever $|x_0-x^*|\le\delta$.
\item If $|f'(x^*)|>1$, then $x^*$ is unstable.
\end{enumerate}
\end{theorem}

\begin{proof}
(i) By continuity of $f'$ there is $\delta>0$ with $|f'(x)|\le q$ on $I\coloneqq[x^*-\delta,x^*+\delta]$. For $x\in I$, the mean value theorem gives $|f(x)-x^*|=|f(x)-f(x^*)|\le q|x-x^*|\le\delta$, so $f$ maps $I$ into itself, and \vcref{1}{prop:fixed-point-contraction} gives the bound; stability follows with $\min(\delta,\varepsilon)$ in the place of $\delta$. (ii) Choose $q$ with $1<q<|f'(x^*)|$ and $\delta>0$ with $|f'(x)|\ge q$ on $I$. If $0<|x_0-x^*|<\delta$, then as long as the iterates stay in $I$, the mean value theorem gives $|x_{n+1}-x^*|\ge q|x_n-x^*|$, so $|x_n-x^*|\ge q^n|x_0-x^*|$, and some iterate leaves $I$. So the definition of stability fails for $\varepsilon=\delta$.
\end{proof}

The theorem is the discrete counterpart of \cref{thm:linearised-stability}, with $|f'(x^*)|<1$ in the place of $f'(y^*)<0$, and the cobweb diagrams of \cref{fig:cobweb-linear} show the four cases of the linearisation. When $|f'(x^*)|=1$, anything can happen. For $f=\sin$ and $x^*=0$, $f'(0)=1$, and the iterates $x_{n+1}=\sin x_n$ from $x_0=1$ decrease, since $0<\sin x<x$ for $0<x\le1$ (\vcref{1}{lem:sin-bounds}), to a limit that is a fixed point and hence $0$; but they do so very slowly, like $\sqrt{3/n}$~\cite[Chap.~8]{debruijn1981}. For $f=\sinh$, with $f'(0)=1$ as well, the iterates from any $x_0>0$ increase without bound, since $\sinh x>x$.

Newton's method for $g(x)=0$ is the recurrence $x_{n+1}=N(x_n)$ with $N(x)\coloneqq x-g(x)/g'(x)$ (\vcref{1}{sec:newton}). If $g$ is twice continuously differentiable and $c$ is a simple zero, then $N'(x)=g(x)g''(x)/g'(x)^2$ vanishes at $c$: the fixed point is \term{superattracting}, and by Taylor's theorem $|N(x)-c|\le C|x-c|^2$ near $c$, the quadratic convergence of \vcref{1}{thm:newton-quadratic}. Iterations with $f'(x^*)=0$ are the fastest that the linearisation allows, and much of the design of iterative methods aims at making the derivative at the fixed point small.

A point $p$ is \term{periodic} with period $m$ if $f^m(p)=p$, where $f^m$ is the $m$-fold composition, and $m$ is the smallest such number; its orbit $p,f(p),\ldots,f^{m-1}(p)$ is a \term{cycle}. A cycle is a fixed point of $f^m$, and by the chain rule $(f^m)'(p)=f'(p)f'(f(p))\cdots f'(f^{m-1}(p))$, the product of the derivatives along the cycle; it is asymptotically stable if this product has modulus less than $1$, by the theorem applied to $f^m$.

\subsection{The Logistic Map}

The \term{logistic map} $f_r(x)\coloneqq rx(1-x)$, $0<r\le4$, maps $[0,1]$ into $[0,r/4]\subseteq[0,1]$, since $x(1-x)\le1/4$. As a model of a population with nonoverlapping generations, it was studied by May~\cite{may1976}, who drew attention to its complicated behaviour. It is also Euler's method for the logistic differential equation of \cref{sec:separable}: Euler's method for $y'=y(1-y)$ with the step $h$ gives $y_{n+1}=y_n+hy_n(1-y_n)$, and the substitution $y_n=((1+h)/h)\,x_n$ turns it into $x_{n+1}=(1+h)\,x_n(1-x_n)$, the logistic map with $r=1+h$.

\begin{proposition}[Fixed points and the first period doubling]
\label{prop:logistic-two-cycle}
Let $0<r\le4$.
\begin{enumerate}
\item If $r<1$, then $f_r^n(x)\to0$ for every $x\in[0,1]$.
\item If $r>1$, the fixed points in $[0,1]$ are $0$, which is unstable, and $x^*=1-1/r$, with $f_r'(x^*)=2-r$; so $x^*$ is asymptotically stable for $1<r<3$ and unstable for $r>3$.
\item For $r>3$, $f_r$ has exactly one cycle of period $2$, the points
\[
p_\pm=\frac{r+1\pm\sqrt{(r+1)(r-3)}}{2r}\in(0,1),
\]
and its multiplier is $f_r'(p_+)f_r'(p_-)=-r^2+2r+4$. The cycle is asymptotically stable for $3<r<1+\sqrt6$ and unstable for $r>1+\sqrt6$.
\end{enumerate}
\end{proposition}

\begin{proof}
(i) $0\le f_r(x)\le rx$ on $[0,1]$, so $f_r^n(x)\le r^nx\to0$. (ii) $rx(1-x)=x$ gives $x=0$ or $x=1-1/r$. $f_r'(x)=r(1-2x)$, so $f_r'(0)=r>1$ and $f_r'(x^*)=r(1-2+2/r)=2-r$, and \cref{thm:map-linearised} applies. (iii) The equation $f_r(f_r(x))=x$ is a polynomial equation of degree $4$ whose roots include the fixed points $0$ and $x^*$; dividing them out, one finds that the points of period $2$ are the roots of
\[
r^2x^2-r(r+1)x+(r+1)=0 ,
\]
which are $p_\pm$; they are real and different exactly when $(r+1)(r-3)>0$, that is, $r>3$, and $f_r$ maps each of them to the other, since $f_r(p_+)$ is again a root of $f_r(f_r(x))=x$ and not a fixed point. They lie in $(0,1)$: $p_->0$ is clear, and $p_+<1$ is equivalent to $\sqrt{(r+1)(r-3)}<r-1$, that is, $r^2-2r-3<r^2-2r+1$. From $p_++p_-=(r+1)/r$ and $p_+p_-=(r+1)/r^2$,
\[
f_r'(p_+)f_r'(p_-)=r^2\bigl(1-2(p_++p_-)+4p_+p_-\bigr)=r^2-2r(r+1)+4(r+1)=-r^2+2r+4 .
\]
This lies in $(-1,1)$ if and only if $r^2-2r-3>0$ and $r^2-2r-5<0$, that is, $3<r<1+\sqrt6$; it is less than $-1$ for $r>1+\sqrt6$. \Cref{thm:map-linearised}, applied to $f_r^2$, gives the stability.
\end{proof}

As $r$ passes $3$, the fixed point loses its stability through the multiplier $-1$, and a stable cycle of period $2$ is born around it: a \term{period-doubling bifurcation}. \Cref{fig:cobweb-logistic} shows the two regimes. At $r=1+\sqrt6=3.449490$ the cycle of period $2$ loses its stability in the same way, through the multiplier $-1$, and a stable cycle of period $4$ appears, then one of period $8$, and so on. We computed the parameters $r_k$ at which the stable cycle of period $2^{k-1}$ loses its stability, by solving $f_r^{p}(x)=x$ and $(f_r^{p})'(x)=-1$ for $(x,r)$ with Newton's method, $p=2^{k-1}$; \cref{tab:period-doubling} lists them. Extrapolating with the last ratio, they accumulate at about $3.5699457$, and the ratios of consecutive distances approach $4.6692$. Feigenbaum discovered that this ratio has the same limit $\delta=4.6692016\ldots$ for a large class of maps with a single quadratic maximum~\cite{feigenbaum1978}. Beyond the accumulation point the orbits look irregular for most $r$, interrupted by windows of periodic behaviour; the most conspicuous of them, near $r=3.83$, contains a stable cycle of period $3$. \Cref{fig:logistic-bifurcation} shows the whole picture.

\begin{figure}
\sdpanel{A}{\begin{tikzpicture}
\begin{axis}[sdplot, width=0.46\textwidth, height=0.42\textwidth, xmin=0, xmax=1, ymin=0, ymax=1,
  xtick={0,0.5,1}, ytick={0,0.5,1}, xlabel={$x_n$}, ylabel={$x_{n+1}$}]
\addplot[black!45, line width=0.4pt, domain=0:1, samples=2] {x};
\addplot[black, line width=0.7pt, domain=0:1, samples=60, smooth] {2.8*x*(1-x)};
\addplot[black, line width=0.35pt] coordinates {(0.1000,0.0000) (0.1000,0.2520) (0.2520,0.2520) (0.2520,0.5278) (0.5278,0.5278) (0.5278,0.6978) (0.6978,0.6978) (0.6978,0.5904) (0.5904,0.5904) (0.5904,0.6771) (0.6771,0.6771) (0.6771,0.6122) (0.6122,0.6122) (0.6122,0.6648) (0.6648,0.6648) (0.6648,0.6240) (0.6240,0.6240) (0.6240,0.6570) (0.6570,0.6570) (0.6570,0.6310) (0.6310,0.6310) (0.6310,0.6519) (0.6519,0.6519) (0.6519,0.6354) (0.6354,0.6354) (0.6354,0.6487) (0.6487,0.6487) (0.6487,0.6381) (0.6381,0.6381) (0.6381,0.6466) (0.6466,0.6466) (0.6466,0.6398) (0.6398,0.6398) (0.6398,0.6453) (0.6453,0.6453) (0.6453,0.6409) (0.6409,0.6409) (0.6409,0.6444) (0.6444,0.6444) (0.6444,0.6416) (0.6416,0.6416) (0.6416,0.6438) (0.6438,0.6438) (0.6438,0.6421) (0.6421,0.6421) (0.6421,0.6435) (0.6435,0.6435) (0.6435,0.6423) (0.6423,0.6423) (0.6423,0.6433) (0.6433,0.6433)};
\end{axis}
\end{tikzpicture}}\hfill
\sdpanel{B}{\begin{tikzpicture}
\begin{axis}[sdplot, width=0.46\textwidth, height=0.42\textwidth, xmin=0, xmax=1, ymin=0, ymax=1,
  xtick={0,0.5,1}, ytick={0,0.5,1}, xlabel={$x_n$}, ylabel={$x_{n+1}$}]
\addplot[black!45, line width=0.4pt, domain=0:1, samples=2] {x};
\addplot[black, line width=0.7pt, domain=0:1, samples=60, smooth] {3.3*x*(1-x)};
\addplot[black, line width=0.35pt] coordinates {(0.1000,0.0000) (0.1000,0.2970) (0.2970,0.2970) (0.2970,0.6890) (0.6890,0.6890) (0.6890,0.7071) (0.7071,0.7071) (0.7071,0.6835) (0.6835,0.6835) (0.6835,0.7139) (0.7139,0.7139) (0.7139,0.6740) (0.6740,0.6740) (0.6740,0.7251) (0.7251,0.7251) (0.7251,0.6577) (0.6577,0.6577) (0.6577,0.7429) (0.7429,0.7429) (0.7429,0.6303) (0.6303,0.6303) (0.6303,0.7690) (0.7690,0.7690) (0.7690,0.5863) (0.5863,0.5863) (0.5863,0.8004) (0.8004,0.8004) (0.8004,0.5271) (0.5271,0.5271)};
\addplot[black, line width=1.1pt] coordinates {(0.4794,0.4794) (0.4794,0.8236) (0.8236,0.8236) (0.8236,0.4794) (0.4794,0.4794)};
\end{axis}
\end{tikzpicture}}
\caption{Cobweb diagrams of the logistic map from $x_0=0.1$. (A) $r=2.8$: the iterates spiral into the fixed point $x^*=0.6429$, where $f_r'(x^*)=-0.8$. (B) $r=3.3$: the fixed point $x^*=0.6970$ is unstable, $f_r'(x^*)=-1.3$, and the iterates spiral out towards the stable cycle $\{0.4794,0.8236\}$ of period $2$ (thick square), whose multiplier is $-0.29$ (\cref{prop:logistic-two-cycle}).\label{fig:cobweb-logistic}}
\end{figure}

\begin{table}
\sdfloat{%
\begin{sdtable}{colspec={rlll}, row{1}={font=\small\bfseries}}
\toprule
$k$ & $r_k$ & $r_k-r_{k-1}$ & $(r_{k-1}-r_{k-2})/(r_k-r_{k-1})$ \\
\midrule
$1$ & $3$ & & \\
$2$ & $3.449490$ & $4.495\times10^{-1}$ & \\
$3$ & $3.544090$ & $9.460\times10^{-2}$ & $4.7514$ \\
$4$ & $3.564407$ & $2.032\times10^{-2}$ & $4.6563$ \\
$5$ & $3.568759$ & $4.352\times10^{-3}$ & $4.6682$ \\
$6$ & $3.569692$ & $9.322\times10^{-4}$ & $4.6687$ \\
$7$ & $3.569891$ & $1.996\times10^{-4}$ & $4.6691$ \\
$8$ & $3.569934$ & $4.276\times10^{-5}$ & $4.6692$ \\
\bottomrule
\end{sdtable}}{The parameters $r_k$ at which the stable cycle of period $2^{k-1}$ of the logistic map loses its stability, computed in binary64; $r_1=3$ and $r_2=1+\sqrt6$ are exact (\cref{prop:logistic-two-cycle}). The ratios of consecutive distances approach Feigenbaum's constant $\delta=4.6692\ldots$\label{tab:period-doubling}}
\end{table}

\begin{figure}
\sdfloat{%
\begin{tikzpicture}
\begin{axis}[sdplot, width=\textwidth, height=0.52\textwidth, xmin=2.8, xmax=4, ymin=0, ymax=1,
  xtick={2.8,3,3.2,3.4,3.6,3.8,4}, xticklabels={$2.8$,$3$,$3.2$,$3.4$,$3.6$,$3.8$,$4$}, ytick={0,0.5,1}, xlabel={$r$}, ylabel={$x$}]
\addplot[only marks, mark=*, mark size=0.22pt, black!55] coordinates {(2.8000,0.6429) (2.8043,0.6434) (2.8086,0.6439) (2.8129,0.6445) (2.8171,0.6450) (2.8214,0.6456) (2.8257,0.6461) (2.8300,0.6466) (2.8343,0.6472) (2.8386,0.6477) (2.8429,0.6482) (2.8471,0.6488) (2.8514,0.6493) (2.8557,0.6498) (2.8600,0.6503) (2.8643,0.6509) (2.8686,0.6514) (2.8729,0.6519) (2.8771,0.6524) (2.8814,0.6529) (2.8857,0.6535) (2.8900,0.6540) (2.8943,0.6545) (2.8986,0.6550) (2.9029,0.6555) (2.9071,0.6560) (2.9114,0.6565) (2.9157,0.6570) (2.9200,0.6575) (2.9243,0.6580) (2.9286,0.6585) (2.9329,0.6590) (2.9371,0.6595) (2.9414,0.6600) (2.9457,0.6605) (2.9500,0.6610) (2.9543,0.6615) (2.9586,0.6620) (2.9629,0.6625) (2.9671,0.6630) (2.9714,0.6635) (2.9757,0.6639) (2.9800,0.6644) (2.9843,0.6649) (2.9886,0.6654) (2.9929,0.6659) (2.9971,0.6663) (3.0014,0.6792) (3.0014,0.6540) (3.0057,0.6915) (3.0057,0.6412) (3.0100,0.6994) (3.0100,0.6328) (3.0143,0.7056) (3.0143,0.6262) (3.0186,0.7109) (3.0186,0.6204) (3.0229,0.7156) (3.0229,0.6152) (3.0271,0.7198) (3.0271,0.6106) (3.0314,0.7236) (3.0314,0.6062) (3.0357,0.7272) (3.0357,0.6022) (3.0400,0.7306) (3.0400,0.5984) (3.0443,0.7338) (3.0443,0.5947) (3.0486,0.7367) (3.0486,0.5913) (3.0529,0.7396) (3.0529,0.5880) (3.0571,0.7423) (3.0571,0.5848) (3.0614,0.7449) (3.0614,0.5817) (3.0657,0.7474) (3.0657,0.5788) (3.0700,0.7498) (3.0700,0.5759) (3.0743,0.7521) (3.0743,0.5732) (3.0786,0.7544) (3.0786,0.5705) (3.0829,0.7565) (3.0829,0.5679) (3.0871,0.7586) (3.0871,0.5653) (3.0914,0.7607) (3.0914,0.5628) (3.0957,0.7626) (3.0957,0.5604) (3.1000,0.7646) (3.1000,0.5580) (3.1043,0.7664) (3.1043,0.5557) (3.1086,0.7683) (3.1086,0.5534) (3.1129,0.7701) (3.1129,0.5512) (3.1171,0.7718) (3.1171,0.5490) (3.1214,0.7735) (3.1214,0.5469) (3.1257,0.7752) (3.1257,0.5448) (3.1300,0.7768) (3.1300,0.5427) (3.1343,0.7784) (3.1343,0.5407) (3.1386,0.7800) (3.1386,0.5387) (3.1429,0.7815) (3.1429,0.5367) (3.1471,0.7830) (3.1471,0.5348) (3.1514,0.7845) (3.1514,0.5329) (3.1557,0.7859) (3.1557,0.5310) (3.1600,0.7873) (3.1600,0.5291) (3.1643,0.7887) (3.1643,0.5273) (3.1686,0.7901) (3.1686,0.5255) (3.1729,0.7914) (3.1729,0.5237) (3.1771,0.7927) (3.1771,0.5220) (3.1814,0.7940) (3.1814,0.5203) (3.1857,0.7953) (3.1857,0.5186) (3.1900,0.7966) (3.1900,0.5169) (3.1943,0.7978) (3.1943,0.5152) (3.1986,0.7991) (3.1986,0.5136) (3.2029,0.8003) (3.2029,0.5120) (3.2071,0.8014) (3.2071,0.5104) (3.2114,0.8026) (3.2114,0.5088) (3.2157,0.8038) (3.2157,0.5072) (3.2200,0.8049) (3.2200,0.5057) (3.2243,0.8060) (3.2243,0.5041) (3.2286,0.8071) (3.2286,0.5026) (3.2329,0.8082) (3.2329,0.5011) (3.2371,0.8093) (3.2371,0.4996) (3.2414,0.8103) (3.2414,0.4982) (3.2457,0.8114) (3.2457,0.4967) (3.2500,0.8124) (3.2500,0.4953) (3.2543,0.8134) (3.2543,0.4938) (3.2586,0.8145) (3.2586,0.4924) (3.2629,0.8155) (3.2629,0.4910) (3.2671,0.8164) (3.2671,0.4896) (3.2714,0.8174) (3.2714,0.4883) (3.2757,0.8184) (3.2757,0.4869) (3.2800,0.8193) (3.2800,0.4856) (3.2843,0.8203) (3.2843,0.4842) (3.2886,0.8212) (3.2886,0.4829) (3.2929,0.8221) (3.2929,0.4816) (3.2971,0.8230) (3.2971,0.4803) (3.3014,0.8239) (3.3014,0.4790) (3.3057,0.8248) (3.3057,0.4777) (3.3100,0.8257) (3.3100,0.4765) (3.3143,0.8265) (3.3143,0.4752) (3.3186,0.8274) (3.3186,0.4739) (3.3229,0.8282) (3.3229,0.4727) (3.3271,0.8291) (3.3271,0.4715) (3.3314,0.8299) (3.3314,0.4703) (3.3357,0.8307) (3.3357,0.4691) (3.3400,0.8315) (3.3400,0.4679) (3.3443,0.8324) (3.3443,0.4667) (3.3486,0.8332) (3.3486,0.4655) (3.3529,0.8339) (3.3529,0.4643) (3.3571,0.8347) (3.3571,0.4631) (3.3614,0.8355) (3.3614,0.4620) (3.3657,0.8363) (3.3657,0.4608) (3.3700,0.8370) (3.3700,0.4597) (3.3743,0.8378) (3.3743,0.4586) (3.3786,0.8385) (3.3786,0.4575) (3.3829,0.8393) (3.3829,0.4563) (3.3871,0.8400) (3.3871,0.4552) (3.3914,0.8407) (3.3914,0.4541) (3.3957,0.8414) (3.3957,0.4530) (3.4000,0.8422) (3.4000,0.4520) (3.4043,0.8429) (3.4043,0.4509) (3.4086,0.8436) (3.4086,0.4498) (3.4129,0.8443) (3.4129,0.4488) (3.4171,0.8449) (3.4171,0.4477) (3.4214,0.8456) (3.4214,0.4467) (3.4257,0.8463) (3.4257,0.4456) (3.4300,0.8470) (3.4300,0.4446) (3.4343,0.8476) (3.4343,0.4436) (3.4386,0.8483) (3.4386,0.4425) (3.4429,0.8489) (3.4429,0.4415) (3.4471,0.8496) (3.4471,0.4405) (3.4514,0.8548) (3.4514,0.4284) (3.4514,0.8451) (3.4514,0.4517) (3.4557,0.8587) (3.4557,0.4193) (3.4557,0.8414) (3.4557,0.4611) (3.4600,0.8613) (3.4600,0.4132) (3.4600,0.8390) (3.4600,0.4675) (3.4643,0.8635) (3.4643,0.4084) (3.4643,0.8370) (3.4643,0.4727) (3.4686,0.8653) (3.4686,0.4042) (3.4686,0.8353) (3.4686,0.4772) (3.4729,0.8670) (3.4729,0.4005) (3.4729,0.8338) (3.4729,0.4812) (3.4771,0.8685) (3.4771,0.3971) (3.4771,0.8325) (3.4771,0.4849) (3.4814,0.8699) (3.4814,0.3941) (3.4814,0.8313) (3.4814,0.4883) (3.4857,0.8712) (3.4857,0.3912) (3.4857,0.8302) (3.4857,0.4914) (3.4900,0.8724) (3.4900,0.3885) (3.4900,0.8291) (3.4900,0.4944) (3.4943,0.8735) (3.4943,0.3860) (3.4943,0.8282) (3.4943,0.4973) (3.4986,0.8746) (3.4986,0.3836) (3.4986,0.8272) (3.4986,0.5000) (3.5029,0.8757) (3.5029,0.3813) (3.5029,0.8264) (3.5029,0.5026) (3.5071,0.8767) (3.5071,0.3791) (3.5071,0.8255) (3.5071,0.5051) (3.5114,0.8777) (3.5114,0.3770) (3.5114,0.8248) (3.5114,0.5075) (3.5157,0.8786) (3.5157,0.3750) (3.5157,0.8240) (3.5157,0.5098) (3.5200,0.8795) (3.5200,0.3731) (3.5200,0.8233) (3.5200,0.5121) (3.5243,0.8804) (3.5243,0.3712) (3.5243,0.8226) (3.5243,0.5143) (3.5286,0.8812) (3.5286,0.3694) (3.5286,0.8220) (3.5286,0.5164) (3.5329,0.8820) (3.5329,0.3676) (3.5329,0.8213) (3.5329,0.5184) (3.5371,0.8828) (3.5371,0.3659) (3.5371,0.8207) (3.5371,0.5204) (3.5414,0.8836) (3.5414,0.3643) (3.5414,0.8201) (3.5414,0.5224) (3.5457,0.8857) (3.5457,0.3590) (3.5457,0.8159) (3.5457,0.5325) (3.5457,0.3672) (3.5457,0.8239) (3.5457,0.5145) (3.5500,0.8874) (3.5500,0.3548) (3.5500,0.8127) (3.5500,0.5405) (3.5500,0.8817) (3.5500,0.3703) (3.5500,0.8278) (3.5500,0.5060) (3.5543,0.8886) (3.5543,0.3519) (3.5543,0.8106) (3.5543,0.5456) (3.5543,0.8812) (3.5543,0.3721) (3.5543,0.8305) (3.5543,0.5004) (3.5586,0.8896) (3.5586,0.3495) (3.5586,0.8091) (3.5586,0.5497) (3.5586,0.8809) (3.5586,0.3735) (3.5586,0.8327) (3.5586,0.4959) (3.5629,0.8905) (3.5629,0.3475) (3.5629,0.8078) (3.5629,0.5531) (3.5629,0.8807) (3.5629,0.3745) (3.5629,0.8346) (3.5629,0.4919) (3.5671,0.8918) (3.5671,0.3443) (3.5671,0.8053) (3.5671,0.5594) (3.5671,0.8792) (3.5671,0.3789) (3.5671,0.8394) (3.5671,0.4808) (3.5671,0.3479) (3.5671,0.8093) (3.5671,0.5506) (3.5671,0.8827) (3.5671,0.3695) (3.5714,0.8906) (3.5714,0.3480) (3.5714,0.8104) (3.5714,0.5488) (3.5714,0.8844) (3.5714,0.3653) (3.5714,0.8280) (3.5714,0.5086) (3.5714,0.3424) (3.5714,0.8041) (3.5714,0.5625) (3.5714,0.8789) (3.5714,0.3801) (3.5757,0.8806) (3.5757,0.3759) (3.5757,0.8389) (3.5757,0.4833) (3.5757,0.8929) (3.5757,0.3418) (3.5757,0.8045) (3.5757,0.5624) (3.5757,0.4797) (3.5757,0.3432) (3.5757,0.5591) (3.5757,0.3737) (3.5757,0.8369) (3.5800,0.8950) (3.5800,0.3366) (3.5800,0.7994) (3.5800,0.5741) (3.5800,0.8753) (3.5800,0.3907) (3.5800,0.8522) (3.5800,0.4509) (3.5800,0.8864) (3.5800,0.3606) (3.5800,0.8255) (3.5800,0.5158) (3.5800,0.3390) (3.5843,0.8865) (3.5843,0.3606) (3.5843,0.8264) (3.5843,0.5143) (3.5843,0.8953) (3.5843,0.3359) (3.5843,0.7995) (3.5843,0.5745) (3.5843,0.8762) (3.5843,0.3889) (3.5843,0.8518) (3.5843,0.4524) (3.5843,0.3567) (3.5886,0.8758) (3.5886,0.3902) (3.5886,0.8539) (3.5886,0.4477) (3.5886,0.8873) (3.5886,0.3588) (3.5886,0.8256) (3.5886,0.5168) (3.5886,0.8961) (3.5886,0.3340) (3.5886,0.7983) (3.5886,0.5779) (3.5886,0.8549) (3.5929,0.7945) (3.5929,0.5866) (3.5929,0.8713) (3.5929,0.4029) (3.5929,0.8643) (3.5929,0.4213) (3.5929,0.8760) (3.5929,0.3904) (3.5929,0.8551) (3.5929,0.4453) (3.5929,0.8875) (3.5929,0.3588) (3.5929,0.8266) (3.5971,0.8780) (3.5971,0.3854) (3.5971,0.8520) (3.5971,0.4535) (3.5971,0.8915) (3.5971,0.3479) (3.5971,0.8161) (3.5971,0.5399) (3.5971,0.3422) (3.5971,0.8097) (3.5971,0.5543) (3.5971,0.8887) (3.5971,0.3559) (3.6014,0.8237) (3.6014,0.5231) (3.6014,0.8984) (3.6014,0.3286) (3.6014,0.7946) (3.6014,0.5878) (3.6014,0.8726) (3.6014,0.4004) (3.6014,0.8646) (3.6014,0.4215) (3.6014,0.8782) (3.6014,0.3853) (3.6014,0.8530) (3.6057,0.9014) (3.6057,0.3204) (3.6057,0.7852) (3.6057,0.6082) (3.6057,0.8592) (3.6057,0.4362) (3.6057,0.8868) (3.6057,0.3621) (3.6057,0.8328) (3.6057,0.5020) (3.6100,0.8754) (3.6100,0.3937) (3.6100,0.8617) (3.6100,0.4301) (3.6100,0.8849) (3.6100,0.3678) (3.6100,0.8394) (3.6100,0.4867) (3.6100,0.9019) (3.6100,0.3195) (3.6100,0.7849) (3.6100,0.6095) (3.6100,0.4367) (3.6143,0.7906) (3.6143,0.5984) (3.6143,0.8686) (3.6143,0.4125) (3.6143,0.8759) (3.6143,0.3928) (3.6143,0.8620) (3.6143,0.4298) (3.6143,0.8858) (3.6143,0.3657) (3.6143,0.8383) (3.6143,0.4898) (3.6143,0.9032) (3.6186,0.8469) (3.6186,0.4692) (3.6186,0.9012) (3.6186,0.3222) (3.6186,0.7902) (3.6186,0.5999) (3.6186,0.8685) (3.6186,0.4132) (3.6186,0.8774) (3.6186,0.3893) (3.6186,0.8603) (3.6186,0.4349) (3.6186,0.8893) (3.6229,0.8494) (3.6229,0.4635) (3.6229,0.9009) (3.6229,0.3235) (3.6229,0.7928) (3.6229,0.5951) (3.6229,0.8730) (3.6229,0.4018) (3.6229,0.4077) (3.6229,0.8748) (3.6229,0.3967) (3.6229,0.8670) (3.6229,0.4177) (3.6271,0.9068) (3.6271,0.3066) (3.6271,0.7711) (3.6271,0.6401) (3.6271,0.8356) (3.6271,0.4984) (3.6314,0.9079) (3.6314,0.3038) (3.6314,0.7680) (3.6314,0.6470) (3.6314,0.8294) (3.6314,0.5138) (3.6314,0.7709) (3.6314,0.6413) (3.6314,0.8353) (3.6314,0.4995) (3.6357,0.8322) (3.6357,0.5077) (3.6357,0.9087) (3.6357,0.3016) (3.6357,0.7658) (3.6357,0.6521) (3.6357,0.8249) (3.6357,0.5252) (3.6357,0.3078) (3.6357,0.7746) (3.6357,0.6347) (3.6357,0.8430) (3.6357,0.4812) (3.6400,0.8509) (3.6400,0.4617) (3.6400,0.9047) (3.6400,0.3140) (3.6400,0.7840) (3.6400,0.6164) (3.6400,0.8607) (3.6400,0.4364) (3.6400,0.8953) (3.6400,0.3412) (3.6400,0.8183) (3.6400,0.5413) (3.6400,0.3165) (3.6443,0.8112) (3.6443,0.5581) (3.6443,0.8988) (3.6443,0.3316) (3.6443,0.8077) (3.6443,0.5659) (3.6443,0.8952) (3.6443,0.3418) (3.6443,0.8199) (3.6443,0.5382) (3.6443,0.9058) (3.6443,0.3111) (3.6443,0.7810) (3.6486,0.9118) (3.6486,0.2933) (3.6486,0.7563) (3.6486,0.6724) (3.6486,0.8037) (3.6486,0.5757) (3.6486,0.8912) (3.6486,0.3536) (3.6486,0.8340) (3.6486,0.5052) (3.6486,0.6743) (3.6486,0.8013) (3.6486,0.5809) (3.6529,0.8444) (3.6529,0.4799) (3.6529,0.9117) (3.6529,0.2939) (3.6529,0.7581) (3.6529,0.6698) (3.6529,0.8078) (3.6529,0.5670) (3.6529,0.8968) (3.6529,0.3381) (3.6529,0.8174) (3.6529,0.5451) (3.6529,0.9058) (3.6571,0.8072) (3.6571,0.5692) (3.6571,0.8968) (3.6571,0.3385) (3.6571,0.8189) (3.6571,0.5423) (3.6571,0.9077) (3.6571,0.3063) (3.6571,0.7770) (3.6571,0.6337) (3.6571,0.8489) (3.6571,0.4690) (3.6571,0.9108) (3.6614,0.8611) (3.6614,0.4381) (3.6614,0.9013) (3.6614,0.3257) (3.6614,0.8041) (3.6614,0.5768) (3.6614,0.8938) (3.6614,0.3476) (3.6614,0.8303) (3.6614,0.5159) (3.6614,0.9144) (3.6614,0.2865) (3.6614,0.7485) (3.6657,0.8138) (3.6657,0.5554) (3.6657,0.9052) (3.6657,0.3146) (3.6657,0.7905) (3.6657,0.6071) (3.6657,0.8743) (3.6657,0.4027) (3.6657,0.8818) (3.6657,0.3822) (3.6657,0.8656) (3.6657,0.4266) (3.6657,0.8967) (3.6700,0.7894) (3.6700,0.6102) (3.6700,0.8730) (3.6700,0.4070) (3.6700,0.8858) (3.6700,0.3713) (3.6700,0.8568) (3.6700,0.4504) (3.6700,0.9085) (3.6700,0.3052) (3.6700,0.7782) (3.6700,0.6335) (3.6700,0.8521) (3.6743,0.7524) (3.6743,0.6845) (3.6743,0.7935) (3.6743,0.6020) (3.6743,0.8803) (3.6743,0.3870) (3.6743,0.8717) (3.6743,0.4110) (3.6743,0.8895) (3.6743,0.3613) (3.6743,0.8479) (3.6743,0.4739) (3.6743,0.9161) (3.6786,0.9027) (3.6786,0.3230) (3.6786,0.8044) (3.6786,0.5788) (3.6786,0.8968) (3.6786,0.3404) (3.6786,0.8259) (3.6786,0.5289) (3.6786,0.9166) (3.6786,0.2813) (3.6786,0.7437) (3.6786,0.7011) (3.6786,0.7709) (3.6829,0.2692) (3.6829,0.7244) (3.6829,0.7352) (3.6829,0.7170) (3.6829,0.7473) (3.6829,0.6955) (3.6829,0.7799) (3.6829,0.6322) (3.6829,0.8564) (3.6829,0.4530) (3.6829,0.9126) (3.6829,0.2939) (3.6829,0.7642) (3.6871,0.5092) (3.6871,0.9215) (3.6871,0.2668) (3.6871,0.7213) (3.6871,0.7412) (3.6871,0.7072) (3.6871,0.7635) (3.6871,0.6659) (3.6871,0.8204) (3.6871,0.5434) (3.6871,0.9149) (3.6871,0.2872) (3.6871,0.7549) (3.6914,0.8773) (3.6914,0.3972) (3.6914,0.8839) (3.6914,0.3789) (3.6914,0.8687) (3.6914,0.4210) (3.6914,0.8998) (3.6914,0.3328) (3.6914,0.8196) (3.6914,0.5458) (3.6914,0.9151) (3.6914,0.2867) (3.6914,0.7550) (3.6957,0.8091) (3.6957,0.5709) (3.6957,0.9054) (3.6957,0.3167) (3.6957,0.7997) (3.6957,0.5920) (3.6957,0.8927) (3.6957,0.3541) (3.6957,0.8452) (3.6957,0.4834) (3.6957,0.9229) (3.6957,0.2629) (3.6957,0.7162) (3.7000,0.3075) (3.7000,0.7879) (3.7000,0.6183) (3.7000,0.8733) (3.7000,0.4095) (3.7000,0.8947) (3.7000,0.3486) (3.7000,0.8402) (3.7000,0.4968) (3.7000,0.9250) (3.7000,0.2568) (3.7000,0.7062) (3.7000,0.7677) (3.7043,0.8096) (3.7043,0.5711) (3.7043,0.9074) (3.7043,0.3113) (3.7043,0.7942) (3.7043,0.6054) (3.7043,0.8849) (3.7043,0.3772) (3.7043,0.8702) (3.7043,0.4183) (3.7043,0.9014) (3.7043,0.3293) (3.7043,0.8182) (3.7086,0.7737) (3.7086,0.6492) (3.7086,0.8446) (3.7086,0.4869) (3.7086,0.9265) (3.7086,0.2525) (3.7086,0.7000) (3.7086,0.7788) (3.7086,0.6390) (3.7086,0.8555) (3.7086,0.4584) (3.7086,0.9207) (3.7086,0.2707) (3.7129,0.7748) (3.7129,0.6478) (3.7129,0.8471) (3.7129,0.4809) (3.7129,0.9269) (3.7129,0.2517) (3.7129,0.6993) (3.7129,0.7808) (3.7129,0.6355) (3.7129,0.8600) (3.7129,0.4470) (3.7129,0.9178) (3.7129,0.2802) (3.7171,0.3447) (3.7171,0.8396) (3.7171,0.5005) (3.7171,0.9293) (3.7171,0.2443) (3.7171,0.6862) (3.7171,0.8004) (3.7171,0.5938) (3.7171,0.8966) (3.7214,0.8058) (3.7214,0.5824) (3.7214,0.9051) (3.7214,0.3197) (3.7214,0.8094) (3.7214,0.5742) (3.7214,0.9099) (3.7214,0.3052) (3.7214,0.7892) (3.7214,0.6192) (3.7214,0.8775) (3.7214,0.4001) (3.7214,0.8932) (3.7257,0.9092) (3.7257,0.3076) (3.7257,0.7935) (3.7257,0.6104) (3.7257,0.8860) (3.7257,0.3763) (3.7257,0.8745) (3.7257,0.4090) (3.7257,0.9006) (3.7257,0.3335) (3.7257,0.8282) (3.7257,0.5301) (3.7257,0.9280) (3.7300,0.3038) (3.7300,0.7890) (3.7300,0.6211) (3.7300,0.8778) (3.7300,0.4000) (3.7300,0.8952) (3.7300,0.3499) (3.7300,0.8484) (3.7300,0.4797) (3.7300,0.9310) (3.7300,0.2398) (3.7300,0.6799) (3.7300,0.8118) (3.7343,0.6318) (3.7343,0.8687) (3.7343,0.4260) (3.7343,0.9131) (3.7343,0.2963) (3.7343,0.7786) (3.7343,0.6437) (3.7343,0.8564) (3.7343,0.4591) (3.7343,0.9273) (3.7343,0.2516) (3.7343,0.7032) (3.7343,0.8582) (3.7386,0.9346) (3.7386,0.2284) (3.7386,0.6589) (3.7386,0.8403) (3.7386,0.5018) (3.7429,0.9357) (3.7429,0.2252) (3.7429,0.6530) (3.7429,0.8481) (3.7429,0.4821) (3.7429,0.2290) (3.7429,0.6609) (3.7429,0.8388) (3.7429,0.5061) (3.7429,0.4992) (3.7471,0.6868) (3.7471,0.8060) (3.7471,0.5860) (3.7471,0.9091) (3.7471,0.3097) (3.7471,0.8011) (3.7471,0.5971) (3.7471,0.9014) (3.7471,0.3330) (3.7471,0.8322) (3.7471,0.5232) (3.7471,0.9348) (3.7471,0.2285) (3.7514,0.6437) (3.7514,0.8604) (3.7514,0.4507) (3.7514,0.9287) (3.7514,0.2483) (3.7514,0.7002) (3.7514,0.7875) (3.7514,0.6277) (3.7514,0.8767) (3.7514,0.4055) (3.7514,0.9044) (3.7514,0.3244) (3.7514,0.8222) (3.7557,0.2593) (3.7557,0.7213) (3.7557,0.7550) (3.7557,0.6947) (3.7557,0.7965) (3.7557,0.6087) (3.7557,0.8946) (3.7557,0.3542) (3.7557,0.8591) (3.7557,0.4548) (3.7557,0.9312) (3.7557,0.2405) (3.7557,0.6860) (3.7600,0.9078) (3.7600,0.3146) (3.7600,0.8108) (3.7600,0.5769) (3.7600,0.9178) (3.7600,0.2837) (3.7600,0.7641) (3.7600,0.6777) (3.7600,0.8212) (3.7600,0.5520) (3.7600,0.9298) (3.7600,0.2453) (3.7600,0.6961) (3.7643,0.7024) (3.7643,0.7868) (3.7643,0.6314) (3.7643,0.8761) (3.7643,0.4086) (3.7643,0.9096) (3.7643,0.3095) (3.7643,0.8044) (3.7643,0.5923) (3.7643,0.3113) (3.7643,0.8070) (3.7643,0.5862) (3.7643,0.9131) (3.7686,0.6206) (3.7686,0.8874) (3.7686,0.3766) (3.7686,0.8848) (3.7686,0.3842) (3.7686,0.8916) (3.7686,0.3643) (3.7686,0.8727) (3.7686,0.4186) (3.7686,0.9172) (3.7686,0.2863) (3.7686,0.7700) (3.7686,0.6674) (3.7729,0.8784) (3.7729,0.4030) (3.7729,0.9077) (3.7729,0.3161) (3.7729,0.8156) (3.7729,0.5673) (3.7729,0.9261) (3.7729,0.2582) (3.7729,0.7225) (3.7729,0.7564) (3.7729,0.6952) (3.7729,0.7994) (3.7729,0.6050) (3.7771,0.8696) (3.7771,0.4282) (3.7771,0.9248) (3.7771,0.2626) (3.7771,0.7313) (3.7771,0.7421) (3.7771,0.7228) (3.7771,0.7567) (3.7771,0.6954) (3.7771,0.8001) (3.7771,0.6041) (3.7771,0.9033) (3.7771,0.3298) (3.7814,0.5514) (3.7814,0.9354) (3.7814,0.2286) (3.7814,0.6669) (3.7814,0.8401) (3.7814,0.5081) (3.7814,0.9451) (3.7814,0.1962) (3.7814,0.5963) (3.7814,0.9103) (3.7814,0.3087) (3.7814,0.8070) (3.7814,0.5889) (3.7857,0.5546) (3.7857,0.9352) (3.7857,0.2295) (3.7857,0.6695) (3.7857,0.8377) (3.7857,0.5148) (3.7857,0.9456) (3.7857,0.1947) (3.7857,0.5937) (3.7857,0.9132) (3.7857,0.3001) (3.7857,0.7951) (3.7857,0.6168) (3.7900,0.2279) (3.7900,0.6670) (3.7900,0.8418) (3.7900,0.5047) (3.7900,0.9474) (3.7900,0.1888) (3.7900,0.5805) (3.7900,0.9230) (3.7900,0.2695) (3.7900,0.7461) (3.7900,0.7179) (3.7900,0.7676) (3.7900,0.6761) (3.7943,0.9354) (3.7943,0.2292) (3.7943,0.6704) (3.7943,0.8384) (3.7943,0.5141) (3.7943,0.9478) (3.7943,0.1877) (3.7943,0.5784) (3.7943,0.9252) (3.7943,0.2625) (3.7943,0.7345) (3.7943,0.7399) (3.7943,0.7302) (3.7986,0.2144) (3.7986,0.6398) (3.7986,0.8754) (3.7986,0.4144) (3.7986,0.9218) (3.7986,0.2738) (3.7986,0.7554) (3.7986,0.7019) (3.7986,0.7947) (3.7986,0.6196) (3.7986,0.8953) (3.7986,0.3561) (3.7986,0.8710) (3.8029,0.7783) (3.8029,0.6561) (3.8029,0.8580) (3.8029,0.4632) (3.8029,0.9456) (3.8029,0.1958) (3.8029,0.5987) (3.8029,0.9137) (3.8029,0.2999) (3.8029,0.7985) (3.8029,0.6118) (3.8029,0.9032) (3.8029,0.3326) (3.8071,0.1798) (3.8071,0.5615) (3.8071,0.9374) (3.8071,0.2235) (3.8071,0.6607) (3.8071,0.8534) (3.8071,0.4763) (3.8071,0.9496) (3.8071,0.1821) (3.8071,0.5670) (3.8071,0.2323) (3.8071,0.6790) (3.8071,0.8297) (3.8114,0.6532) (3.8114,0.8634) (3.8114,0.4496) (3.8114,0.9432) (3.8114,0.2043) (3.8114,0.6196) (3.8114,0.8984) (3.8114,0.3480) (3.8114,0.4455) (3.8114,0.9415) (3.8114,0.2098) (3.8114,0.6319) (3.8114,0.8866) (3.8157,0.1989) (3.8157,0.6080) (3.8157,0.9095) (3.8157,0.3142) (3.8157,0.8222) (3.8157,0.5578) (3.8157,0.9412) (3.8157,0.2113) (3.8157,0.6358) (3.8157,0.8836) (3.8157,0.3925) (3.8157,0.3129) (3.8157,0.8204) (3.8200,0.2165) (3.8200,0.6480) (3.8200,0.8713) (3.8200,0.4282) (3.8200,0.9353) (3.8200,0.2311) (3.8200,0.6788) (3.8200,0.8329) (3.8200,0.5316) (3.8200,0.9512) (3.8200,0.1773) (3.8200,0.5573) (3.8200,0.9425) (3.8243,0.2499) (3.8243,0.7169) (3.8243,0.7761) (3.8243,0.6645) (3.8243,0.8526) (3.8243,0.4806) (3.8243,0.9546) (3.8243,0.1656) (3.8243,0.5285) (3.8243,0.9530) (3.8243,0.1714) (3.8243,0.5432) (3.8243,0.9489) (3.8286,0.9566) (3.8286,0.1588) (3.8286,0.5114) (3.8329,0.9582) (3.8329,0.1535) (3.8329,0.4981) (3.8371,0.9590) (3.8371,0.1508) (3.8371,0.4915) (3.8414,0.9596) (3.8414,0.1488) (3.8414,0.4864) (3.8457,0.9614) (3.8457,0.1427) (3.8457,0.4704) (3.8457,0.1545) (3.8457,0.5024) (3.8500,0.9625) (3.8500,0.1390) (3.8500,0.4607) (3.8500,0.9565) (3.8500,0.1600) (3.8500,0.5176) (3.8500,0.9613) (3.8500,0.1432) (3.8500,0.4723) (3.8500,0.1494) (3.8500,0.4893) (3.8500,0.4650) (3.8500,0.1557) (3.8543,0.9595) (3.8543,0.1497) (3.8543,0.4906) (3.8543,0.9632) (3.8543,0.1365) (3.8543,0.4543) (3.8543,0.9555) (3.8543,0.1638) (3.8543,0.5279) (3.8543,0.1460) (3.8543,0.4805) (3.8543,0.1405) (3.8543,0.4655) (3.8586,0.1728) (3.8586,0.5516) (3.8586,0.9544) (3.8586,0.1680) (3.8586,0.5394) (3.8586,0.9586) (3.8586,0.1530) (3.8586,0.5000) (3.8586,0.9646) (3.8586,0.1316) (3.8586,0.4410) (3.8586,0.9512) (3.8586,0.1791) (3.8629,0.1339) (3.8629,0.4480) (3.8629,0.9553) (3.8629,0.1651) (3.8629,0.5325) (3.8629,0.9616) (3.8629,0.1425) (3.8629,0.4721) (3.8629,0.9627) (3.8629,0.1387) (3.8629,0.4615) (3.8629,0.1484) (3.8629,0.4881) (3.8671,0.9610) (3.8671,0.1451) (3.8671,0.4797) (3.8671,0.9652) (3.8671,0.1299) (3.8671,0.4371) (3.8671,0.9515) (3.8671,0.1785) (3.8671,0.5671) (3.8671,0.9494) (3.8671,0.1859) (3.8671,0.5852) (3.8671,0.9387) (3.8714,0.9510) (3.8714,0.1805) (3.8714,0.5726) (3.8714,0.9475) (3.8714,0.1927) (3.8714,0.6022) (3.8714,0.9274) (3.8714,0.2606) (3.8714,0.7459) (3.8714,0.7337) (3.8714,0.7563) (3.8714,0.7135) (3.8714,0.7914) (3.8757,0.5213) (3.8757,0.9672) (3.8757,0.1230) (3.8757,0.4182) (3.8757,0.9430) (3.8757,0.2083) (3.8757,0.6393) (3.8757,0.8938) (3.8757,0.3680) (3.8757,0.9014) (3.8757,0.3446) (3.8757,0.8753) (3.8757,0.4230) (3.8800,0.8612) (3.8800,0.4638) (3.8800,0.9649) (3.8800,0.1313) (3.8800,0.4427) (3.8800,0.9573) (3.8800,0.1587) (3.8800,0.5182) (3.8800,0.9687) (3.8800,0.1176) (3.8800,0.4025) (3.8800,0.9331) (3.8800,0.2421) (3.8843,0.3071) (3.8843,0.8266) (3.8843,0.5568) (3.8843,0.9586) (3.8843,0.1543) (3.8843,0.5069) (3.8843,0.9709) (3.8843,0.1098) (3.8843,0.3797) (3.8843,0.9148) (3.8843,0.3027) (3.8843,0.8198) (3.8843,0.5738) (3.8886,0.1504) (3.8886,0.4968) (3.8886,0.9721) (3.8886,0.1055) (3.8886,0.3668) (3.8886,0.9032) (3.8886,0.3400) (3.8886,0.8726) (3.8886,0.4322) (3.8886,0.9543) (3.8886,0.1697) (3.8886,0.5480) (3.8886,0.9632) (3.8929,0.2782) (3.8929,0.7816) (3.8929,0.6645) (3.8929,0.8679) (3.8929,0.4463) (3.8929,0.9620) (3.8929,0.1424) (3.8929,0.4754) (3.8929,0.9709) (3.8929,0.1101) (3.8929,0.3815) (3.8929,0.9186) (3.8929,0.2912) (3.8971,0.5465) (3.8971,0.9659) (3.8971,0.1285) (3.8971,0.4365) (3.8971,0.9586) (3.8971,0.1547) (3.8971,0.5097) (3.8971,0.9739) (3.8971,0.0990) (3.8971,0.3476) (3.8971,0.8838) (3.8971,0.4004) (3.8971,0.9356) (3.9014,0.4237) (3.9014,0.9526) (3.9014,0.1761) (3.9014,0.5661) (3.9014,0.9583) (3.9014,0.1558) (3.9014,0.5132) (3.9014,0.9747) (3.9014,0.0963) (3.9014,0.3395) (3.9014,0.8749) (3.9014,0.4271) (3.9014,0.9546) (3.9057,0.9764) (3.9057,0.0899) (3.9057,0.3195) (3.9057,0.8492) (3.9057,0.5001) (3.9100,0.8779) (3.9100,0.4191) (3.9100,0.9519) (3.9100,0.1791) (3.9100,0.5748) (3.9100,0.9556) (3.9100,0.1658) (3.9100,0.5407) (3.9100,0.9710) (3.9100,0.1101) (3.9100,0.3830) (3.9100,0.9239) (3.9100,0.2748) (3.9143,0.0841) (3.9143,0.3014) (3.9143,0.8242) (3.9143,0.5672) (3.9143,0.9609) (3.9143,0.1471) (3.9143,0.4911) (3.9143,0.9783) (3.9143,0.8199) (3.9143,0.5780) (3.9143,0.9548) (3.9143,0.1690) (3.9143,0.5497) (3.9186,0.9592) (3.9186,0.1532) (3.9186,0.5084) (3.9186,0.9794) (3.9186,0.0792) (3.9186,0.2857) (3.9186,0.7997) (3.9186,0.6276) (3.9186,0.9158) (3.9186,0.3021) (3.9186,0.8262) (3.9186,0.5626) (3.9186,0.9643) (3.9229,0.2962) (3.9229,0.8178) (3.9229,0.5844) (3.9229,0.9527) (3.9229,0.1766) (3.9229,0.5705) (3.9229,0.9612) (3.9229,0.1462) (3.9229,0.4897) (3.9229,0.9803) (3.9229,0.0758) (3.9229,0.2747) (3.9229,0.7816) (3.9271,0.9302) (3.9271,0.2551) (3.9271,0.7462) (3.9271,0.7437) (3.9271,0.7392) (3.9271,0.7571) (3.9271,0.7222) (3.9271,0.7880) (3.9271,0.6561) (3.9271,0.8860) (3.9271,0.3965) (3.9271,0.9397) (3.9271,0.2224) (3.9314,0.0742) (3.9314,0.2702) (3.9314,0.7752) (3.9314,0.6851) (3.9314,0.8481) (3.9314,0.5063) (3.9314,0.9827) (3.9314,0.0668) (3.9314,0.2452) (3.9314,0.7277) (3.9314,0.7791) (3.9314,0.6766) (3.9314,0.8602) (3.9357,0.6512) (3.9357,0.8940) (3.9357,0.3730) (3.9357,0.9204) (3.9357,0.2883) (3.9357,0.8075) (3.9357,0.6117) (3.9357,0.9348) (3.9357,0.2399) (3.9357,0.7176) (3.9357,0.7976) (3.9357,0.6354) (3.9357,0.9117) (3.9400,0.0606) (3.9400,0.2244) (3.9400,0.6858) (3.9400,0.8490) (3.9400,0.5051) (3.9400,0.9849) (3.9400,0.2174) (3.9400,0.6702) (3.9400,0.8708) (3.9400,0.4433) (3.9400,0.9723) (3.9400,0.1060) (3.9400,0.3735) (3.9443,0.8483) (3.9443,0.5076) (3.9443,0.9858) (3.9443,0.0550) (3.9443,0.2051) (3.9443,0.6431) (3.9443,0.9053) (3.9443,0.3382) (3.9443,0.8828) (3.9443,0.4080) (3.9443,0.9527) (3.9443,0.1777) (3.9443,0.5763) (3.9486,0.9853) (3.9486,0.0572) (3.9486,0.2128) (3.9486,0.6614) (3.9486,0.8842) (3.9486,0.4041) (3.9486,0.9509) (3.9486,0.1845) (3.9486,0.5941) (3.9486,0.1797) (3.9486,0.5820) (3.9486,0.9606) (3.9486,0.1494) (3.9529,0.0863) (3.9529,0.3116) (3.9529,0.8479) (3.9529,0.5098) (3.9529,0.9878) (3.9529,0.0475) (3.9529,0.1788) (3.9529,0.5804) (3.9529,0.9626) (3.9529,0.1421) (3.9529,0.4820) (3.9529,0.0510) (3.9529,0.1912) (3.9571,0.8310) (3.9571,0.5556) (3.9571,0.9770) (3.9571,0.0888) (3.9571,0.3202) (3.9571,0.8614) (3.9571,0.4726) (3.9571,0.9863) (3.9571,0.0534) (3.9571,0.2001) (3.9571,0.6334) (3.9571,0.9188) (3.9571,0.2951) (3.9614,0.9904) (3.9614,0.0378) (3.9614,0.1442) (3.9614,0.4889) (3.9614,0.9899) (3.9614,0.0397) (3.9614,0.1512) (3.9614,0.5083) (3.9614,0.1481) (3.9614,0.4998) (3.9657,0.1377) (3.9657,0.4708) (3.9657,0.9880) (3.9657,0.0469) (3.9657,0.1771) (3.9657,0.5779) (3.9657,0.9673) (3.9657,0.1253) (3.9657,0.4346) (3.9657,0.9745) (3.9657,0.0986) (3.9657,0.3525) (3.9657,0.9051) (3.9700,0.1533) (3.9700,0.5153) (3.9700,0.9916) (3.9700,0.0332) (3.9700,0.1274) (3.9700,0.4412) (3.9700,0.9788) (3.9700,0.0824) (3.9700,0.3003) (3.9700,0.8341) (3.9700,0.5493) (3.9700,0.9829) (3.9700,0.0669) (3.9743,0.6351) (3.9743,0.9211) (3.9743,0.2889) (3.9743,0.8165) (3.9743,0.5955) (3.9743,0.9573) (3.9743,0.1624) (3.9743,0.5406) (3.9743,0.9870) (3.9743,0.0510) (3.9743,0.1922) (3.9743,0.6171) (3.9743,0.9391) (3.9786,0.4651) (3.9786,0.9898) (3.9786,0.0402) (3.9786,0.1536) (3.9786,0.5172) (3.9786,0.9935) (3.9786,0.0258) (3.9786,0.1001) (3.9786,0.3584) (3.9786,0.9149) (3.9786,0.3097) (3.9786,0.8505) (3.9786,0.5058) (3.9829,0.9437) (3.9829,0.2117) (3.9829,0.6647) (3.9829,0.8876) (3.9829,0.3973) (3.9829,0.9537) (3.9829,0.1759) (3.9829,0.5775) (3.9829,0.9718) (3.9829,0.1091) (3.9829,0.3871) (3.9829,0.2071) (3.9829,0.6541) (3.9871,0.2135) (3.9871,0.6695) (3.9871,0.8822) (3.9871,0.4143) (3.9871,0.9675) (3.9871,0.1253) (3.9871,0.4369) (3.9871,0.9809) (3.9871,0.0747) (3.9871,0.2755) (3.9871,0.7958) (3.9871,0.6480) (3.9871,0.9094) (3.9914,0.1376) (3.9914,0.4735) (3.9914,0.9951) (3.9914,0.0196) (3.9914,0.0768) (3.9914,0.2829) (3.9914,0.8098) (3.9914,0.6148) (3.9914,0.9453) (3.9914,0.2065) (3.9914,0.6540) (3.9914,0.9032) (3.9914,0.3488) (3.9957,0.2217) (3.9957,0.6894) (3.9957,0.8556) (3.9957,0.4937) (3.9957,0.9988) (3.9957,0.0049) (3.9957,0.0195) (3.9957,0.0764) (3.9957,0.2820) (3.9957,0.8090) (3.9957,0.6175) (3.9957,0.9438) (3.9957,0.2120) (4.0000,0.0000)};
\end{axis}
\end{tikzpicture}}{The bifurcation diagram of the logistic map: for $281$ values of $r$ between $2.8$ and $4$, the iterates $x_n$, $6001\le n\le6064$, from $x_0=1/2$, with nearby points merged and at most $13$ points per value of $r$ ($1611$ points in all). A single point for $r<3$ is the stable fixed point; it splits into the cycle of period $2$ at $r=3$, of period $4$ at $r=3.4495$, and so on (\cref{tab:period-doubling}); beyond $3.57$ the iterates spread over intervals, apart from windows such as the cycle of period $3$ near $r=3.83$.\label{fig:logistic-bifurcation}}
\end{figure}

\subsection{Chaos at \texorpdfstring{$r=4$}{r = 4}}

For $r=4$ the logistic map can be solved explicitly, and the solution shows what lies behind the irregular orbits.

\begin{proposition}[The logistic map at $r=4$]
\label{prop:logistic-four}
Let $x_0=\sin^2\theta$ with a real number $\theta$. The iterates of $f_4(x)=4x(1-x)$ are $x_n=\sin^2(2^n\theta)$. Consequently:
\begin{enumerate}
\item for every $m\ge1$, the points $\sin^2(j\pi/(2^m-1))$, $j\in\mathbb Z$, are periodic with a period dividing $m$, and the periodic points are dense in $[0,1]$;
\item if $x_0=\sin^2\theta$ and $\tilde x_0=\sin^2\tilde\theta$, the iterates $x_n$ and $\tilde x_n$ are the values of $\sin^2$ at angles whose difference $2^n(\theta-\tilde\theta)$ doubles in every step.
\end{enumerate}
\end{proposition}

\begin{proof}
$f_4(\sin^2\varphi)=4\sin^2\varphi\cos^2\varphi=\sin^22\varphi$ by the double angle formula (\vcref{1}{thm:addition}), and induction gives the formula. (i) For $\theta=j\pi/(2^m-1)$, $2^m\theta=\theta+j\pi$, and $\sin^2$ has the period $\pi$, so $x_m=x_0$. Given $x\in[0,1]$ and $\varepsilon>0$, write $x=\sin^2\varphi$; since $\sin^2$ is uniformly continuous and the points $j\pi/(2^m-1)$ have the spacing $\pi/(2^m-1)$, some such point gives a value within $\varepsilon$ of $x$ when $m$ is large. (ii) follows from the formula.
\end{proof}

Two initial values that differ by $\eta$ in the angle lead to angles that differ by $2^n\eta$, so the iterates become unrelated after about $\log_2(1/\eta)$ steps, although each of them is determined by its initial value: \term{sensitive dependence on initial conditions}. In binary64 arithmetic, $\eta$ is of the order of the unit roundoff, and after some fifty steps a computed orbit has nothing to do with the exact orbit of the same initial value. We iterated $f_4$ from the binary64 number nearest to $0.1$ in binary64 arithmetic and, for comparison, in the ball arithmetic of python-flint with $2000$ bits. The difference was $2.0\times10^{-14}$ after $10$ steps, $2.2\times10^{-11}$ after $20$, $2.8\times10^{-8}$ after $30$ and $1.8\times10^{-5}$ after $40$, growing roughly by the factor $2^{10}=1024$ every ten steps, and it first exceeded $0.1$ at step $53$. The computed orbit is still an orbit of the map, to within rounding, but not the one that was asked for. Dense periodic orbits, sensitive dependence and orbits that come close to every point are the ingredients of a common definition of \term{chaos} for maps~\cite[Sect.~15.4]{hirsch2013}.

\begin{note}{Period three implies chaos}
For continuous maps of an interval, periodic points obey remarkable laws. Li and Yorke proved in 1975 that if a continuous $f\colon I\to I$ has a point of period $3$, then it has points of every period, and an uncountable set of points whose orbits do not approach any periodic orbit~\cite{li1975}; the title of their paper is ``Period three implies chaos''. Sharkovskii had proved in 1964 a sharper statement: order the positive integers as $3\succ5\succ7\succ\cdots\succ2\cdot3\succ2\cdot5\succ\cdots\succ4\cdot3\succ\cdots\succ2^3\succ2^2\succ2\succ1$, odd numbers first, then their doubles, their quadruples, and so on, and finally the powers of $2$ in decreasing order; if $f$ has a point of period $m$, it has points of every period that follows $m$ in this order~\cite{sharkovskii1964}. The logistic map has a cycle of period $3$ in the window near $r=3.83$, so for such $r$ it has cycles of every period, almost all of them unstable. Strogatz~\cite[Chap.~10]{strogatz2015} introduces the logistic map, its periodic windows and Feigenbaum's universality~\cite{feigenbaum1978}.
\end{note}

The contrast with the differential equation is complete. The logistic equation $y'=y(1-y)$ has the explicit solutions of \eqref{eq:logistic-solution}, which approach the equilibrium $1$ monotonically from every positive initial value, as every solution of an autonomous equation of the first order is monotone (\cref{prop:autonomous-monotone}). Euler's method for it is the logistic map with $r=1+h$, so it reproduces this behaviour for $h<1$, approaches the equilibrium with alternation for $1<h<2$, and for larger steps produces cycles of period $2$, $4$, $8$, \ldots, and from $h=2.57$ on the irregular orbits of \cref{fig:logistic-bifurcation}, none of which has anything to do with the differential equation. A scalar autonomous differential equation of the first order cannot be chaotic, and in the plane the Poincar\'e--Bendixson theorem (\cref{thm:poincare-bendixson}) excludes chaos as well; a one-dimensional recurrence can be.

\section{Computing with Recurrences: A First Look}
\label{sec:recurrence-computation}

A recurrence is an algorithm, and this section runs the recurrences of the chapter in floating-point arithmetic and compares the results with exact values. All computations were done in binary64 arithmetic with plain Python, on the machine and with the version of the stardust library described in \cref{sec:systems-practice}; exact reference values were computed in rational arithmetic or with the ball arithmetic of python-flint (version 0.9.0), with error bounds far below the digits shown. Sums were formed from the first term to the last. A relative error is $|\tilde y-y|/|y|$, where $\tilde y$ is the computed and $y$ the exact value.

\subsection{Fibonacci Numbers in Three Ways}

\Vcref{1}{sec:recurrences} found that Binet's formula, evaluated as $\phi^n/\sqrt5$ in binary64 and rounded to the nearest integer, first gives a wrong Fibonacci number at $n=71$. \Cref{tab:fibonacci} shows the reason: the relative error of the computed $\phi^n/\sqrt5$ stays between $10^{-15}$ and $4\times10^{-15}$ for $n\ge50$, but since $F_n$ grows, the absolute error passes $1/2$ between $n=70$ and $n=71$. The recurrence $F_{n+1}=F_n+F_{n-1}$ in binary64 adds integers, and the sum of two integers is exact as long as it is representable, that is, below $2^{53}$ or a multiple of a suitable power of $2$; the computed values were exact up to $F_{78}=8\,944\,394\,323\,791\,464<2^{53}$, and the first inexact one was $F_{79}$, with the relative error $6.9\times10^{-17}$. These are the bounds that \code{linalg.matrix\_power} met in \cref{sec:computing-matrix-functions}, where the power $F^k$ of the Fibonacci matrix contains $F_{k+1}$ and was exact up to $k=77$. Computed forward, the Fibonacci numbers are a dominant solution, and their relative error did not grow: at $n=100$ it was $2.3\times10^{-16}$.

\begin{table}
\sdfloat{%
\begin{sdtable}{colspec={rrlll}, row{1}={font=\small\bfseries}}
\toprule
$n$ & $F_n$ & $|\phi^n/\sqrt5-F_n|$ & relative & recurrence \\
\midrule
$30$ & $832040$ & $2.4\times10^{-7}$ & $2.9\times10^{-13}$ & $0$ \\
$50$ & $12586269025$ & $1.9\times10^{-5}$ & $1.5\times10^{-15}$ & $0$ \\
$70$ & $190392490709135$ & $4.4\times10^{-1}$ & $2.3\times10^{-15}$ & $0$ \\
$71$ & $308061521170129$ & $6.9\times10^{-1}$ & $2.2\times10^{-15}$ & $0$ \\
$78$ & $8944394323791464$ & $2.4\times10^{1}$ & $2.7\times10^{-15}$ & $0$ \\
$79$ & $14472334024676221$ & $3.9\times10^{1}$ & $2.7\times10^{-15}$ & $6.9\times10^{-17}$ \\
$100$ & $354224848179261915075$ & $1.2\times10^{6}$ & $3.4\times10^{-15}$ & $2.3\times10^{-16}$ \\
\bottomrule
\end{sdtable}}{Fibonacci numbers in binary64: the distance of the computed $\phi^n/\sqrt5$ from $F_n$, absolute and relative, and the relative error of the forward recurrence. Rounding $\phi^n/\sqrt5$ gives $F_n$ up to $n=70$ and fails from $n=71$ on. For $n=30$ the distance consists mainly of the term $\hat\phi^{30}/\sqrt5$ that the formula omits.\label{tab:fibonacci}}
\end{table}

Run backwards, $F_{n-1}=F_{n+1}-F_n$, the same recurrence computes a minimal solution, since $\hat\phi^n$, which was minimal going forward, becomes dominant going backward, and the result depends on whether the starting values are exact. From the binary64 values of $F_{101}$ and $F_{100}$, which carry rounding errors of relative size about $10^{-16}$, the backward recurrence produced $F_{50}\approx5.2\times10^{14}$ instead of $1.26\times10^{10}$ and $F_0\approx1.5\times10^{25}$ instead of $0$: the initial error, about $10^4$ in absolute terms, was amplified by roughly $\phi^{100}=7.9\times10^{20}$, as \cref{thm:recurrence-direction} predicts. From the exactly representable $F_{71}$ and $F_{70}$, the backward recurrence returned exactly $F_1=1$ and $F_0=0$, since every intermediate result was an integer computed without rounding. The instability lies in the amplification of errors, not in the recurrence itself, and with no error to amplify there is nothing to see.

\subsection{The Integrals \texorpdfstring{$I_n$}{In}}

\Cref{tab:integral-recurrence} shows the integrals $I_n=\int_0^1x^ne^{x-1}\,dx$ of \cref{sec:three-term}, computed forward from the binary64 value of $1-e^{-1}$ and backward from $\tilde I_{30}\coloneqq0$. The initial value has the error $-1.24\times10^{-17}$, and the forward errors agree to all digits shown with $(-1)^nn!$ times this error, as \cref{prop:integral-two-directions}\,(i) predicts: the rounding errors of the later steps are negligible beside the amplified initial error. From $n=18$ on, the computed values are wrong in sign or in magnitude. The backward recurrence has relative errors below $1.5\times10^{-16}$ up to $n=18$; its errors grow towards the starting index, $6.2\times10^{-15}$ at $n=20$ and $4.9\times10^{-8}$ at $n=25$, in accordance with the bound $n!/31!$ of \cref{prop:integral-two-directions}\,(ii), which is $3.0\times10^{-16}$ and $1.9\times10^{-9}$ in absolute terms.

\begin{table}
\sdfloat{%
\begin{sdtable}{colspec={rllll}, row{1}={font=\small\bfseries}}
\toprule
$n$ & $I_n$ & forward & error, forward & error, backward \\
\midrule
$0$ & $0.6321206$ & $0.6321206$ & $2.0\times10^{-17}$ & $2.0\times10^{-17}$ \\
$5$ & $0.1455329$ & $0.1455329$ & $1.0\times10^{-14}$ & $1.4\times10^{-16}$ \\
$10$ & $0.0838771$ & $0.0838771$ & $5.4\times10^{-10}$ & $3.5\times10^{-17}$ \\
$15$ & $0.0590175$ & $0.0590338$ & $2.8\times10^{-4}$ & $1.9\times10^{-17}$ \\
$16$ & $0.0557193$ & $0.0554593$ & $4.7\times10^{-3}$ & $5.6\times10^{-17}$ \\
$17$ & $0.0527711$ & $0.0571919$ & $8.4\times10^{-2}$ & $8.3\times10^{-17}$ \\
$18$ & $0.0501199$ & $-0.0294537$ & $1.6$ & $5.8\times10^{-17}$ \\
$20$ & $0.0455449$ & $-3.02\times10^{1}$ & $6.6\times10^{2}$ & $6.2\times10^{-15}$ \\
$25$ & $0.0370862$ & $1.93\times10^{8}$ & $5.2\times10^{9}$ & $4.9\times10^{-8}$ \\
\bottomrule
\end{sdtable}}{The integrals $I_n=\int_0^1x^ne^{x-1}\,dx$ computed in binary64 by $I_n=1-nI_{n-1}$ forward from $I_0=1-e^{-1}$, and by $I_{n-1}=(1-I_n)/n$ backward from $\tilde I_{30}=0$: the exact values, the forward values, and the relative errors of both (\cref{prop:integral-two-directions}).\label{tab:integral-recurrence}}
\end{table}

\subsection{Bessel Functions of Integral Order}

\Cref{tab:bessel-recurrence} compares three ways of computing $J_n(1)$ and $J_n(10)$: the forward recurrence from the correctly rounded values of $J_0(x)$ and $J_1(x)$, Miller's backward recurrence (\cref{alg:miller}) with $N=60$ and the normalisation \eqref{eq:bessel-normalisation}, and the library function \code{special.jv}. The forward recurrence loses accuracy as soon as $n$ exceeds $x$, at a rate that matches the growth of $|Y_n(x)/J_n(x)|$ in \cref{fig:bessel-minimal}: at $x=1$ the relative error is $2.7\times10^{-4}$ at $n=8$ and $27$ at $n=10$, at $x=10$ it is $5.3\times10^{-8}$ at $n=20$ and $0.10$ at $n=25$. Miller's algorithm, run as stated in \cref{alg:miller} with the start $y_N=1$, gave relative errors of at most $3.3\times10^{-16}$ at $x=1$ and $1.9\times10^{-15}$ at $x=10$ for all $n\le40$; at $x=10$ the largest occurred at $n=6$, where $J_6(10)=-0.0145$ is small compared with its neighbours $J_5(10)=-0.234$ and $J_7(10)=0.217$. Starting lower gave visibly worse results near the top: with $N=40$, the error at $x=10$ reached $2.6\times10^{-4}$ at $n=39$, and with $N=50$ the largest error for $n\le40$ was $3.3\times10^{-15}$ again, in accordance with the error term $f_{N+1}/g_{N+1}$ in the proof of \cref{thm:miller}. The library function had relative errors of at most $3.2\times10^{-16}$ at $x=1$ and $1.2\times10^{-15}$ at $x=10$ for $n\le40$.

\begin{table}
\sdfloat{%
\begin{sdtable}{colspec={rllll}, row{1}={font=\small\bfseries}}
\toprule
$n$ & $J_n(x)$ & error, forward & error, Miller & error, \code{jv} \\
\midrule
\SetCell[c=5]{l} $x=1$ & & & & \\
$1$ & $0.4400506$ & $4.1\times10^{-17}$ & $1.7\times10^{-16}$ & $8.5\times10^{-17}$ \\
$2$ & $0.1149035$ & $7.8\times10^{-16}$ & $1.8\times10^{-16}$ & $6.0\times10^{-17}$ \\
$5$ & $2.497577\times10^{-4}$ & $6.2\times10^{-11}$ & $4.6\times10^{-17}$ & $2.6\times10^{-16}$ \\
$8$ & $9.422344\times10^{-8}$ & $2.7\times10^{-4}$ & $6.6\times10^{-19}$ & $6.6\times10^{-19}$ \\
$10$ & $2.630615\times10^{-10}$ & $2.7\times10^{1}$ & $2.8\times10^{-16}$ & $1.1\times10^{-16}$ \\
$15$ & $2.297532\times10^{-17}$ & $2.4\times10^{15}$ & $8.5\times10^{-17}$ & $5.0\times10^{-17}$ \\
$20$ & $3.873503\times10^{-25}$ & $6.3\times10^{30}$ & $8.5\times10^{-17}$ & $3.4\times10^{-17}$ \\
$30$ & $3.482870\times10^{-42}$ & $5.2\times10^{64}$ & $2.2\times10^{-16}$ & $1.4\times10^{-16}$ \\
$40$ & $1.107916\times10^{-60}$ & $3.8\times10^{101}$ & $2.4\times10^{-16}$ & $1.4\times10^{-16}$ \\
\midrule
\SetCell[c=5]{l} $x=10$ & & & & \\
$5$ & $-0.2340615$ & $6.1\times10^{-17}$ & $1.8\times10^{-16}$ & $6.1\times10^{-17}$ \\
$10$ & $0.2074861$ & $5.9\times10^{-16}$ & $5.4\times10^{-17}$ & $1.9\times10^{-16}$ \\
$15$ & $4.507973\times10^{-3}$ & $5.3\times10^{-13}$ & $4.0\times10^{-16}$ & $1.7\times10^{-17}$ \\
$20$ & $1.151337\times10^{-5}$ & $5.3\times10^{-8}$ & $1.5\times10^{-16}$ & $1.5\times10^{-16}$ \\
$25$ & $7.214635\times10^{-9}$ & $1.0\times10^{-1}$ & $4.6\times10^{-16}$ & $1.2\times10^{-16}$ \\
$30$ & $1.551096\times10^{-12}$ & $1.8\times10^{6}$ & $2.5\times10^{-16}$ & $5.1\times10^{-16}$ \\
$40$ & $6.030895\times10^{-21}$ & $8.6\times10^{22}$ & $1.1\times10^{-16}$ & $3.6\times10^{-16}$ \\
\bottomrule
\end{sdtable}}{The Bessel functions $J_n(1)$ and $J_n(10)$ in binary64: the relative errors of the forward recurrence from the correctly rounded $J_0(x)$ and $J_1(x)$, of Miller's backward recurrence from $N=60$ with the normalisation \eqref{eq:bessel-normalisation}, and of the library function \code{special.jv}.\label{tab:bessel-recurrence}}
\end{table}

\subsection{Multistep Methods}

\Cref{tab:leapfrog} shows the leapfrog method \eqref{eq:leapfrog} for $y'=-y$, $y(0)=1$, with $h=0.1$ and the exact starting value $y_1=e^{-0.1}$, for which \cref{prop:leapfrog-roots} gives $B=7.47\times10^{-5}$. The error at $t=1$ is $7.9\times10^{-4}$, of which the parasitic component $B\mu_2^{10}$ contributes $2.0\times10^{-4}$; the parasitic component then grows like $e^{t}$ while the solution decays like $e^{-t}$, it overtakes the solution before $t=5$, and from $t=10$ on the computed values are the parasitic solution, to all digits shown. Starting with an Euler step, $y_1=1-h$, gave $B=2.5\times10^{-3}$, and the same behaviour with errors larger by the factor $33$. With $h=0.05$ and $h=0.025$ the relative error at $t=10$ was $4.8\times10^3$ and $6.2\times10^2$: smaller steps postpone the takeover, as convergence on bounded intervals requires, but do not prevent it.

\begin{table}
\sdfloat{%
\begin{sdtable}{colspec={rllll}, row{1}={font=\small\bfseries}}
\toprule
$t$ & $e^{-t}$ & $y_n$ & $y_n-e^{-t}$ & $B\mu_2^n$ \\
\midrule
$1$ & $3.6788\times10^{-1}$ & $3.6867\times10^{-1}$ & $7.86\times10^{-4}$ & $2.03\times10^{-4}$ \\
$2$ & $1.3534\times10^{-1}$ & $1.3633\times10^{-1}$ & $9.90\times10^{-4}$ & $5.50\times10^{-4}$ \\
$5$ & $6.7379\times10^{-3}$ & $1.7788\times10^{-2}$ & $1.11\times10^{-2}$ & $1.10\times10^{-2}$ \\
$10$ & $4.5400\times10^{-5}$ & $1.6183$ & $1.62$ & $1.62$ \\
$15$ & $3.0590\times10^{-7}$ & $2.3819\times10^{2}$ & $2.38\times10^{2}$ & $2.38\times10^{2}$ \\
$20$ & $2.0612\times10^{-9}$ & $3.5059\times10^{4}$ & $3.51\times10^{4}$ & $3.51\times10^{4}$ \\
\bottomrule
\end{sdtable}}{The leapfrog method for $y'=-y$, $y(0)=1$, with $h=0.1$ and $y_1=e^{-0.1}$: the exact solution, the computed $y_n$ at $t=nh$, its error, and the parasitic component $B\mu_2^n$ of \cref{prop:leapfrog-roots}.\label{tab:leapfrog}}
\end{table}

\Cref{tab:unstable-multistep} shows the zero-unstable two-step method of \cref{sec:ode-to-difference} for the same equation, with exact starting values, beside the zero-stable second-order Adams--Bashforth method $y_{n+2}=y_{n+1}+h\bigl(3f(t_{n+1},y_{n+1})-f(t_n,y_n)\bigr)/2$, whose $\rho(\zeta)=\zeta^2-\zeta$ has the roots $0$ and $1$. The unstable method, although of order $3$, gives errors at $t=1$ that grow explosively as the step decreases, by a factor per additional step that rises from about $4$ to about $5$; the Adams--Bashforth errors decrease by a factor close to $4$ with each halving of $h$, as befits a convergent method of order $2$.

\begin{table}
\sdfloat{%
\begin{sdtable}{colspec={rrll}, row{1}={font=\small\bfseries}}
\toprule
$h$ & $n$ & unstable, order $3$ & Adams--Bashforth \\
\midrule
$0.1$ & $10$ & $7.0$ & $1.5\times10^{-3}$ \\
$0.05$ & $20$ & $4.7\times10^{6}$ & $3.8\times10^{-4}$ \\
$0.025$ & $40$ & $2.9\times10^{19}$ & $9.5\times10^{-5}$ \\
$0.0125$ & $80$ & $1.7\times10^{46}$ & $2.4\times10^{-5}$ \\
$0.00625$ & $160$ & $8.7\times10^{100}$ & $6.0\times10^{-6}$ \\
\bottomrule
\end{sdtable}}{Absolute errors at $t=1$ for $y'=-y$, $y(0)=1$, with $n=1/h$ steps, of the method $y_{n+2}+4y_{n+1}-5y_n=h(4f_{n+1}+2f_n)$, which violates the root condition, and of the second-order Adams--Bashforth method (\cref{thm:dahlquist}).\label{tab:unstable-multistep}}
\end{table}

\subsection{Chebyshev Polynomials}

\Cref{sec:constant-coefficient-recurrences} gave three formulas for $T_n(x)$: the recurrence \eqref{eq:chebyshev-recurrence}, the closed form $\cos(n\arccos x)$ of \cref{prop:chebyshev-closed}, and the polynomial with its integer coefficients, evaluated by Horner's scheme. \Cref{tab:chebyshev-evaluation} compares them at the binary64 number nearest to $0.9$, against exact rational values, as \cref{tab:legendre-evaluation} did for the Legendre polynomials. The coefficients have the same defect as there: they alternate in sign and grow like $2^{n-1}$, the terms $c_kx^k$ reach $8\times10^{54}$ for $n=160$, and Horner's scheme has no correct digit from $n=80$ on. The recurrence, for which no solution is minimal on $(-1,1)$, had relative errors below $6\times10^{-16}$ at all five degrees. The closed form was as accurate except at $n=80$, where $T_{80}(0.9)=-0.046$ is small and the error $3.5\times10^{-14}$ appears in relative terms. At $x=0.3$ the closed form had the relative error $4.1\times10^{-13}$ for $n=160$, and at $x=0.99$ the recurrence had errors up to $7.6\times10^{-15}$ for these degrees; Volume~XI analyses such errors.

\begin{table}
\sdfloat{%
\begin{sdtable}{colspec={rlllll}, row{1}={font=\small\bfseries}}
\toprule
$n$ & $T_n(0.9)$ & largest $|c_kx^k|$ & recurrence & $\cos(n\arccos x)$ & coefficients \\
\midrule
$10$ & $-0.2007475$ & $5.95\times10^{2}$ & $5.7\times10^{-16}$ & $9.5\times10^{-16}$ & $5.7\times10^{-14}$ \\
$20$ & $-0.9194009$ & $1.50\times10^{6}$ & $2.7\times10^{-17}$ & $1.5\times10^{-16}$ & $2.5\times10^{-11}$ \\
$40$ & $0.6905961$ & $1.12\times10^{13}$ & $5.8\times10^{-16}$ & $9.0\times10^{-16}$ & $8.9\times10^{-4}$ \\
$80$ & $-0.0461542$ & $9.23\times10^{26}$ & $3.9\times10^{-17}$ & $3.5\times10^{-14}$ & $1.1\times10^{12}$ \\
$160$ & $-0.9957396$ & $8.26\times10^{54}$ & $4.1\times10^{-16}$ & $3.0\times10^{-16}$ & $3.4\times10^{38}$ \\
\bottomrule
\end{sdtable}}{The Chebyshev polynomial $T_n(x)=\sum_kc_kx^k$ at the binary64 number $x$ nearest to $0.9$: the largest term of the coefficient form, and the relative errors of the recurrence \eqref{eq:chebyshev-recurrence}, of $\cos(n\arccos x)$ with the library functions of Python, and of Horner's scheme with the integer coefficients rounded to binary64.\label{tab:chebyshev-evaluation}}
\end{table}

\subsection{The Lesson}

The experiments repeat one lesson in several forms: a recurrence that defines a sequence exactly may or may not compute it accurately, and what decides is the relation between the wanted solution and the other solutions of the recurrence. A dominant solution, such as the Fibonacci numbers, $P_n(x)$ for $x>1$ or the Chebyshev polynomials outside $[-1,1]$, is computed accurately in the forward direction. A solution that is comparable with all others, such as $P_n(x)$ and $T_n(x)$ for $|x|<1$, is computed with errors that grow at most slowly. A minimal solution, such as $J_n(x)$ for $n>x$ or the integrals $I_n$, is destroyed in the forward direction and recovered in the backward direction, from almost arbitrary starting values. The same analysis explains the parasitic solutions of multistep methods, where the wanted solution is minimal and the method must keep the others from growing. Volume~XI computes special functions with recurrences chosen and run in this way, and with the continued fractions to which they are related; Volume~VI analyses the accumulation of rounding errors in recurrences and sums.

\bigskip
A difference equation determines its solutions from initial values without any condition on the right side, and its linear theory parallels that of linear differential equations: the shift $E$ and the difference $\Delta=E-I$ take the place of $D$, falling factorials that of powers, sums that of integrals, and Newton's formula that of Taylor's. The solutions of a linear recurrence of order $k$ form a space of dimension $k$; with constant coefficients they are spanned by the sequences $n^jr^n$ for the characteristic roots $r$, together with sequences that vanish from some index on if $0$ is a root, and for real coefficients by $n^j\rho^n\cos n\theta$ and $n^j\rho^n\sin n\theta$. The Casorati determinant follows the discrete form of Abel's formula, and inputs that are polynomials times geometric sequences have particular solutions of the same kind, multiplied by $n^s$ at resonance. Generating functions of solutions are exactly the rational functions, whose poles give the growth of the solutions, and exponential generating functions turn recurrences into differential equations. A linear system $\mathbf x_{n+1}=A\mathbf x_n$ is asymptotically stable exactly when every eigenvalue has modulus less than $1$; sampling a differential system preserves this, while Euler's method does so only for small steps, and Lyapunov's equation, Duhamel's formula and the criteria on the coefficients all have their discrete versions. A three-term recurrence may have a minimal solution, which the forward recurrence cannot compute and Miller's backward recurrence can, and continued fractions, orthogonal polynomials and Jacobi matrices are tied to such recurrences. Multistep methods add parasitic solutions, and the root condition decides whether they converge. Nonlinear recurrences have fixed points whose stability is decided by $|f'(x^*)|<1$, and the logistic map passes through a cascade of period doublings into chaos.

This chapter closes Part~II and with it Volume~II. The four parts of Volumes~I and~II have laid the foundations on which the series builds: the language of logic, sets and numbers, and the combinatorics of counting; the calculus of one and several variables, with series, approximation and vector calculus; linear algebra from linear systems to the singular value decomposition, the Jordan form and functions of matrices; and differential and difference equations. Throughout, results were stated in full and the important ones proved, while some proofs were only sketched, and some facts were taken from analysis without proof: the completeness of the real numbers as an axiom, the exchange of limits under uniform convergence, Peano's existence theorem and the existence and uniqueness theorem for nonlinear systems, and the complex analysis behind radii of convergence and the Schur--Cohn test. The Lebesgue integral was only mentioned.

Volume~III supplies these foundations. Its real analysis proves the completeness of the real numbers and its consequences, develops sequences and series of functions with uniform convergence, and proves the contraction mapping theorem together with the existence and uniqueness theorem for differential equations; its complex analysis gives Cauchy's theorem, Laurent series, the residue theorem and the argument principle with Rouch\'e's theorem; and its measure theory builds the Lebesgue integral and the $L^p$ spaces. Volume~IV builds probability and statistics on that foundation, from probability spaces and limit theorems to Markov chains and random walks, of which the gambler's ruin was a first example, and to statistical inference. Volume~V turns to computing: data structures, the analysis of algorithms with the recurrences of divide and conquer and the master theorem, and the architecture of computers, including the floating-point numbers whose rounding errors appeared in every computational section of Volumes~I and~II. With these five volumes in hand, the topical volumes from Volume~VI on take up the algorithms of the stardust library, beginning with computer arithmetic.

\chapternotes

Leonardo of Pisa posed the rabbit problem that gives the Fibonacci numbers in his \emph{Liber abaci} of 1202, which Sigler translated~\cite{sigler2002}. De Moivre treated recurrent series by what we now call generating functions, and the closed form of the Fibonacci numbers, published by Binet in 1843, was known to him~\cite[Sect.~1.2.8]{knuth1997taocp1}. Boole's treatise of 1860~\cite{boole1860} presented the calculus of finite differences and difference equations as a calculus of operators, and Milne-Thomson's book~\cite{milnethomson1933} is a classical account of the subject. Poincar\'e's theorem on the ratios of solutions appeared in 1885~\cite{poincare1885}. Gautschi's survey of 1967~\cite{gautschi1967} presents the theory of minimal solutions and its use in the computation of special functions by recurrences, with Pincherle's theorem and Miller's backward recurrence. Clenshaw introduced his algorithm in 1955~\cite{clenshaw1955}, Favard proved his theorem in 1935~\cite{favard1935}, and Golub and Welsch computed Gauss rules from Jacobi matrices in 1969~\cite{golubwelsch1969}. Leslie introduced his matrix model of populations in 1945~\cite{leslie1945}. Dahlquist's equivalence theorem dates from 1956~\cite{dahlquist1956}. May's article of 1976~\cite{may1976} brought the complicated dynamics of the logistic map to wide attention; Li and Yorke's theorem appeared in 1975~\cite{li1975}, Sharkovskii's in 1964~\cite{sharkovskii1964}, and Feigenbaum's discovery of universality in 1978~\cite{feigenbaum1978}.

Elaydi~\cite{elaydi2005} and Kelley and Peterson~\cite{kelley2001} are introductions to difference equations at the level of this chapter and beyond, with the stability theory and asymptotic methods; Elaydi also treats the $z$-transform, the theorems of Poincar\'e and Perron, and the relation to continued fractions and orthogonal polynomials~\cite[Chaps.~6, 8 and~9]{elaydi2005}; Graham, Knuth and Patashnik~\cite{graham1994concrete} develop the finite calculus and generating functions for the analysis of algorithms. Gautschi's survey~\cite{gautschi1967} and Wimp's book~\cite{wimp1984} treat the computation of solutions of recurrences; Chihara~\cite{chihara1978} and Szeg\H{o}~\cite{szego1975} treat orthogonal polynomials and their recurrences. Henrici~\cite{henrici1962discrete} and Hairer, N{\o}rsett and Wanner~\cite[Chap.~III]{hairer1993} develop multistep methods and their stability, Henrici~\cite{henrici1974} and Jury~\cite{jury1964} the location of roots in the unit disc, and Strogatz~\cite[Chap.~10]{strogatz2015} gives an introduction to one-dimensional maps and chaos; Hirsch, Smale and Devaney~\cite[Chap.~15]{hirsch2013} treat discrete dynamical systems, with the logistic map and symbolic dynamics.

\chapterreferences
